\documentclass[11pt]{article}
\usepackage[T1]{fontenc}
\usepackage[utf8]{inputenc}
\usepackage{lmodern,amsmath,amssymb,amsthm,mathtools,mathrsfs}
\usepackage[margin=26mm]{geometry}
\usepackage{booktabs,longtable,array}
\usepackage[hidelinks]{hyperref}
\usepackage{microtype}
\setlength{\emergencystretch}{2em}
\raggedbottom
\numberwithin{equation}{section}
\newtheorem{theorem}{Theorem}[section]
\newtheorem{lemma}[theorem]{Lemma}
\newtheorem{proposition}[theorem]{Proposition}
\newtheorem{corollary}[theorem]{Corollary}
\theoremstyle{definition}
\newtheorem{definition}[theorem]{Definition}
\newtheorem{remark}[theorem]{Remark}
\newcommand{\C}{\mathbb C}
\newcommand{\R}{\mathbb R}
\newcommand{\Q}{\mathbb Q}
\newcommand{\Z}{\mathbb Z}
\newcommand{\T}{\mathbb T}
\newcommand{\ii}{\mathrm i}
\newcommand{\dd}{\,\mathrm d}
\newcommand{\eps}{\varepsilon}
\renewcommand{\DH}{\mathcal D_H}
\newcommand{\DL}{\mathcal D_L}
\newcommand{\XH}{\mathcal X_H}
\newcommand{\XL}{\mathcal X_L}
\newcommand{\cO}{\mathcal O}
\newcommand{\norm}[1]{\left\lVert#1\right\rVert}
\DeclareMathOperator{\Pf}{Pf}
\DeclareMathOperator{\Res}{Res}
\DeclareMathOperator{\dist}{dist}
\hypersetup{pdftitle={Natural boundaries of the full square-lattice Ising susceptibility in two complex coupling variables: Spectral domains and physical-slice restrictions},pdfauthor={Dehao Lin},pdfsubject={Natural boundaries for two independent complex Ising couplings}}
\title{Natural boundaries of the full square-lattice Ising\newline
susceptibility in two complex coupling variables\newline
\large Spectral domains and physical-slice restrictions}
\author{Dehao Lin\\School of Physics, Sun Yat-sen University\\\small ORCID: \href{https://orcid.org/0009-0001-4551-8490}{0009-0001-4551-8490}\\\small\texttt{lindh9@mail2.sysu.edu.cn}}
\date{October 2026}
\begin{document}
\maketitle
\begin{abstract}
We prove natural-boundary results for the complete connected zero-field
response of the infinite nearest-neighbor square-lattice Ising model
with two independent complex couplings. We first identify explicit
high- and low-temperature domains in \(\C^2\) as components of a
spectral complement. On each domain the physical full sum has a
single-valued holomorphic continuation, and this continuation fails to
extend through every boundary point.

Three local arguments produce the boundary obstructions. At regular
positive edges, a first singular derivative survives the entire
higher-particle tail. At negative odd-particle edges, exact pole periods
make one particle index dominate selected high Taylor coefficients of
the complete sum. On the collinear isotropic restriction, a separate
even--odd tail estimate supplies the obstruction. The regular-edge
arguments hold along arbitrary nonramified holomorphic transverse
curves and locally uniformly in complex dispersion parameters.

Exact pullbacks recover both prescribed temperature domains for every
finite positive coupling ratio, including their exceptional boundary
points. The map \(s=\sinh z\) gives the full isotropic interior and
exterior descendants with their distinct physical germs. Every result
retains the onsite term, all lattice sites and particle indices, the
plus-phase subtraction, and the relation \(\chi=\beta\mathcal X\).
The natural-boundary statements concern the physical frontiers; no
singularity assertion is made throughout the interior of the spectral set.
\end{abstract}

\section{Main results and the physical observables}
\label{main:statement}

The zero-field susceptibility sums spin correlations over the entire
square lattice. To derive a natural boundary from its form-factor
expansion, singularities of finite-particle integrals must be combined
with estimates for the remaining infinite sum
\cite[Section 3.2]{ONGP}\cite[Section 6]{BHM}.
With independent complex couplings, nonextension on physical temperature
lines also requires a separate argument.

The variables \(K_x,K_y\) are independent coupling coordinates.
Put
\[
 \begin{gathered}
 z=2K_x=x+\ii y,\qquad w=2K_y=u+\ii v,\\
 S=\sinh z,\quad T=\sinh w,\quad A=\cosh z\cosh w,\quad p=ST.
 \end{gathered}
\]
For real positive couplings, \(\mathcal X_H\) is the sum of all
two-spin correlations in the unique high-temperature Gibbs state;
\(\mathcal X_L\) is the sum of all connected two-spin correlations
in the infinite-volume plus phase. Their precise definitions and the
physical interface are in Section~\ref{ana:physical}.

Define independently of any susceptibility
\[
 \Sigma=\{(K_x,K_y):A=Sa+Tb\text{ for some }a,b\in[-1,1]\}.
\]
The two domains used in this paper are explicitly
\begin{align}
 \DH=\bigl\{&|y|<\pi/2,\ |v|<\pi/2,\nonumber\\[-2pt]
       &\cosh x\tanh|u|<\cos y,\quad
        \cosh u\tanh|x|<\cos v\bigr\},\label{main:DH}\\
 \DL=\bigl\{&x>0,\ u>0,\quad
       \cosh x\tanh u>|\cos y|,\quad
       \cosh u\tanh x>|\cos v|\bigr\}.\label{main:DL}
\end{align}
These are open sets in the original coupling space, using the displayed
invertible linear change of coordinates. The HT domain is unbounded
along decoupled directions; boundedness is asserted only for its
positive-ratio temperature slices.

\begin{theorem}[Joint physical natural boundaries]\label{main:joint}
The following assertions hold for both phases of the full bulk reduced
susceptibility.
\begin{enumerate}
\item The sets \(\DH\) and \(\DL\) are precisely the components of
\(\C^2\setminus\Sigma\) containing, respectively, the real positive
regions \(0<p<1\) and \(p>1\). They are contractible and
pseudoconvex. On them \(|p|<1\) and \(|p|>1\), respectively.
\item The physical germs extend to single-valued jointly holomorphic
functions \(\XH\in\mathcal O(\DH)\) and
\(\XL\in\mathcal O(\DL)\). Their complete odd and connected-even
form-factor sums converge normally on compact subsets, with every
fixed mixed complex derivative.
\item At every point of the respective boundary, there is no
holomorphic function on a neighborhood agreeing with that physical
continuation on an accumulating connected component of its domain
side. These are thus natural boundaries also for continuation in
unramified Riemann domains over \(\C^2\), reached from the given
physical germ.
\item The proof gives more than an ambient obstruction. At a dense
set on every regular boundary face, the complete physical sum is
nonextendable along every holomorphic transverse curve on which the
local edge parameter is nonramified. The estimates are locally
uniform for a fixed selected phase and compact admissible subfamilies
of bounded holomorphic complex-parameter families on common polydisks
in a regular-edge chart. No uniformity over the primes or through
\(\Im B=0\) is required. The degenerate equal-coupling restriction has the additional
dense local obstructions of Proposition~\ref{iso:local-obstruction}.
\end{enumerate}
Here regular faces and their edge parameters are specified by
Definition~\ref{geo:regular-definition} and
Proposition~\ref{geo:local-coordinates}. The final assertion is part
of the parent construction needed for the exact restrictions below.
\end{theorem}

For every finite \(\lambda>0\), put
\(\iota_\lambda(\zeta)=(\zeta/2,\lambda\zeta/2)\). Let
\[
 P_\lambda(a)=\cosh a\tanh(\lambda a),\qquad
 Q_\lambda(a)=\cosh(\lambda a)\tanh a,\qquad
 Y_\lambda=\frac{\pi}{2\max(1,\lambda)}.
\]
Both functions on \([0,\infty)\) have increasing inverses. Define
\begin{align*}
 H_\lambda(b)&=\min\{P_\lambda^{-1}(\cos b),
                      Q_\lambda^{-1}(\cos(\lambda b))\},
                 &&|b|<Y_\lambda,\\
 L_\lambda(b)&=\max\{P_\lambda^{-1}(|\cos b|),
                      Q_\lambda^{-1}(|\cos(\lambda b)|)\},
                 &&b\in\R,\\
 \Omega_{H,\lambda}&=\{a+\ii b:|b|<Y_\lambda,\ |a|<H_\lambda(b)\},\\
 \Omega_{L,\lambda}&=\{a+\ii b:a>L_\lambda(b)\}.
\end{align*}

\begin{theorem}[Exact recovery of every positive-ratio slice]
\label{main:slices}
For every finite positive \(\lambda\),
\[
 \iota_\lambda^{-1}(\DH)=\Omega_{H,\lambda},\qquad
 \iota_\lambda^{-1}(\DL)=\Omega_{L,\lambda}.
\]
The restrictions of the two joint physical functions are the respective
physical reduced susceptibilities, holomorphic on these exact domains,
and nonextendable through every boundary point in the componentwise
sense of Theorem~\ref{main:joint}. They retain the complete lattice
sum and every odd or nonzero even particle order. On a physical line
with fixed \(J_x>0,\ J_y=\lambda J_x\), the same conclusions hold
for the energy-field response
\[
       \chi(\zeta)=\frac{\zeta}{2J_x}\,
                         \mathcal X(\iota_\lambda(\zeta)).
\]
The proof uses the curve-level part of Theorem~\ref{main:joint}
and its analytic construction, rather than inferring a restriction
obstruction solely from ambient nonextension.
\end{theorem}

\begin{corollary}[Exact isotropic descendants]\label{main:isotropic}
Set \(\lambda=1\) and \(s=\sinh\zeta\).
The reduced HT physical germ is single valued and holomorphic on
\(|s|<1\), with value one at zero; the reduced connected-plus-phase
LT physical germ is single valued and holomorphic on \(|s|>1\).
Every point of \(|s|=1\) is a natural-boundary point for each germ.
Their complete normalized representations are
\[
 \XH(s)=\frac{(1-s^4)^{1/4}}{s}
                      \sum_{N\ge1,\ N\ {\rm odd}}C_N(1/s),\qquad
 \XL(s)=(1-s^{-4})^{1/4}
                      \sum_{N\ge2,\ N\ {\rm even}}C_N(s),
\]
where the coefficients are the full-site integrals in
Section~\ref{iso:section}. The LT onsite term is
\(1-(1-s^{-4})^{1/4}\).
The map \(\sinh\) is a biholomorphism from \(\Omega_{H,1}\)
to the disk and is the universal covering from \(\Omega_{L,1}\)
to the exterior, with deck translations \(\zeta\mapsto\zeta+2\pi\ii k\).
The energy-field results follow with \(\beta=\zeta/(2J)\) on that
physical continuation. No single-valued inverse temperature on the
whole exterior is required.
\end{corollary}

\subsection{Related work and proof outline}

The starting geometry is related to the unsquared anisotropic Nickel
equation studied by Boukraa, Hassani and Maillard
\cite[Section 6]{BHM}. That paper distinguishes multi-variable
finite-integral singularities from the full-susceptibility question.
The passage from these finite integrals to the susceptibility requires
an obstruction for the complete particle sum. Sections~\ref{ana:edges}
and \ref{neg:section} establish that obstruction on regular faces. In particular the set \(\Sigma\) has
interior; neither its interior nor every finite-integral candidate is
declared a singularity of the physical sum.

Sections~\ref{geo:section}--\ref{main:closure} establish the joint
continuation and infinite-tail estimates required for
Theorems~\ref{main:joint}--\ref{main:slices} and
Corollary~\ref{main:isotropic}. Section~\ref{ana:interface} specifies
the real-temperature exact-solution input.

The joint domains contain points outside all positive-ratio
temperature lines. For example
\(z=\operatorname{arsinh}(1/2),\ w=\ii\pi/6\) is an interior point
of \(\DH\) and lies on no positive-real-ratio temperature line.
The independent complex coordinates \((C,B)=(A/S,T/S)\) have
nonzero Jacobian at every regular edge. Their common polydisks,
constant-\(B\) transversals, and variable-\(B\) curve estimates
are used explicitly. Conversely the exact identity
\(\XH(K_x,0)=e^{2K_x}\) shows why restriction nonextension needs
its own proof.

The proof has three stages. Section~\ref{geo:section} determines the
domains and the geometry of their frontiers.
Section~\ref{ana:physical} constructs the holomorphic functions from
the complete physical sums. Sections~\ref{ana:edges}--\ref{iso:section}
then produce local obstructions that survive those sums.
Section~\ref{main:closure} combines them to prove
Theorems~\ref{main:joint}--\ref{main:slices} and
Corollary~\ref{main:isotropic}.



\par\medskip\noindent\textbf{AI Declaration.} ChatGPT was used extensively in the mathematical research,
proof development, and preparation of this manuscript.\par

\clearpage
\tableofcontents
\clearpage

\section{Independent complex spectral geometry}
\label{geo:section}

We determine the spectral complement before constructing the
susceptibility. The calculation supplies the physical components, their
local edge coordinates and their exact temperature slices. Its input
is the spectral equation alone. Write
\[
 \begin{gathered}
 z=2K_x=x+\ii y,\qquad w=2K_y=u+\ii v,\\
 S=\sinh z,\quad T=\sinh w,\quad A=\cosh z\cosh w,\quad p=ST,
 \end{gathered}
\]
and define the closed set
\begin{equation}\label{geo:sigma}
 \Sigma=\{(K_x,K_y):A=S a+T b
                    \text{ for some }(a,b)\in[-1,1]^2\}.
\end{equation}
The two real functions that describe this set are
\begin{equation}\label{geo:F}
 F_1=\cosh^2u\cos^2y-\cosh^2x\sinh^2u,\qquad
 F_2=\cosh^2x\cos^2v-\cosh^2u\sinh^2x.
\end{equation}
The use of \(z,w\) instead of \(K_x,K_y\) is just an invertible linear
coordinate change; all domains below are domains in the original
\(\C^2\) coupling space.

\begin{theorem}[Exact spectral set]\label{geo:exact-sigma}
For two independent complex couplings, including zero, opposite-sign,
and collinear real-part configurations,
\begin{equation}\label{geo:sigma-F}
                         \Sigma=\{F_1F_2\leq0\}.
\end{equation}
\end{theorem}

\begin{proof}
\emph{Reduction and the noncollinear case.}
The independent reflections \(z\mapsto-z\) and \(w\mapsto-w\)
preserve \(\Sigma\), by changing the sign of the corresponding
representing coefficient.  They also preserve \(F_1,F_2\).
We may therefore first suppose \(x,u>0\).  Set
\[
 \begin{gathered}
 h=\sinh x,\quad c=\cosh x,\quad H=\sinh u,\quad C_0=\cosh u,\\
 t=\cos y,\quad q=\sin y,\quad V=\cos v,\quad Q=\sin v.
 \end{gathered}
\]
For
\[
 D=\Im(\overline S T),\qquad
 N_a=\Im(\overline A T),\qquad N_b=\Im(\overline S A),
\]
direct multiplication gives
\begin{align}
 D&=hC_0tQ-cHqV,\nonumber\\
 N_a&=ctVQ-hHC_0q,\qquad
 N_b=hcHQ-C_0tqV. \label{geo:determinants}
\end{align}
When \(D\ne0\), the unique real solution of \(A=Sa+Tb\) is
\(a=N_a/D,\ b=N_b/D\).

For the sign analysis put
\[
 \begin{gathered}
 a_0=h^2,\quad b_0=H^2,\quad A_0=1+a_0,\quad B_0=1+b_0,\\
 X=t^2,\quad Y=V^2,\quad
 X_0=\frac{A_0b_0}{B_0},\quad Y_0=\frac{B_0a_0}{A_0}.
 \end{gathered}
\]
Expanding \eqref{geo:determinants} and cancelling the common cross
terms gives the identities
\begin{align}
 F_1&=B_0(X-X_0),& F_2&=A_0(Y-Y_0),\nonumber\\
 N_a^2-D^2&=F_2K,&K&=(B_0-Y)X-b_0,\label{geo:factorization}\\
 N_b^2-D^2&=-F_1L,&L&=a_0-(A_0-X)Y.\nonumber
\end{align}
Both \(K\) and \(L\) vanish at \((X_0,Y_0)\), increase with
\(X\), and decrease with \(Y\).  If \(F_1,F_2\) have strictly
opposite signs, both differences of squares in
\eqref{geo:factorization} are strictly negative.  Consequently
\(D\ne0\) and \(|a|,|b|<1\); the point belongs to \(\Sigma\).

Suppose instead that \(F_1,F_2\) have the same strict sign and
\(D\ne0\).  For \(0<Y<1\), membership would impose
\[
 X\leq\frac{b_0}{B_0-Y},\quad
 X\geq A_0-\frac{a_0}{Y}
 \quad\text{if }F_1,F_2>0,
\]
and the reverse two inequalities if \(F_1,F_2<0\).  But
\begin{equation}\label{geo:threshold-difference}
 \frac{b_0}{B_0-Y}-\left(A_0-\frac{a_0}{Y}\right)
 =\frac{A_0(Y-1)(Y-Y_0)}{Y(B_0-Y)}.
\end{equation}
For the positive-sign case \(Y>Y_0\), this expression is negative;
for the negative-sign case \(Y<Y_0\), it is positive.  Each case
contradicts the required inequalities.  At \(Y=0\), \(F_2<0\)
and \(L=a_0>0\), so the negative-sign case is impossible directly
from \eqref{geo:factorization}.  At \(Y=1\),
\(K=b_0(X-1)\) and \(L=X-1\); either same-sign case requires
\(X=1\), which gives \(D=0\).

It remains to exclude a same-strict-sign point with \(D=0\)
from \(\Sigma\).  Membership would give \(N_a=N_b=0\).
Since \(F_1F_2\ne0\), \eqref{geo:factorization} implies
\(K=L=0\).  Here \(Y>0\), and \eqref{geo:threshold-difference}
then implies \(Y=Y_0\) or \(Y=1\).  The first contradicts
\(F_2\ne0\), whereas the second gives \(X=1\).  At \(X=Y=1\)
the segment condition is
\[
                         cC_0\leq h+H.
\]
As
\[
                 (cC_0)^2-(h+H)^2=(hH-1)^2,
\]
this requires \(hH=1\), and hence \(F_1=F_2=0\), again a
contradiction.

\emph{Equality faces.} Endpoint roots give a direct description of
the equality cases and will also provide the later edge coordinates.
At a point with \(F_1=0\), \(S\ne0\), and
\(\tanh(z/2)\) is finite and nonzero.  The two endpoint roots
\begin{equation}\label{geo:endpoint-roots}
 r_+=e^w\tanh(z/2),\qquad r_-=e^{-w}\tanh(z/2)
\end{equation}
satisfy
\begin{equation}\label{geo:endpoint-identities}
 \frac{r_++r_+^{-1}}2=\frac{A-T}{S},\qquad
 \frac{r_-+r_-^{-1}}2=\frac{A+T}{S}.
\end{equation}
Moreover,
\begin{equation}\label{geo:tanh-modulus}
 |\tanh(z/2)|^2=\frac{\cosh x-\cos y}{\cosh x+\cos y}.
\end{equation}
Thus \(F_1=0\) means that at least one of \(r_+,r_-\) has
modulus one: the respective equations are
\(\cos y=\cosh x\tanh u\) and
\(\cos y=-\cosh x\tanh u\).  The corresponding right side
of \eqref{geo:endpoint-identities} belongs to \([-1,1]\).
This supplies a representing pair with \(b=1\) or \(b=-1\).
Interchanging \(z,w\) proves membership when \(F_2=0\).

\emph{Vanishing real parts.} The remaining cases follow directly from the real and imaginary
parts of the spectral equation.  Let \(x=0,u>0\).  If \(\cos v\ne0\),
the real part of the equation \(A=Sa+Tb\) forces
\[
 b=\frac{\cosh u}{\sinh u}\cos y.
\]
Its necessary condition is \(|\cos y|\leq\tanh u\).
Under that condition \(\sin y\ne0\), and
\[
 a=-\frac{\cos y\sin v}{\sinh u\sin y}
\]
has modulus at most one, so the condition is also sufficient.
Indeed
\(\cos^2y\leq\sinh^2u/\cosh^2u
\leq\sinh^2u/(\sinh^2u+\sin^2v)\).
Here \(F_2=\cos^2v>0\), and the condition is precisely \(F_1\leq0\).
If \(\cos v=0\), all terms are imaginary, and the choice
\[
                 a=0,\qquad b=\tanh u\cos y
\]
always represents \(A\); here \(F_2=0\).
The case \(u=0,x>0\) follows by exchange.  Finally, if \(x=u=0\),
the coefficients \(S,T\) are imaginary and \(A=\cos y\cos v\)
is real.  Membership is equivalent to \(A=0\), represented by
\(a=b=0\), and is exactly \(F_1F_2=0\).
The initial reflections now cover all real-part signs.
\end{proof}

\subsection{The physical components and their branches}
\label{geo:components-section}

The sign calculation now becomes a description of connected domains.
The contractions used in the proof also fix the branches of the
physical fourth-root factors.

\begin{theorem}[All components; explicit physical domains]
\label{geo:components}
The components of \(\C^2\setminus\Sigma\) with \(F_1,F_2>0\)
are indexed by \(m,n\in\Z\) and are
\[
 F_1>0,\quad F_2>0,\quad
 |y-m\pi|<\pi/2,\quad |v-n\pi|<\pi/2.
\]
Each is star-shaped about the point
\((z,w)=(\ii m\pi,\ii n\pi)\).
The components with \(F_1,F_2<0\) are the four sets specified
by the independent signs of \(x,u\); each is contractible.

The component containing the origin is exactly
\begin{equation}\label{geo:DH}
 \DH=\left\{
 \begin{array}{l}
 |y|<\pi/2,\quad |v|<\pi/2,\\
 \cosh x\tanh|u|<\cos y,\quad
 \cosh u\tanh|x|<\cos v
 \end{array}\right\}.
\end{equation}
The component containing the positive real low-temperature region is
\begin{equation}\label{geo:DL}
 \DL=\left\{
 \begin{array}{l}
 x>0,\quad u>0,\\
 \cosh x\tanh u>|\cos y|,\quad
 \cosh u\tanh x>|\cos v|
 \end{array}\right\}.
\end{equation}
They satisfy the strict inequalities
\begin{equation}\label{geo:p-domains}
                         |p|<1\text{ on }\DH,\qquad
                         |p|>1\text{ on }\DL.
\end{equation}
\end{theorem}

\begin{proof}
The sign description follows from Theorem~\ref{geo:exact-sigma}.
For \(F_1,F_2>0\), both cosines are nonzero, so a connected
component remains in a single indicated cosine strip in each
variable.  Within a fixed pair of strips, set
\(\eta=y-m\pi,\ \theta=v-n\pi\).  The inequalities are
\[
 \cosh x\tanh|u|<\cos\eta,\qquad
 \cosh u\tanh|x|<\cos\theta.
\]
Replacing \((x,u,\eta,\theta)\) by
\(t(x,u,\eta,\theta)\), \(0\leq t\leq1\), decreases the
left sides and increases the right sides.  This proves
star-shapedness and excludes further component subdivision.

For \(F_1,F_2<0\), neither \(x\) nor \(u\) can vanish.
Writing \(a=|x|,\ b=|u|\), the conditions are
\[
 \cosh a\tanh b>|\cos y|,\qquad
 \cosh b\tanh a>|\cos v|.
\]
Each left side is nondecreasing in both \(a,b\), and strictly
increasing in the variable inside its hyperbolic tangent.
Choose once and for all \(R>0\) with \(\sinh R>1\).
First increase \((a,b)\) continuously to
\((\max(a,R),\max(b,R))\).  Then move these two coordinates
linearly to \((R,R)\), throughout staying in \([R,\infty)^2\).
During the second move both left sides are at least
\(\cosh R\tanh R=\sinh R>1\).  Finally move \(y,v\) linearly
to zero.  These three moves give a continuous contraction
within the fixed sign component.  The signs cannot change
without leaving the set.  This proves all component claims
and the physical identifications \eqref{geo:DH}--\eqref{geo:DL}.

The modulus identity, with the notation in the proof of
Theorem~\ref{geo:exact-sigma}, is
\begin{equation}\label{geo:p-modulus}
                  |p|^2=(A_0-X)(B_0-Y).
\end{equation}
For \(F_1,F_2>0\),
\[
 A_0-X<\frac{A_0}{B_0},\qquad
 B_0-Y<\frac{B_0}{A_0},
\]
so \(|p|<1\).  For \(F_1,F_2<0\), both inequalities reverse,
so \(|p|>1\).  These arguments allow zero real parts as well.
On the positive real quadrant the common sign of \(F_1,F_2\)
is the sign of \(1-(\sinh x\sinh u)^2\), which proves the
claimed physical identification.
\end{proof}

\begin{remark}\label{geo:algebraic-branches}
In particular \(\{|p|=1\}\subset\Sigma\).  The more specific
corner identity is
\[
 A^2-(\epsilon S+\delta T)^2=(p-\epsilon\delta)^2,
 \qquad \epsilon,\delta\in\{1,-1\},
\]
and therefore
\[
 \bigcup_{\epsilon,\delta\in\{1,-1\}}
 \{A=\epsilon S+\delta T\}=\{p^2=1\}\subset\Sigma.
\]
The branches \((1-p^2)^{1/4}\) on \(\DH\) and
\((1-p^{-2})^{1/4}\) on \(\DL\) are defined directly by the
binomial series and their positive real physical values.
No additional covering is needed for these factors.
On \(\DL\), \(S,T\) are nonzero.  On \(\DH\) the zeros of
\(S\) and \(T\) are respectively \(z=0\) and \(w=0\);
these are genuine complex subvarieties of the joint domain.
Any particle representation used there must account for
its removable normalization at these axes.
\end{remark}

\begin{lemma}[A holomorphic modulus description]
\label{geo:holomorphic-inequalities}
Within the two central strips in \eqref{geo:DH}, the four
inequalities defining \(\DH\) can equivalently be written
\[
 |e^w\tanh(z/2)|<1,\quad |e^{-w}\tanh(z/2)|<1,\quad
 |e^z\tanh(w/2)|<1,\quad |e^{-z}\tanh(w/2)|<1.
\]
The domain \(\DL\) is equivalently described by
\[
 \begin{split}
 |e^{-w}\tanh(z/2)|<1,\quad &|e^{-w}\coth(z/2)|<1,\\
 |e^{-z}\tanh(w/2)|<1,\quad &|e^{-z}\coth(w/2)|<1 .
 \end{split}
\]
The functions in the second display are evaluated where they
are finite; its inequalities themselves force \(x,u>0\).
\end{lemma}

\begin{proof}
For the first pair in the HT assertion,
\eqref{geo:tanh-modulus} gives
\[
 e^{2|u|}\frac{\cosh x-\cos y}{\cosh x+\cos y}<1,
\]
equivalent to \(\cos y>\cosh x\tanh|u|\).  Exchange variables
for the other pair.  In the first LT pair, multiplication
forces \(e^{-2u}<1\), hence \(u>0\); the second pair similarly
forces \(x>0\).  If
\(r=|\tanh(z/2)|^2\), the first pair is
\(e^{-2u}<r<e^{2u}\), which is equivalent to
\(|\cos y|<\cosh x\tanh u\).  Exchange variables again.
\end{proof}

\begin{lemma}[Canonical decaying root]\label{geo:canonical-root}
For \(K=(K_x,K_y)\notin\Sigma\) and real \(q\), the equation
\begin{equation}\label{geo:root-equation}
               S\rho^2-2(A-T\cos q)\rho+S=0
\end{equation}
has exactly one root of modulus less than one, interpreting
that root as zero when \(S=0\).  This root is holomorphic in
the independent coupling variables and locally jointly
holomorphic with the complex angle \(q\).  On a compact
subset of \(\C^2\setminus\Sigma\), its modulus has a uniform
gap below one for all real \(q\).  On both positive real
physical components it is the positive decaying root.
\end{lemma}

\begin{proof}
When \(S\ne0\), put \(W=(A-T\cos q)/S\).
By the definition of \(\Sigma\), \(W\notin[-1,1]\).
The two roots are reciprocal, neither has modulus one,
and exactly one lies in the unit disk.  Equivalently, it is
\(\rho=W-\sqrt{W^2-1}\), with the square-root branch
asymptotic to \(W\) at infinity.  Its uniqueness makes these
local holomorphic roots agree on overlaps.  At \(S=0\),
\(A-T\cos q\ne0\); the implicit function theorem at
\(\rho=0\) supplies the same holomorphic root.
Compactness gives both the modulus gap and a common thin
complex angular neighborhood on any compact parameter set.
For positive real \(S,T\),
\[
              (A-S-T)(A+S+T)=(p-1)^2>0
\]
on either physical component, and hence \(W>1\).
The interior root is then real and belongs to \((0,1)\).
\end{proof}

\subsection{Exact fixed-ratio pullbacks}
\label{geo:slices-section}

A temperature line is obtained by substituting the fixed ratio into
the joint inequalities. We perform this substitution exactly so that
the later nonextension theorem has the correct one-variable domain.

\begin{theorem}[Exact slice domains]\label{geo:exact-pullbacks}
For every finite \(\lambda>0\), let
\[
 \iota_\lambda(\zeta)=(\zeta/2,\lambda\zeta/2),\qquad
 Y_\lambda=\frac{\pi}{2\max(1,\lambda)},
\]
and, for \(a\geq0\), let
\[
 P_\lambda(a)=\cosh a\,\tanh(\lambda a),\qquad
 Q_\lambda(a)=\cosh(\lambda a)\,\tanh a.
\]
Both functions are strictly increasing bijections of
\([0,\infty)\) onto itself.  Define
\begin{align*}
 H_\lambda(b)&=\min\{
  P_\lambda^{-1}(\cos b),\
  Q_\lambda^{-1}(\cos(\lambda b))\},
       &&|b|<Y_\lambda,\\
 L_\lambda(b)&=\max\{
  P_\lambda^{-1}(|\cos b|),\
  Q_\lambda^{-1}(|\cos(\lambda b)|)\},
       &&b\in\R.
\end{align*}
Then the exact equalities of open sets are
\begin{align}
 \iota_\lambda^{-1}(\DH)
 &=\Omega_{H,\lambda}
   =\{a+\ii b:|b|<Y_\lambda,\ |a|<H_\lambda(b)\},
                         \label{geo:HT-pullback}\\
 \iota_\lambda^{-1}(\DL)
 &=\Omega_{L,\lambda}
   =\{a+\ii b:a>L_\lambda(b)\}.
                         \label{geo:LT-pullback}
\end{align}
Both domains are connected and simply connected.
There is a unique \(z_c>0\) with
\[
 \sinh z_c\sinh(\lambda z_c)=1,\qquad
 P_\lambda(z_c)=Q_\lambda(z_c)=1.
\]
The function \(H_\lambda\) is positive inside its interval
and extends continuously by zero at both endpoints,
and \(0\leq L_\lambda\leq z_c\).
\end{theorem}

\begin{proof}
The derivatives
\[
 P_\lambda'(a)=\sinh a\tanh(\lambda a)
              +\lambda\cosh a\,\operatorname{sech}^2(\lambda a),
\quad
 Q_\lambda'(a)=\lambda\sinh(\lambda a)\tanh a
              +\cosh(\lambda a)\operatorname{sech}^2a
\]
are positive on \([0,\infty)\), and both functions tend to
infinity.  On the slice,
\((x,y,u,v)=(a,b,\lambda a,\lambda b)\).
The two strip conditions in \eqref{geo:DH} are exactly
\(|b|<Y_\lambda\); its remaining two inequalities are
\(P_\lambda(|a|)<\cos b\) and
\(Q_\lambda(|a|)<\cos(\lambda b)\).
This proves \eqref{geo:HT-pullback}.
In \eqref{geo:DL} the positive-real-part conditions are
exactly \(a>0\), and the other inequalities invert to
\eqref{geo:LT-pullback}.

The positive function \(\sinh a\sinh(\lambda a)\) increases
strictly from zero to infinity.  The identities
\[
 P_\lambda(a)^2-1
 =\frac{\sinh^2a\sinh^2(\lambda a)-1}{\cosh^2(\lambda a)},
\quad
 Q_\lambda(a)^2-1
 =\frac{\sinh^2a\sinh^2(\lambda a)-1}{\cosh^2a}
\]
prove the assertions about \(z_c\) and \(L_\lambda\).
At either endpoint of the HT interval at least one of its
two cosine arguments is zero, which proves the endpoint
limit for \(H_\lambda\).
The HT domain is star-shaped by Theorem~\ref{geo:components},
or directly by simultaneous scaling of \(a,b\).
The LT domain is homeomorphic to
\((0,\infty)\times\R\) via
\((a,b)\mapsto(a-L_\lambda(b),b)\).
\end{proof}

\begin{lemma}[Isotropic covering geometry]\label{geo:isotropic-cover}
At \(\lambda=1\),
\[
 \iota_1^{-1}(\Sigma)=\{\zeta:|\sinh\zeta|=1\},
\]
\[
 \Omega_{H,1}
 =\{|b|<\pi/2,\ |a|<\operatorname{arsinh}(\cos b)\},
\qquad
 \Omega_{L,1}
 =\{a>\operatorname{arsinh}(|\cos b|)\}.
\]
The map \(s=\sinh\zeta\) maps \(\Omega_{H,1}\)
biholomorphically onto \(\{|s|<1\}\), and maps
\(\Omega_{L,1}\) as the universal covering of
\(\{|s|>1\}\), with deck translations
\(\zeta\mapsto\zeta+2\pi\ii k\), \(k\in\Z\).
\end{lemma}

\begin{proof}
Here \(S=T=s\), \(A=1+s^2\).  For \(s\ne0\), the spectral
condition is \(s+s^{-1}\in[-2,2]\), whose preimage is
exactly the unit circle; \(s=0\) is not spectral.
Also \(P_1=Q_1=\sinh\), giving the displayed domains.
On the central strip \(|\Im\zeta|<\pi/2\), \(\sinh\)
is injective: the alternatives for equal hyperbolic sines
are \(\zeta_1=\zeta_2+2\pi\ii k\) or
\(\zeta_1=-\zeta_2+(2k+1)\pi\ii\), and only the first
with \(k=0\) is possible there.  The domain
\(\Omega_{H,1}\) is bounded, its boundary maps into
\(|s|=1\), and its derivative has no zero inside.
Its image is consequently both open and relatively closed
in the disk, and contains zero; it is the whole disk.

For the exterior, put \(\omega=e^\zeta\).
The equation \(\omega-\omega^{-1}=2s\) has, for \(|s|>1\),
exactly one root with \(|\omega|>1\).
A root of modulus one would give \(s\in\ii[-1,1]\).
The outside root varies holomorphically, for example by
using \(\sqrt{s^2+1}=s\sqrt{1+s^{-2}}\).
Its domain is the continuous radial epigraph
\[
 \{\,r e^{\ii b}:r>\exp(\operatorname{arsinh}|\cos b|)\,\}.
 \]
Exponentiation from the simply connected epigraph
\(\Omega_{L,1}\) is its universal covering, and the
outside-root map is a biholomorphism to \(|s|>1\).
The stated deck group follows.
\end{proof}

\subsection{Regular boundary faces in two complex variables}
\label{geo:regular-section}

The local picture is a parallelogram in the dispersion coordinate:
one representing coefficient reaches an endpoint while the other
remains free. The next two propositions first show that such faces
are dense and then give coordinates in which one can cross them.

\begin{definition}\label{geo:regular-definition}
A point of \(\partial\DH\) or \(\partial\DL\) will be called
a regular joint edge point if \(S,T\ne0\),
\(D=\Im(\overline S T)\ne0\), exactly one of its unique
representing coefficients \(a,b\) is \(1\) or \(-1\),
the other lies in \((-1,1)\), and \(p^2\ne1\).
\end{definition}

\begin{proposition}[Density including exceptional strata]
\label{geo:joint-density}
Regular joint edge points are dense in each of
\(\partial\DH\) and \(\partial\DL\).
More precisely, the points on either boundary for which
\(x,u\ne0\) and exactly one of \(F_1,F_2\) vanishes
are regular and are dense.  Thus the excluded zero-real-part,
collinear, and intersection strata contain no relatively
open boundary patch.
\end{proposition}

\begin{proof}
First suppose \(x,u\ne0\), \(F_1=0\), \(F_2\ne0\).
The algebra \eqref{geo:factorization}, valid also for
negative real parts because it only uses squares, reduces to
\begin{equation}\label{geo:strict-face}
 N_a^2-D^2=-\frac{\sinh^2u}{\cosh^2u}F_2^2<0,
 \qquad N_b^2-D^2=0.
\end{equation}
Thus \(D\ne0\), \(|a|<1\), and \(|b|=1\).
Furthermore \eqref{geo:p-modulus} gives
\begin{equation}\label{geo:p-edge}
 |p|^2=1-\frac{F_2}{\cosh^2u}\quad(F_1=0),\qquad
 |p|^2=1-\frac{F_1}{\cosh^2x}\quad(F_2=0).
\end{equation}
At an HT face the other \(F_i\) is positive, and at an LT
face it is negative.  Consequently \(p^2\ne1\).
The case \(F_2=0,F_1\ne0\) follows by exchange.

Here is an explicit approximation argument covering all
remaining boundary configurations.  For \(a,b\geq0\), set
\[
 r_1(a,b)=\cosh a\tanh b,\qquad
 r_2(a,b)=\cosh b\tanh a.
\]
Both \(r_i<1\), both equal one, or both exceed one
according as \(\sinh a\sinh b<1\), \(=1\), or \(>1\).
The closure of \(\DH\) is
\begin{equation}\label{geo:HT-closure}
 |y|,|v|\leq\pi/2,\qquad
 r_1(|x|,|u|)\leq\cos y,\qquad
 r_2(|x|,|u|)\leq\cos v.
\end{equation}
Indeed, contracting all four coordinates toward zero
strictly opens every active inequality, proving the
reverse inclusion in this closure description.
If \(\sinh|x|\sinh|u|<1\), the angular fiber of
\eqref{geo:HT-closure} is the closed rectangle
\[
 |y|\leq\arccos r_1,\qquad |v|\leq\arccos r_2.
\]
If the product is one, both intervals collapse to zero.
Every boundary point has at least one active rectangular
face.  Approximate its real-part magnitudes by
\(a_j,b_j>0\) with \(\sinh a_j\sinh b_j<1\).
For an active first face choose
\(|y_j|=\arccos r_1(a_j,b_j)\), with the sign giving
\(y_j\to y\), and choose \(v_j\) strictly inside the second
interval with \(v_j\to v\).  If the first interval has
collapsed at the limit, either sign is available.
Preserve the original signs of nonzero real parts; choose
either sign at a zero real part.  This gives \(x_j,u_j\ne0\),
\(F_1=0,F_2>0\), tending to the given point.
Use the other face if it is the active one.
This also covers corners, strip endpoints, and the collapsed
fiber at real product one.

For \(\DL\), the closure is exactly
\begin{equation}\label{geo:LT-closure}
 x,u\geq0,\qquad r_1(x,u)\geq|\cos y|,\qquad
 r_2(x,u)\geq|\cos v|.
\end{equation}
Increasing both real parts by any positive amount opens
all its inequalities and proves the closure assertion.
A boundary point must have
\(\sinh x\sinh u\leq1\): if the product exceeds one,
both \(r_i>1\) and the point is interior.
It must also have at least one active cosine inequality.
For example, \(x=0\) forces the second inequality to be
active, while \(u=0\) forces the first.
Again choose \(x_j,u_j>0\) with product less than one,
tending to \(x,u\).  At an active first face choose
\[
 |\cos y_j|=r_1(x_j,u_j),\quad y_j\to y,
 \qquad
 |\cos v_j|<r_2(x_j,u_j),\quad v_j\to v .
 \]
These choices are always possible in the appropriate
local cosine intervals.  When \(r_i\to0\), these intervals
shrink to a cosine zero; when \(r_i\to1\), their endpoints
approach a multiple of \(\pi\).  The strict choice in the
second interval remains possible since \(r_2(x_j,u_j)>0\).
This gives \(F_1=0,F_2<0\), tending to the prescribed point.
The other active face is treated by exchange.
\end{proof}

\begin{proposition}[Holomorphic coordinates on a regular face]
\label{geo:local-coordinates}
Where \(S\ne0\), set
\[
                         C=A/S,\qquad B=T/S.
\]
Then
\begin{equation}\label{geo:Jacobian}
             \det\frac{\partial(C,B)}{\partial(z,w)}
                    =\frac{p^2-1}{S^3}.
\end{equation}
At a regular joint edge, \((C,B)\) are therefore local
holomorphic coordinates and \(\Im B\ne0\).
If the endpoint is \(b=\epsilon\), use
\(h=C-\epsilon B\): locally the boundary is
\(\Im h=0\), with \(\Re h\in(-1,1)\), and the physical
domain occupies exactly one of its two sides.
If the endpoint is \(a=\epsilon\), exchange the transfer
axes for the same conclusion.

On each regular boundary patch, points at which its free
coefficient is
\[
       \cos(2\pi k/q),\qquad q\text{ an odd prime},\quad
                         1\leq k\leq(q-1)/2,
\]
are dense.  This assertion specifies a dense geometric
family only; it does not assert singularity of a sum there.
\end{proposition}

\begin{proof}
The four derivatives are
\[
 C_z=-\frac{\cosh w}{S^2},\quad
 C_w=\frac{\cosh z\,T}{S},\quad
 B_z=-\frac{T\cosh z}{S^2},\quad
 B_w=\frac{\cosh w}{S},
\]
which give \eqref{geo:Jacobian}.
Also \(D=|S|^2\Im B\).  In these coordinates the spectral
condition is simply
\[
                         C\in[-1,1]+B[-1,1].
\]
For nonreal \(B\), this is a closed parallelogram in the
\(C\)-plane.  Near an interior point of its edge with
\(b=\epsilon\), the condition for the corresponding
representing coefficient is
\[
          b-\epsilon=\frac{\Im(C-\epsilon B)}{\Im B};
\]
the other coefficient remains strictly between its
endpoints.  Hence one side is spectral and the other
is the unique local complementary side.  Its intersection
with the physical component is that whole side after
shrinking the neighborhood.
Holding \(B\) fixed permits \(h\) to range over an open
real interval on the edge.  The stated prime cosine
values are dense in \((-1,1)\), since for arbitrarily
large odd primes the angles \(2\pi k/q\) approximate
every angle in \((0,\pi)\).  The coordinate description
then proves the density assertion.
\end{proof}

\begin{remark}\label{geo:endpoint-signs}
On an HT regular face \(F_1=0\), the endpoint is
\(b=\operatorname{sign}u\); on an HT regular face \(F_2=0\),
it is \(a=\operatorname{sign}x\).
This follows from \eqref{geo:endpoint-identities}, since
both cosines are positive in the physical strips.
On an LT regular face \(F_1=0\), the endpoint is
\(b=\operatorname{sign}(\cos y)\); on \(F_2=0\), it is
\(a=\operatorname{sign}(\cos v)\).
In particular both signs of endpoint must be treated in
any proof for the full physical boundaries.
\end{remark}

\subsection{Regular arcs on each fixed-ratio slice}
\label{geo:slice-regularity-section}

To prove the slice theorem we need dense regular points on each
specified temperature line. The following argument establishes that
one-variable density for every \(\lambda\ne1\), including the
local finiteness of its exceptional points.

\begin{proposition}[Every anisotropic slice has dense regular arcs]
\label{geo:slice-density}
Fix \(\lambda>0,\lambda\ne1\), and restrict \(S,T,A,p,D\)
to \((z,w)=(\zeta,\lambda\zeta)\).
On each boundary \(\partial\Omega_{H,\lambda}\) and
\(\partial\Omega_{L,\lambda}\), all but a locally finite
set of points have \(D\ne0\), exactly one endpoint
coefficient, and a nonzero derivative of the corresponding
meromorphic edge map
\begin{equation}\label{geo:slice-edge-maps}
 h_\epsilon(\zeta)=\frac{A-\epsilon S}{T}
       \quad(a=\epsilon),\qquad
 g_\epsilon(\zeta)=\frac{A-\epsilon T}{S}
       \quad(b=\epsilon).
\end{equation}
Such points lie on real analytic boundary arcs, and are
dense on both full boundaries.  On each such arc, the
prime cosine values in Proposition~\ref{geo:local-coordinates}
are dense in its free coefficient.
\end{proposition}

\begin{proof}
We first prove the required local finiteness of the
collinear candidates.  Off the imaginary axis, both
real parts \(x,\lambda x\) are nonzero.  The strict
opposite-sign case in \eqref{geo:factorization} gives
\(D\ne0\); the strict face identity
\eqref{geo:strict-face} excludes \(D=0\) on a single
zero face as well.  Consequently a collinear candidate
off the imaginary axis has \(F_1=F_2=0\), hence
\(|p|=1\).  Conversely \(|p|=1\) implies membership
in \(\Sigma\), by \eqref{geo:p-domains} applied to all
complementary components.  Thus
\begin{equation}\label{geo:collinear-equality}
 \{\,\zeta\in\iota_\lambda^{-1}(\Sigma):
                       \Re\zeta\ne0,\ D=0\,\}
 =
 \{\,\zeta:\Re\zeta\ne0,\ D=0,\ |p|=1\,\}.
\end{equation}

For \(x>0\), the modulus equation is
\begin{equation}\label{geo:modulus-graph}
 (\sinh^2x+\sin^2y)
       (\sinh^2(\lambda x)+\sin^2(\lambda y))=1.
\end{equation}
If \(|\sin y\sin(\lambda y)|<1\), its left side starts
below one and strictly increases to infinity with \(x\).
Its positive \(x\) derivative at the root gives a unique
positive real analytic solution \(x=X(y)\).
At a simultaneous zero of \(\cos y,\cos(\lambda y)\),
the root tends continuously to zero.

If there are no simultaneous cosine zeros, \(X\) is
analytic on all of \(\R\).  Since \(X\) is even near
zero, \(X'(0)=0\), and \(X(0)=z_c\).
Direct differentiation therefore gives
\begin{equation}\label{geo:collinear-derivative}
 \left.\frac{\dd}{\dd y}D(X(y)+\ii y)\right|_{y=0}
       =\lambda\coth(\lambda z_c)-\coth z_c\ne0.
\end{equation}
For the last inequality, \(t\coth t\) is strictly
increasing on \(t>0\), because its derivative has the
positive numerator \(\sinh t\cosh t-t\).
The analytic function \(D(X(y)+\ii y)\) is therefore
nonzero identically and has locally isolated zeros.

Simultaneous cosine zeros occur exactly when
\(\lambda=r/s\) in lowest positive terms with \(r,s\)
both odd; their locations are
\(y_j=(2j+1)\pi s/2\).
Both \(X\) and \(D(X(y)+\ii y)\) have period \(\pi s\).
Each interval between successive \(y_j\)'s is a
translate of the interval containing zero, so
\eqref{geo:collinear-derivative} proves nonidentity
on every such interval.  It remains to exclude
accumulation at its endpoints.  Put
\(\zeta=\ii y_j+\eta\), \(\eta=x+\ii\delta\).
For signs \(\sigma,\tau\in\{1,-1\}\),
\[
 S=\ii\sigma\cosh\eta,\qquad
 T=\ii\tau\cosh(\lambda\eta),
\]
and direct Taylor expansion gives
\begin{equation}\label{geo:odd-odd-expansion}
 D=\sigma\tau(\lambda^2-1)x\delta+O(|\eta|^4).
\end{equation}
Equation \eqref{geo:modulus-graph}, written analytically
in \(x^2,\delta^2\), has linear part
\((1+\lambda^2)(x^2-\delta^2)\).
The implicit function theorem gives
\[
                  X(y_j+\delta)=|\delta|+O(|\delta|^3).
\]
Inserting this into \eqref{geo:odd-odd-expansion}
shows \(D\ne0\) for all sufficiently small
nonzero \(\delta\), since \(\lambda\ne1\).
There is no endpoint accumulation.

On the imaginary axis the candidate condition is
\(\cos y\cos(\lambda y)=0\), a locally finite set.
At a point where the cosine zeros are not simultaneous,
\(|p|=|\sin y\sin(\lambda y)|<1\); continuity and
\eqref{geo:collinear-equality} exclude nearby off-axis
collinear candidates.  The simultaneous case was just
handled.  Reflection \(x\mapsto-x\), leaving \(y\)
fixed, supplies the same conclusion on the other side.
The complete set of collinear candidates is consequently
locally finite in the \(\zeta\)-plane.

The remaining exceptions are also locally finite.
The zeros of \(S,T\), and of \(p^2-1\), are discrete.
Each of the four functions in \eqref{geo:slice-edge-maps}
is nonconstant meromorphic: at zero its leading term
is \(1/(\lambda\zeta)\) or \(1/\zeta\).
Its poles and critical points are locally finite.
Let \(E_\lambda\) be the union of these sets and the
collinear candidates.

At a boundary point outside \(E_\lambda\), the unique
representing coefficients cannot both be strict:
otherwise continuity would place a whole neighborhood
inside the spectral set, contrary to being approached
by the complementary domain.  They cannot both be
endpoints, by the corner identity in
Remark~\ref{geo:algebraic-branches}.  Exactly one is
an endpoint.  Its edge map has nonzero derivative.
For example, at \(a=\epsilon\),
\[
 \Im h_\epsilon
          =\Im(S/T)(a-\epsilon)
\]
for the unique real representing coefficients near
the point; \(\Im(S/T)\ne0\).  Hence the spectral set
is locally one closed side of the nonsingular
real analytic arc \(\Im h_\epsilon=0\), and the
physical domain occupies the other side.

Neither physical boundary has an isolated point:
the HT boundary is the two continuous graphs
\(a=\pm H_\lambda(b)\), completed at their endpoints,
and the LT boundary is the continuous graph
\(a=L_\lambda(b)\).
Removing the locally finite set \(E_\lambda\)
therefore leaves a dense subset of both boundaries.
The edge map is a local biholomorphism and maps each
small regular arc to a real interval.  Prime cosine
density now proves the final assertion.
\end{proof}

\begin{remark}\label{geo:slice-signs}
The regular right boundary of the HT lens has positive
endpoint, and its regular left boundary has negative
endpoint, since \(x\) and \(u=\lambda x\) have the
same sign.  On the LT boundary the endpoint signs
are those in Remark~\ref{geo:endpoint-signs}, with
\(v=\lambda y\).  Thus the negative odd endpoint
problem occurs on the actual HT physical boundary
for every anisotropic positive ratio; it is not
removed by selecting the physical component.
The ratio \(\lambda=1\) is deliberately excluded from
Proposition~\ref{geo:slice-density}; its collinear
degeneracy is the entire isotropic boundary, described
instead by Lemma~\ref{geo:isotropic-cover}.
\end{remark}

\subsection{Exceptional boundary points and componentwise continuation}
\label{geo:continuation-section}

A dense set of obstructions will suffice once nearby points on the
physical side can be joined inside a common local component. The
collar lemma establishes precisely this connectivity, including at
exceptional boundary points.

\begin{lemma}[Uniform local collar]\label{geo:collar}
Let \(D=\DH\) or \(\DL\), and let \(P\in\partial D\).
For every open neighborhood \(U\) of \(P\), there is
an open neighborhood \(V\subset U\) of \(P\) such
that \(D\cap V\) is contained in a single connected
component of \(D\cap U\).
Consequently every component \(E\) of \(D\cap U\)
with \(P\in\overline E\) contains \(D\cap V\).
The same assertion holds for every
\(\Omega_{H,\lambda}\) and \(\Omega_{L,\lambda}\).
\end{lemma}

\begin{proof}
It suffices to work in a small ball contained in \(U\).
Use real coordinates \((x,u,y,v)\), and first consider
\(\DH\).  Its two nontrivial defining functions in the
central strips are
\[
 g_1=\cos y-\cosh x\tanh|u|,\qquad
 g_2=\cos v-\cosh u\tanh|x|.
\]
At \(P=(x_0,u_0,y_0,v_0)\), use the fixed direction
\(\nu=-P\).  For an active inequality \(g_1(P)=0\),
its derivative in that direction is
\begin{equation}\label{geo:inward-derivative}
 y_0\sin y_0+
 x_0\sinh x_0\tanh|u_0|+
 |u_0|\cosh x_0\operatorname{sech}^2u_0>0.
\end{equation}
All terms are nonnegative in the closed central strips.
If the last term vanishes then \(u_0=0\), and activity
forces \(|y_0|=\pi/2\), making the first term positive.
If \(u_0\ne0\), the last term is positive.
Although \(|u|\) is not differentiable at zero, when
\(u_0=0\) the chosen direction has zero \(u\) component.
The directional derivative along \(\nu\) is then
well-defined and continuous in a neighborhood of \(P\).
The exchanged calculation applies to an active \(g_2\).
Thus there is a fixed small \(t>0\) and a neighborhood
of \(P\) on which every active inequality increases
uniformly along \(Q+s\nu\), \(0\leq s\leq t\).
Inactive inequalities remain positive after shrinking
this neighborhood and \(t\).  The central strip
conditions are preserved by the same translation
for nearby points and sufficiently small \(t\);
at a strip endpoint it moves strictly inward.
It follows that \(P+t\nu\in D\) and, for all
sufficiently near \(Q\in D\), the entire segment
\(Q+[0,t]\nu\) lies in \(D\cap U\).

For \(\DL\), use instead the direction
\(\nu=(1,1,0,0)\) and the defining functions
\[
 g_1=\cosh x\tanh u-|\cos y|,\qquad
 g_2=\cosh u\tanh x-|\cos v|.
\]
On \(x,u\geq0\), their derivatives in this direction are
\[
 \sinh x\tanh u+\cosh x\operatorname{sech}^2u>0,
 \qquad
 \sinh u\tanh x+\cosh u\operatorname{sech}^2x>0.
\]
The same translation also opens \(x,u>0\).
The absolute cosine terms are fixed during the
translation, so they create no differentiability issue.
This proves the same segment property.

In either case choose a small ball about \(P+t\nu\)
contained in \(D\cap U\).  Shrink \(V\) about \(P\)
until every translated endpoint \(Q+t\nu\), \(Q\in D\cap V\),
lies in that ball.  Join any two such endpoints by a
straight segment in the ball and use their translation
segments to join the original points.  Thus \(D\cap V\)
lies in one component of \(D\cap U\).
A component accumulating at \(P\) meets \(D\cap V\)
and hence is that component.

On a fixed-ratio slice the HT direction \(-P\)
preserves the slice; the same derivative proof applies.
For LT use the real direction
\((1,\lambda,0,0)\), or equivalently the horizontal
translation \(\zeta\mapsto\zeta+s\) of the epigraph
\(a>L_\lambda(b)\).  The preceding ball-and-segment
argument is then entirely within the complex line.
\end{proof}

\begin{definition}\label{geo:componentwise-extension}
For \(f\in\mathcal O(D)\) and \(P\in\partial D\), a local
holomorphic continuation through \(P\) means a
holomorphic function \(F\) on an open neighborhood \(U\)
of \(P\) which agrees with \(f\) on some connected
component \(E\) of \(D\cap U\) satisfying
\(P\in\overline E\).
This definition allows agreement on only the component
reached by a continuation path; it does not require
agreement on every component of \(D\cap U\).
\end{definition}

\begin{corollary}[Propagation from dense genuine obstructions]
\label{geo:propagate-obstruction}
Let \(D\) be any of the domains in
Lemma~\ref{geo:collar}, and let \(f\in\mathcal O(D)\).
Suppose a dense subset \(G\subset\partial D\) consists
of points at which no holomorphic function on a
neighborhood can agree with \(f\) on the full
intersection with \(D\).  Then \(f\) has no local
holomorphic continuation through any point of
\(\partial D\), even in the componentwise sense of
Definition~\ref{geo:componentwise-extension}.
\end{corollary}

\begin{proof}
If \(F,U,E\) gave such a continuation at \(P\),
Lemma~\ref{geo:collar} would give \(V\subset U\)
with \(D\cap V\subset E\), and therefore \(F=f\)
on \(D\cap V\).  Choose \(Q\in G\cap V\) and a
neighborhood of \(Q\) contained in \(V\).  Restricting
\(F\) gives the forbidden continuation at \(Q\).
\end{proof}

\begin{remark}\label{geo:scope-nonextension}
Proposition~\ref{geo:joint-density},
Proposition~\ref{geo:slice-density}, and
Corollary~\ref{geo:propagate-obstruction} are a geometric
bridge.  To apply them to a susceptibility, the dense
obstructions must be proved for the entire physical sum,
with its normalization and tail included.
For a fixed-ratio corollary, its local singularity
calculation must remain valid along the corresponding
analytic curve \(B=B(t)\) in the coordinates of
Proposition~\ref{geo:local-coordinates}; a bare joint
nonextension assertion does not imply the same for
its restriction.  Nothing here claims that all of
\(\Sigma\), which has interior, is a singularity set.
\end{remark}

\subsection{Pseudoconvexity of the spectral components}
\label{geo:psh-section}

\begin{proposition}[An explicit plurisubharmonic exhaustion]
\label{geo:pseudoconvex}
Every component of \(\C^2\setminus\Sigma\) has a continuous
plurisubharmonic exhaustion.  In particular the explicit
physical components above are pseudoconvex.
\end{proposition}

\begin{proof}
Set
\[
 m(K)=\min_{(a,b)\in[-1,1]^2}|A-Sa-Tb|.
\]
Compactness of the square makes \(m\) continuous, and
its zero set is exactly \(\Sigma\).
On a complementary component, each
\(-\log|A-Sa-Tb|\) is pluriharmonic, and
\[
 -\log m(K)=
       \sup_{(a,b)\in[-1,1]^2}-\log|A-Sa-Tb|
\]
is continuous and locally bounded above.
The submean inequality on every complex line, or the
usual locally bounded supremum criterion, shows that
it is plurisubharmonic.  Consequently
\[
 \Phi(K)=|K_x|^2+|K_y|^2+
                      \max\{0,-\log m(K)\}
\]
is continuous and plurisubharmonic.
It tends to infinity at every finite boundary point,
because \(m\to0\), and at infinity by its quadratic
term.  More explicitly, any sequence in a bounded
sublevel set has a convergent subsequence with
\(m\) bounded away from zero; its limit lies in the
same complementary component.  Thus its sublevel
sets are relatively compact there.
\end{proof}

\begin{remark}\label{geo:Hartogs-scope}
The exhaustion establishes pseudoconvexity of the complementary
components and excludes a compulsory Hartogs filling caused by
failure of pseudoconvexity. Nonextension of the physical function
is established separately by the whole-sum obstructions and the
local continuation argument above.
\end{remark}

\section{The physical observable and joint holomorphy}
\label{ana:physical}

We now identify the function to which the boundary estimates will
apply. The order is physical correlations, complete lattice sums,
joint holomorphic continuation, and finally the energy-field
normalization. Keeping these identifications together makes the
onsite contribution and the ordered-phase subtraction explicit.

Throughout this paper \(z=2K_x\) and \(w=2K_y\) are independent
complex coordinates. Write
\[
 S=\sinh z,\qquad T=\sinh w,\qquad
 A=\cosh z\cosh w,\qquad p=ST.
\]
No inverse temperature is introduced as an additional independent
variable in \(\C^2\). For real positive \(K_x,K_y\), the
dimensionless uniform field is denoted by \(\eta\). The zero-field
responses are
\begin{align}
 \XH(K_x,K_y)
 &=\sum_{(m,l)\in\Z^2}\langle\sigma_{0,0}\sigma_{m,l}\rangle,
 &&0<p<1, \label{ana:physicalHT}\\
 \XL(K_x,K_y)
 &=\sum_{(m,l)\in\Z^2}
   \bigl(\langle\sigma_{0,0}\sigma_{m,l}\rangle_+-M^2\bigr),
 &&p>1, \label{ana:physicalLT}\\
 M^2&=(1-p^{-2})^{1/4}. \nonumber
\end{align}
The high-temperature state is the unique Gibbs state. The subscript
\(+\) means the infinite-volume plus state; the fourth root in the last
line is positive. These definitions include the origin, both axes and
all quadrants. In particular the low-temperature onsite contribution
is \(1-M^2\), not \(1\).

\subsection{Real-temperature form factors and state identification}
\label{ana:interface}

Only the real, noncritical spin form factors are imported from the
exact solution. The relevant finite-cylinder formulas at arbitrary
separation are \cite[(35)--(39)]{BL}; their finite-lattice matrix-element
factorization is proved in \cite[Theorem 6.1]{IL}, with the
paramagnetic vacuum change in its Remark 6.2. In our convention put
\[
 B=T/S>0,\quad C=A/S,\quad
 \cosh\gamma(q)=C-B\cos q,\quad \gamma(q)>0.
\]
Then the infinite-plane coefficient is
\begin{equation}\label{ana:ff}
 f_N(m,l)=\frac{B^{N^2/2}}{N!}
 \int_{\T^N}\prod_{i=1}^N
 \frac{e^{-|m|\gamma(q_i)+\ii lq_i}}{\sinh\gamma(q_i)}
 \prod_{i<j}
 \left\{\frac{\sin((q_i-q_j)/2)}
 {\sinh((\gamma(q_i)+\gamma(q_j))/2)}\right\}^{\!2}
 \prod_{i=1}^N\frac{\dd q_i}{2\pi}.
\end{equation}
With positive real square roots the two correlations are
\begin{align}
 \langle\sigma_0\sigma_{m,l}\rangle
 &=\xi_H\sum_{\substack{N\ge1\\N\ {\rm odd}}}f_N(m,l),
 &\xi_H&=(1-p^2)^{1/4}p^{-1/2},\label{ana:corHT}\\
 \langle\sigma_0\sigma_{m,l}\rangle_+
 &=\xi_L\left(1+\sum_{\substack{N\ge2\\N\ {\rm even}}}f_N(m,l)\right),
 &\xi_L&=(1-p^{-2})^{1/4}=M^2.\label{ana:corLT}
\end{align}
This is the normalization of \cite[(1.8)--(1.10)]{ONGP}.
In particular the anisotropy exponent in \eqref{ana:ff} is \(N^2/2\).

Two identifications connect these formulas with the infinite-plane
observable: the particle sum must pass through the cylinder limit,
and the limiting correlations must belong to the specified phase. On a real noncritical compact set the dispersion has a
uniform positive gap and is analytic in an angular strip. The
finite-circumference corrections \(\eta,\Lambda^{-1},\log\xi_T\)
in \cite{BL} tend uniformly to zero, by the exponentially convergent
trapezoidal rule applied to their sum-minus-integral expressions.
The following uniform estimates allow the particle index to grow
with the circumference. If \(L\) is the circumference, the unnormalized Ramond and
Neveu--Schwarz sums of a strip-analytic function differ by
\(O(Le^{-hL})\): their Fourier expansions select the multiples of
\(L\), with respective factors \(1\) and \((-1)^k\).
Apply this to \(\gamma\) and to
\(\log(1-e^{-\gamma(q)-\gamma(q')})\). The mixed double-grid
difference of \(\log\sinh((\gamma(q)+\gamma(q'))/2)\) in
\cite[(25)]{BL} gives
\[
 |\Lambda_L^{-1}|+\sup_q|\eta_L(q)|\le CLe^{-hL},\qquad
 |\log\xi_{T,L}|\le CL^2e^{-2hL}.
\]
Thus at each fixed separation the entire correction is bounded by
\(C_{m,l}C^N\), uniformly in \(L,N\). Terms with \(N>L\)
vanish on the discrete grid. In particular no estimate at a fixed
particle number is being used to pass an unbounded particle sum
through the thermodynamic limit.
For the real roots \(0<\rho(q)=e^{-\gamma(q)}\le r<1\),
\[
 \left|\frac{\rho(q)-\rho(q')}{1-\rho(q)\rho(q')}\right|\le r<1.
\]
The matching argument in Lemma~\ref{ana:matching} is an algebraic
identity on discrete angular grids as well as on continuous circles.
It gives a common bound
\(\exp\{-cN^2+O(N\log(N+1))\}\) at fixed separation, including
the finite-size one-body corrections and factorial normalization.
Consequently the particle sum passes through the cylinder limit;
each fixed-index Riemann sum gives \eqref{ana:ff}.

Every local cylinder limit is a translation-invariant and spin-flip
invariant Gibbs state. The Gibbs-state theorem
\cite[Theorem 1]{Raoufi} applies with the vertex-transitive group
chosen to be translations: the two unequal positive nearest-neighbor
couplings satisfy its nonnegativity, invariance, irreducibility and
summability hypotheses. That state is the equal mixture of the plus
and minus states. Its two-spin correlations equal those of the plus
state. This proves the state identification in \eqref{ana:corLT};
subtracting \(M^2=\xi_L\) removes exactly its zero-particle term.
No response of a zero-mean mixed state is substituted for
\eqref{ana:physicalLT}.

\subsection{Double contours and the full lattice}
\label{ana:double}

The double-contour representation serves two purposes: its spatial
geometric series counts every lattice site, and its two Schur products
give summable bounds uniform on compact coupling sets. We first
check its residues against the real form factors.

Let \(\zeta,\eta\) now denote independent contour variables, and set
\[
 \mathcal D(\zeta,\eta;z,w)=
 A-\frac S2(\zeta+\zeta^{-1})
  -\frac T2(\eta+\eta^{-1}),\qquad
 a(v,v')=\frac{v-v'}{1-vv'}.
\]
If \(X=\prod_i\zeta_i\) and \(Y=\prod_i\eta_i\), the same coefficient
has the symmetric representation
\begin{equation}\label{ana:doublecorr}
 f_N(m,l)=\frac{p^{N/2}}{N!}\oint\cdots\oint
 X^{|m|-1}Y^{|l|-1}
 \prod_{i<j}a(\zeta_i,\zeta_j)a(\eta_i,\eta_j)
 \prod_{i=1}^N
 \frac{\dd\zeta_i\,\dd\eta_i}{(2\pi\ii)^2\mathcal D(\zeta_i,\eta_i)}.
\end{equation}
The circles have radii strictly less than and sufficiently close to
one. This formula can be checked directly on the real physical
region, as follows. At \(\eta=e^{\ii q}\) the inside root is
\(\zeta=\rho=e^{-\gamma(q)}\), and
\[
 \mathcal D_\zeta(\rho,\eta)
   =\frac{S(1-\rho^2)}{2\rho^2},\qquad
 \frac{\sqrt p}{\mathcal D_\zeta(\rho,\eta)}
   =\sqrt B\,\frac{\rho}{\sinh\gamma}.
\]
Subtracting two dispersion equations gives the exact canceled pair
\begin{equation}\label{ana:pair}
 a(\rho_i,\rho_j)a(\eta_i,\eta_j)
 =-B\frac{(\eta_i-\eta_j)^2\rho_i\rho_j}
          {\eta_i\eta_j(1-\rho_i\rho_j)^2}
 =B\left\{\frac{\sin((q_i-q_j)/2)}
       {\sinh((\gamma_i+\gamma_j)/2)}\right\}^{\!2}
 =:P_{ij}.
\end{equation}
The residues therefore give \(B^{N/2}\), and the pairs give
\(B^{N(N-1)/2}\). This proves the full exponent and the normalization
in \eqref{ana:ff}. Reflection of either spatial coordinate is
legitimate; equality under interchanging the separations while holding
unequal couplings fixed is not required.

The exact product of the two spatial geometric series is
\begin{equation}\label{ana:fullsites}
 \sum_{m,l\in\Z}X^{|m|-1}Y^{|l|-1}
 =\frac{(1+X)(1+Y)}{XY(1-X)(1-Y)}
 =\frac1{XY}+\frac{2(X^{-1}+Y^{-1})}{(1-X)(1-Y)}.
\end{equation}
The first term is the onsite coefficient. Denoting by \(T_N\)
the integral with the second numerator but without its factor two
gives \(C_N=f_N(0,0)+2T_N\), where \(C_N\) sums all sites.
The angular form of this complete coefficient is
\begin{equation}\label{ana:angular}
 C_N=\frac1{N!}
 \int_{\sum_iq_i=0\ ({\rm mod}\ 2\pi)}
 G_N(q)\frac{1+Z}{1-Z}
 \prod_{i=1}^{N-1}\frac{\dd q_i}{2\pi},\qquad
 G_N=B^{N/2}\prod_i\frac1{\sinh\gamma_i}\prod_{i<j}P_{ij},
 \quad Z=\prod_i\rho_i.
\end{equation}
There are \(N-1\), not \(N\), independent normalized measures in this
formula; \(q_N=-\sum_{i<N}q_i\) modulo \(2\pi\).

\begin{proposition}[Joint normal convergence]\label{ana:normal}
On each of the two domains \(\DH,\DL\) specified by the independent
spectral geometry, the physical germ has a single-valued jointly
holomorphic continuation. It is given by
\begin{align}
 \XH&=(1-p^2)^{1/4}
       \sum_{\substack{N\ge1\\N\ {\rm odd}}}p^{(N-1)/2}F_N,
       \label{ana:globalHT}\\
 \XL&=(1-p^{-2})^{1/4}
       \sum_{\substack{N\ge2\\N\ {\rm even}}}p^{N/2}F_N,
       \label{ana:globalLT}
\end{align}
where \(F_N\) is the integral \eqref{ana:doublecorr} with its
\(p^{N/2}\) factor removed and with the multiplier
\eqref{ana:fullsites}. The sums converge normally, with every fixed
mixed complex derivative, on compact subsets. The fourth roots are
the binomial branches supplied by \(|p|<1\) on \(\DH\) and
\(|p|>1\) on \(\DL\).
\end{proposition}
\begin{proof}
On a compact parameter neighborhood in \(\C^2\setminus\Sigma\),
the polynomial \(\mathcal D\) has a nonzero minimum modulus on the
product of unit circles. By uniform continuity that gap persists on
a common thin product annulus and a slightly larger parameter
polydisk. Choose fixed radii \(r_x,r_y<1\) in the annulus.
All Schur entries on these circles are bounded by a constant
independent of \(N\). For even \(N\), the Schur Pfaffian identity
and Hadamard's determinant inequality give
\[
 \left|\prod_{i<j}a(\zeta_i,\zeta_j)\right|
 =|\operatorname{Pf}(a(\zeta_i,\zeta_j))|
 \le C^{N/2}N^{N/4}.
\]
For odd \(N\), adjoin a ghost variable \(-1\); its entries are
the removable constant one and the dimension becomes \(N+1\).
Apply the estimate separately to the two contour-variable families.
The one-body factors and contour lengths cost \(C^N\).
Also \(1-X\) and \(1-Y\) are bounded away from zero uniformly in
\(N\), whereas \(X^{-1},Y^{-1}\) cost only \(C^N\).
Thus, after increasing constants,
\begin{equation}\label{ana:compactbound}
 |F_N|\le C_0 C_1^N\frac{N^{(N+1)/2}}{N!}.
\end{equation}
The right side is summable after multiplication by any fixed
exponential in \(N\). Multivariate Cauchy estimates on the larger
polydisk give the same assertion for every fixed mixed derivative.
On overlaps, radii sufficiently close to one can be deformed into
one another inside the common gap annulus; all pair and product
denominators stay nonzero. The local integrals therefore glue.
All powers of \(p\) within \eqref{ana:globalHT}--\eqref{ana:globalLT}
are integers. The binomial fourth roots introduce no monodromy,
and the series prove joint holomorphy, including the HT axes
\(S=0\) or \(T=0\).

On a real physical compact set the same contours in
\eqref{ana:doublecorr} additionally give
\(r_x^{N|m|}r_y^{N|l|}\). Their double geometric sum is bounded
uniformly in \(N\ge1\). Thus the lattice and particle sums are
absolutely interchangeable, and \eqref{ana:fullsites} gives exactly
\eqref{ana:physicalHT}--\eqref{ana:physicalLT}, with the stated
zero-particle subtraction. Equality on a real open physical set
determines the joint germ: iterating the one-variable identity
theorem over a real rectangle proves the required uniqueness.
\end{proof}

\paragraph{Matching the angular and double-contour branches.}
The boundary calculations use angular variables. To identify their
sum with the holomorphic function just constructed, we first specify
the canonical transfer root.
For real \(q\), outside \(\Sigma\) the equation
\[
 S\rho^2-2(A-T\cos q)\rho+S=0
\]
has a unique root of modulus less than one. For \(S\ne0\) it is the
Joukowski inside root
\[
 \rho=W-\sqrt{W^2-1},\qquad W=(A-T\cos q)/S,
 \qquad \sqrt{W^2-1}\sim W\quad(W\longrightarrow\infty).
\]
At \(S=0\) the implicit-function theorem gives its removable value
zero, since \(A-T\cos q\ne0\). This construction is jointly
holomorphic locally in the couplings and in a complex angular
neighborhood. On positive real noncritical parameters,
\((A-S-T)(A+S+T)=(p-1)^2\) shows \(W(q)>1\), so it is the
physical decaying root. We take \(\gamma=-\log\rho\) only on
local patches where \(\rho\ne0\); no global logarithm through a root
zero is asserted.

We next match each normalized angular coefficient with its
double-contour counterpart throughout the physical domain. Indeed
\(1/\sinh\gamma=2\rho/(1-\rho^2)\) and the rational expression
for \(P_{ij}\) in \eqref{ana:pair} contain no logarithm. After
including the physical normalization, their one-body prefactors are
\[
 (1-p^2)^{1/4}\frac{p^{(N-1)/2}}{S^N}
 \quad(N\text{ odd}),\qquad
 (1-p^{-2})^{1/4}\frac{p^{N/2}}{S^N}
 \quad(N\text{ even}).
\]
All powers here are integers. Hence the normalized angular coefficient
is single valued and holomorphic wherever \(S\ne0\) in the
corresponding physical domain; the real angular torus has uniform
root and denominator gaps on compact subsets. The set
\(\DH\setminus\{S=0\}\) is connected: radial contraction enters
a small polydisk without crossing \(z=0\), and that polydisk minus
\(z=0\) is connected. The domain \(\DL\) already has \(S\ne0\).
The identity theorem from the real physical rectangle therefore gives
equality with the double-contour coefficients throughout these sets.
Near a regular edge \(S,T\ne0\), this also fixes the compatible
local odd square root \(\sqrt B=\sqrt p/S\). Normal convergence
of the identified series follows from Proposition~\ref{ana:normal}.
Thus every boundary calculation below is a calculation of the global
physical sum on its actual interior side.

\begin{remark}[The decoupled-axis restriction]\label{ana:axis}
On \(K_y=0\), the lattice splits into independent horizontal chains.
For \(t=\tanh K_x\) near zero its response is
\[
 1+2\sum_{m\ge1}t^m=\frac{1+t}{1-t}=e^{2K_x}.
\]
This agrees with \eqref{ana:globalHT}: \(p=0\), only \(N=1\)
survives, and \(F_1=(A-S)^{-1}=e^{2K_x}\).
The restriction is entire, although the joint domain has boundary
points on this line. This exact example is why a bare joint
nonextension assertion cannot be used to prove the slice part
of the main theorem.
\end{remark}

\subsection{Energy-field normalization}

For fixed real \(J_x,J_y>0\) and physical temperature
\((K_x,K_y)=(\beta J_x,\beta J_y)\), a field of energy units enters
the Hamiltonian as \(-h\sum_v\sigma_v\), and \(\eta=\beta h\).
Finite-volume differentiation gives \(\chi=\beta\mathcal X\).
The same statement holds after taking the prescribed infinite-volume
phase limit. To justify the latter without exchanging an uncontrolled
field limit and infinite sum, use the GHS inequality for arbitrary
nonnegative couplings and fields \cite[Theorem 1.5]{CGN}.
Pinning removed spins by fields tending to \(+\infty\) shows that
plus-box covariances increase with the box. Their nonnegative sums
increase to the finite sum proved above. GHS concavity gives, for
the plus-box magnetization,
\[
 m_V(h)-m_V(0)\le h\,m_V'(0)
       \le\beta h\,\mathcal X,\qquad h\ge0.
\]
The plus-state limit yields \(m'(0+)\le\beta\mathcal X\).
Conversely turn on \(h\) only on a finite set \(E\). The resulting
bounded local change of measure has derivative
\(\beta\sum_{v\in E}\operatorname{Cov}_{+,0}(\sigma_0,\sigma_v)\).
Ferromagnetic monotonicity bounds its magnetization above by the
uniform-field magnetization, with equality at \(h=0\).
Increasing \(E\) proves the reverse inequality. Thus
\(m'(0+)=\beta\mathcal X\). Spin-flip symmetry gives the two-sided
HT derivative; in LT the derivative is the plus-phase right
derivative (or its minus-phase counterpart).
This statement is made on temperature lines only.

\section{Regular spectral edges and transverse curves}
\label{ana:edges}

This section proves the positive-edge obstruction in independent
complex dispersion parameters. A strict contraction controls the
numerators, the higher-index tail remains bounded after differentiation,
and the first resonant coefficient has a nonzero leading amplitude.
Keeping the anisotropy holomorphic along the curve makes the result
applicable to the physical temperature slices.
On a regular edge, exchange the two lattice axes if necessary.
In the local coordinates \(C=A/S,\ B=T/S\) of the geometry section,
one edge is \(C-B=b\), \(b=\cos\phi\in(-1,1)\).
Choose the side for which \(\Im B_0>0\); conjugation handles the
other side. Let \(B(u)\) be any holomorphic function near zero with
\(B(0)=B_0\), and set
\begin{equation}\label{ana:edgecurve}
 W(q;u)=b+u+B(u)(1-\cos q).
\end{equation}
Both a constant-\(B\) complex transversal in \(\C^2\) and each
nonramified physical temperature-line transversal have this form.

At \(u=0\), the inside root lies in the closed lower half disk,
with \(\rho(0)=e^{-\ii\phi}\). Compactness gives
\[
 \Im\rho(q)\le-\delta<0,\qquad |\rho(q)|\le1.
\]
There is no \(W=\pm1\) on the real angular circle. The roots,
their chosen logarithms, the one-body factors, and the canceled
pairs \eqref{ana:pair} are analytic on one fixed complex
neighborhood in \((u,B,q)\). On \(u=\ii\eps\), \(\eps>0\),
\begin{equation}\label{ana:attenuation}
 \Re\gamma(q;\ii\eps)
   \ge c\{\eps+\operatorname{dist}(q,0)^2\}.
\end{equation}
Indeed
\(\Im W=\eps+\Im B(\ii\eps)(1-\cos q)\) has this lower bound;
the identity
\(\Im W=\sinh(\Re\gamma)\sin(\Im\gamma)\) on the compact branch
then gives \eqref{ana:attenuation}. This is the complement side
of the parallelogram edge and hence the side of the bordering
physical component.

\begin{lemma}[Strict contraction on a complex polydisk]
\label{ana:matching}
There are \(\kappa>0\), \(C<\infty\), and a fixed complex parameter
polydisk about this edge such that, for every fixed mixed derivative
\(\partial^\nu\) in the parameters,
\begin{equation}\label{ana:gaussian}
 |\partial^\nu G_N(q)|
 \le C_\nu\exp\{-\kappa N^2+C N\log(N+1)\}
\end{equation}
for all real angles and every \(N\ge1\).
\end{lemma}
\begin{proof}
The identity
\[
 |1-vv'|^2-|v-v'|^2
 =(1-|v|^2)(1-|v'|^2)+4\Im v\,\Im v'
\]
gives a strict bound \(|a(\rho_i,\rho_j)|\le q_0<1\).
The denominator is separated from zero. By continuity a slightly
larger constant below one still works on a common \emph{complex}
parameter polydisk; individual root moduli need not stay below one
there. Every canceled pair \(P_{ij}\) remains bounded by \(C\).

For even \(N\), expand the Schur Pfaffian in the
\(\eta_i=e^{\ii q_i}\) variables only. For each matching, combine
its \(N/2\) edges with the corresponding \(\rho\) factors using
\eqref{ana:pair}. The other \(N(N-2)/2\) root factors retain their
strict contraction. Thus
\[
 \left|\prod_{i<j}P_{ij}\right|
 \le(N-1)!!\,C^{N/2}q_0^{N(N-2)/2}.
\]
For odd \(N\), adjoining the constant ghost entry gives
\[
 \left|\prod_{i<j}P_{ij}\right|
 \le N!!\,C^{(N-1)/2}q_0^{(N-1)^2/2}.
\]
These identities first hold away from reciprocal \(\eta\) pairs and
extend by the analytic canceled formula. There is no division by
a root-Schur zero. One-body factors cost \(C^N\).
Cauchy's inequalities on the larger polydisk prove all the derivative
bounds. They are applied to \(G_N\), never across a pole of
\((1+Z)/(1-Z)\).
\end{proof}

\begin{lemma}[The entire differentiated higher-index tail]
\label{ana:positivetail}
Suppose \(\phi=2\pi a/\ell\), where \(\ell\) is an odd prime and
\(1\le a<\ell/2\). For odd parity let \(n=\ell\) and
\(k=(n^2-1)/2\); for even parity let \(n=2\ell\) and
\(k=n^2/2-1\). For every \(0\le j\le k\),
\[
 \sum_{\substack{N>n\\N\ {\rm of\ the\ chosen\ parity}}}
 |\partial_u^j C_N(\ii\eps)|=O_j(1)
 \qquad(\eps\downarrow0).
\]
All lower indices of the chosen parity are holomorphic near \(u=0\).
\end{lemma}
\begin{proof}
Write \(Q_0=\sum_i\operatorname{dist}(q_i,0)^2\).
Equation \eqref{ana:attenuation} gives
\[
 |1-Z|\ge c\min\{1,N\eps+Q_0\}.
\]
Since all fixed derivatives of \(\gamma\) are bounded, differentiating
the exponential and rational function gives
\[
 \left|\partial_u^r\frac{1+Z}{1-Z}\right|
 \le C_j N^j(1+Q_0^{-j-1}),\qquad 0\le r\le j.
\]
If \(N\ge2j+4\), retain the first \(2j+3\) independent angle
squares in \(Q_0\). Their inverse \((j+1)\)-st power is integrable
in that fixed real dimension; all unused normalized angular measures
have mass one. Therefore Lemma~\ref{ana:matching} gives
\[
 |\partial_u^jC_N(\ii\eps)|
 \le\frac{C_jN^j}{N!}
       \exp\{-\kappa N^2+C N\log(N+1)\},
\]
uniformly in \(\eps\). This covers an actual infinite index range.

For each remaining fixed \(N\), the only configuration with
\(|Z|=1\) is \(q_1=\cdots=q_N=0\). There \(Z=e^{-\ii N\phi}\);
the first resonant odd index is \(\ell\), and the first resonant
even index is \(2\ell\). Lower indices have a denominator gap and
hence are holomorphic. Near the collision of a higher index use
\(q_N=-\sum_{i<N}q_i\), \(Q=\sum_iq_i^2\), and
\(\Delta(q)=\prod_{i<j}(q_i-q_j)\).
Every fixed numerator derivative has the factor \(\Delta^2\).
The worst radial bound for the \(j\)-th derivative is
\[
 r^{N(N-1)}r^{N-2}(\eps+r^2)^{-j-1}\dd r.
\]
Its limiting exponent \(N^2-2j-4\) is greater than \(-1\) for
every larger same-parity \(N\) and \(j\le k\). This handles the
finite intervening range and completes the entire tail estimate.
\end{proof}

\begin{proposition}[Nonzero full-sum edge obstruction on every
holomorphic transversal]\label{ana:positive}
Under \eqref{ana:edgecurve} and the rational-phase hypotheses of
Lemma~\ref{ana:positivetail}, the complete sum of the selected
parity has no holomorphic extension through \(u=0\).
The assertion is locally uniform in independent complex \(B\)
parameters. It remains valid after multiplication by a
holomorphic nonvanishing prefactor and after a nonramified change
of transversal coordinate. It holds for both endpoint signs in
the even case.
\end{proposition}
\begin{proof}
\emph{The first amplitude.} We evaluate the constrained Gaussian
period with its exact measure. Its nonzero value supplies the
leading coefficient that the bounded higher-index tail cannot cancel.
Put \(\alpha=(n^2-1)/2\) and
\[
 P_n=\prod_{r=1}^{n-1}\frac{r!}{2^r}.
\]
The Gaussian Vandermonde period on the \emph{independent-coordinate}
measure is
\begin{equation}\label{ana:mehta}
 \int_{q_n=-\sum_{i<n}q_i}e^{-Q}\Delta(q)^2
               \dd q_1\cdots\dd q_{n-1}
 =\frac{n!\,\pi^{(n-1)/2}P_n}{\sqrt n}=:\mathcal G_n.
\end{equation}
To derive it, expand the ordinary \(n\)-variable Vandermonde
determinant in monic Hermite polynomials. Their Gaussian norms are
\(\sqrt\pi\,r!/2^r\), so that integral is
\(n!\pi^{n/2}P_n\). Setting \(x_i=q_i+m\), \(\sum q_i=0\),
has Jacobian \(n\), and the \(m\) integral is \(\sqrt{\pi/n}\).
Division proves \eqref{ana:mehta}.

On the unique collision chart,
\[
 \gamma(q;u)=\ii\phi-\frac{\ii}{\sin\phi}
       \left(u+\frac{B_0q^2}{2}\right)
       +O(u^2+u q^2+q^4).
\]
The leading singular integral for \(C_n\) is consequently
\begin{equation}\label{ana:firstintegral}
 \frac{2\ii^{1-n^2}B_0^{n^2/2}}
 {n!(2\pi)^{n-1}2^{n(n-1)}\sin^{n^2-1}\phi}
 \int\frac{\Delta(q)^2}{nu+(B_0/2)Q}\,
             \dd q_1\cdots\dd q_{n-1}.
\end{equation}
For odd \(n\) the local square root of \(B\) is the one used in
the physical angular coefficient. Its sign can change the displayed
amplitude's sign but cannot make it zero. For even \(n\) the
power in the original integrand is an integer; the square root
below is fixed by the Gaussian period.

After \(k\) derivatives, the model integral is evaluated from
\[
 \int\frac{\Delta(q)^2}{(a+bQ)^{k+1}}\dd q
 =\mathcal G_n\,b^{-\alpha}
    \frac{\Gamma(k+1-\alpha)}{\Gamma(k+1)}
    a^{\alpha-k-1},
\]
initially in a common right half-plane and then by analytic
continuation to \(a=nu,\ b=B_0/2\).
More explicitly, rotate both coefficients by \(-\ii\):
\(a'=n\eps>0\) and \(b'=-\ii B_0/2\) has positive real part.
The absolutely convergent Laplace--Gaussian calculation for
\(a',b'\) fixes the continuation and its square root. No assumption
\(\Re B_0>0\) is made.
For odd \(n\), \(k=\alpha\) is a multiple of four; for even
\(n\), \(k=\alpha-\tfrac12\) is odd. Substitution gives
\begin{align}
 \lim_{\eps\downarrow0}u\,\partial_u^k C_n(u)
 &=\frac{2^{1-k}\sqrt{B_0}\,P_n n^{k-3/2}}
         {\pi^{(n-1)/2}\sin^{2k}\phi}\ne0,
 &&u=\ii\eps,\quad n\ {\rm odd},\label{ana:oddamp}\\
 \lim_{\eps\downarrow0}u^{1/2}\partial_u^k C_n(u)
 &=-\ii\,\frac{2^{1-\alpha}\sqrt{B_0}\,P_n n^{k-1}}
         {\pi^{(n-2)/2}\sin^{n^2-1}\phi}\ne0,
 &&u=\ii\eps,\quad n\ {\rm even}.\label{ana:evenamp}
\end{align}
In the second line \(u^{1/2}=e^{\ii\pi/4}\sqrt\eps\).

\emph{Comparison with the exact integral.} We now justify these
limits for the actual fixed-\(n\) coefficient by controlling the
difference from its quadratic model. On its collision chart write the exact integrand as
\(\Delta^2 f(q,u)/H(q,u)\). It is even under simultaneous
\(q\mapsto-q\), and, with \(c=-\ii/\sin\phi\),
\[
 H=c\{nu+(B_0/2)Q\}+O((|u|+Q)^2),\quad
 H_u=cn+O(|u|+Q),\quad f=f_0+O(|u|+Q).
\]
Both the exact and model denominators have modulus at least
\(c_1(\eps+Q)\) on \(u=\ii\eps\). The top derivative has
leading term \((-1)^kk! fH_u^k/H^{k+1}\); all other terms have
denominator power at most \(k\). Subtracting the constant quadratic
model therefore gives the pointwise bound
\[
 |\partial_u^k(\hbox{exact integrand}-\hbox{model integrand})|
       \le C\Delta^2(\eps+Q)^{-k}.
\]
Its radial integral is \(O(\log(1/\eps))\) for odd \(n\), where
\(k=\alpha\), and \(O(1)\) for even \(n\), where \(k=\alpha-1/2\).
Extending the differentiated model from the chart to the whole
tangent plane costs \(O(1)\), since its radial tails are
\(R^{-3}\) and \(R^{-2}\), respectively. Outside the chart the
exact denominator has a gap. These bounds prove the two claimed
limits, including for the actual holomorphic function \(B(u)\),
whose contribution to the denominator remainder is \(O(uQ)\).

\emph{Passage to the full sum.}
Lemma~\ref{ana:positivetail} bounds all higher derivatives in the
entire tail, and lower indices are analytic. Lower derivatives of
the even first term are bounded; the next lower odd derivative
may be logarithmic, but all are \(o(u^{-1})\).
Thus the nonzero top amplitudes cannot be canceled in the full
sum or by product-rule terms from an analytic nonvanishing
physical prefactor. The chain rule proves the statement for every
nonramified curve parameter.

At the negative endpoint replace \(q_i\) by \(\pi+q_i'\).
For even \(N\) the total momentum constraint is unchanged and
\(B^{N^2/2}=(-B)^{N^2/2}\). Apply the preceding proof with
\(-B\), conjugating if necessary to have positive imaginary part.
Exchanging transfer axes handles the other two edges.
For odd \(N\) this shift instead makes the total reduced angle
an odd multiple of \(\pi\); that distinct case is addressed by
the full boundary-jet theorem in the next section.
\end{proof}

\section{The full odd sum at a regular negative spectral edge}
\label{neg:section}

At a negative edge, the odd momentum constraint prevents the
all-endpoint configuration at any fixed particle number. Each inward
derivative of the complete sum therefore has a finite boundary value.
The obstruction appears in their growth: selected Taylor coefficients
grow too rapidly to belong to a holomorphic function.

We prove this for an arbitrary holomorphic anisotropy along a
transverse analytic curve, retaining the exact particle-dependent
anisotropy factor, the complete constrained angular integral and every
odd particle number. The proof reduces the coefficients to pole
periods, evaluates the first two windings, selects one dominant index,
and bounds all remaining indices and windings.

\subsection{Statement, branch, and the physical side}
\label{neg:setup}

Fix
\begin{equation}
 \phi=\frac{2\pi r}{p}\in(0,\pi),\qquad
 (r,p)=1,\qquad p\ge3\ \hbox{odd},\qquad
 b=\cos\phi,\qquad s_\phi=\sin\phi>0 .
 \label{neg:rational}
\end{equation}
Let \(B\) be holomorphic on a disk about zero and suppose that
\begin{equation}
 \Im B(0)<0 .
 \label{neg:orientation}
\end{equation}
Shrink this disk so that \(B\ne0\) and \(\Im B<0\) throughout it, and
fix one holomorphic logarithm of \(B\).  For \(q=\pi+\delta\), define
\begin{equation}
 W(t,\pi+\delta)
   =b+B(t)\bigl(t-1+\cos\delta\bigr),\qquad
 g(t,\delta)=\operatorname{arcosh}W(t,\pi+\delta),\qquad
 g(0,0)=\ii\phi .
 \label{neg:dispersion}
\end{equation}
At \(t=0\), \(W\) is in the open upper half-plane except at
\(\delta=0\bmod 2\pi\), where it equals \(b\in(-1,1)\).  Consequently
the branch in \eqref{neg:dispersion} extends to a fixed complex
neighborhood of all real angles and a fixed parameter disk.  It is
even and \(2\pi\)-periodic in \(\delta\).  Put
\(\gamma(t,q)=g(t,q-\pi)\) and \(\rho(t,q)=e^{-\gamma(t,q)}\).
At \(t=0\), \(|\rho|\le1\), and the compact set of these roots lies
strictly in the lower half-plane:
\begin{equation}
 \Im\rho(0,q)\le-\delta_0<0 .
 \label{neg:root-compact}
\end{equation}
The logarithm used for \(\gamma\) is fixed by
\(0<\Im\gamma(0,q)<\pi\).  In particular all the one-body
denominators and the canceled pair denominators below are nonzero on
a fixed complex neighborhood.

For odd \(n\ge1\), let
\[
 \mathcal T_n=\{(q_1,\ldots,q_n)\in(\R/2\pi\Z)^n:
                         q_1+\cdots+q_n=0\bmod 2\pi\},
 \qquad q_n=-\sum_{i<n}q_i\bmod 2\pi .
 \]
The measure on this torus is \(\dd q_1\cdots\dd q_{n-1}\).  Include
the normalization in the numerator:
\begin{align}
 H_n(t,q)
 &=\frac{B(t)^{n^2/2}}{n!(2\pi)^{n-1}}
   \prod_{i=1}^n\frac1{\sinh\gamma(t,q_i)}
   \prod_{i<k}
   \left[
    \frac{\sin((q_i-q_k)/2)}
         {\sinh((\gamma(t,q_i)+\gamma(t,q_k))/2)}
   \right]^2 ,
 \label{neg:numerator}\\
 \Gamma_n(t,q)&=\sum_{i=1}^n\gamma(t,q_i),\\
 C_n(t)&=\int_{\mathcal T_n}H_n(t,q)
                          \coth\!\left(\frac{\Gamma_n(t,q)}2\right)
                          \dd q .
 \label{neg:coefficient}
\end{align}
The power of \(B\) uses the chosen common logarithm.  Changing its
square-root sign changes every odd coefficient by the same sign.

The relevant open side of the boundary is
\begin{equation}
 \mathcal U^-=\{\,|t|<r_0:\ \Im(B(t)t)>0\,\}.
 \label{neg:physical-side}
\end{equation}
Here
\[
 \Im W(t,\pi+\delta)
   =\Im(B(t)t)-\Im B(t)(1-\cos\delta)>0
 \]
for every real angle, so the chosen roots satisfy \(|\rho|<1\).
This is the local physical side when the regular negative endpoint
has the lower-half-plane root \(\rho(0,\pi)=e^{-\ii\phi}\).
It contains \(t=-\epsilon\) for every sufficiently small
\(\epsilon>0\).  When these coordinates are obtained from
\(W(q)=C-B\cos q\), they are exactly
\[
 t=\frac{C+B-b}{B}.
 \]
The curve hypothesis is that this \(t\) is a nonramified local
parameter.  Thus any transverse analytic curve through the selected
endpoint is included, with its actual function \(B(t)\).  The opposite
orientation is obtained by complex conjugating all the data.

\begin{theorem}[Universal analytic-curve obstruction at a negative edge]
\label{neg:theorem}
Under \eqref{neg:rational}--\eqref{neg:orientation}, the series
\[
 F(t)=\sum_{\substack{n\ge1\\n\ {\rm odd}}}C_n(t)
 \]
converges normally on \(\mathcal U^-\).  Every inward derivative has
a finite limit
\begin{equation}
 a_j=\lim_{\epsilon\downarrow0}\frac{F^{(j)}(-\epsilon)}{j!},
 \qquad j=0,1,2,\ldots .
 \label{neg:jets}
\end{equation}
Nevertheless its formal boundary Taylor series has radius zero.
More precisely, as \(n\) tends to infinity through odd multiples of
\(p\), there are integers \(j_n\) such that
\begin{align}
 j_n&=n^2\left(\log n+\tfrac12\log\log n-c_*+o(1)\right),
 \label{neg:jn-statement}\\
 |a_{j_n}|^{1/j_n}
   &\sim \frac{4j_n}{e\pi^2\log j_n},
 \label{neg:growth-statement}\\
 c_*&=-\frac34-\frac32\log2+\log\pi
            +\frac12\log|B(0)|-\log s_\phi .
 \label{neg:cstar}
\end{align}
It follows that \(F\) has no holomorphic extension through \(t=0\)
from \(\mathcal U^-\).  If \(P\) is holomorphic near zero and
\(P(0)\ne0\), the same nonextension conclusion holds for \(P F\).

The hypotheses and all estimates are locally uniform for complex
holomorphic families \(B(t,\eta)\): it suffices to work on a common
parameter disk and a relatively compact parameter polydisk on which
\(\Im B(0,\eta)\le-\beta<0\) and \(B\) has a common holomorphic bound.
The constants can depend on that disk, that bound, \(\beta\), and the
 fixed rational phase.  The selected integers are allowed to depend
 on the family parameter, \(j_n=j_n(\eta)\); a common sequence of
 selected integers for all \(\eta\) is not asserted.
\end{theorem}

We prove the theorem by identifying a single particle index that
dominates selected, increasingly high Taylor coefficients in modulus.
The other indices have arbitrary complex phases; no positivity or
phase alignment of the infinite sum is assumed.

\subsection{A uniform Gaussian numerator bound and all boundary jets}
\label{neg:jets-section}

There are two estimates to establish first. Gaussian decay in the
particle number makes each fixed derivative summable. The odd
momentum constraint gives a positive gap for each individual particle
integral, although that gap shrinks as the particle number increases.

Write \(\mathcal A(v,w)=(v-w)/(1-vw)\) and \(y_i=e^{\ii q_i}\).
Subtraction of two dispersion equations gives the canceled identity
\begin{equation}
 \mathcal A(\rho_i,\rho_k)\mathcal A(y_i,y_k)
 =B(t)\left[
    \frac{\sin((q_i-q_k)/2)}
         {\sinh((\gamma_i+\gamma_k)/2)}
       \right]^2 .
 \label{neg:canceled}
\end{equation}
This is an identity of holomorphic canceled functions, including
reciprocal \(y\)-pairs.  The elementary identity
\begin{equation}
 |1-vw|^2-|v-w|^2
 =(1-|v|^2)(1-|w|^2)+4\Im v\Im w
 \label{neg:schur}
\end{equation}
and \eqref{neg:root-compact} imply
\(|\mathcal A(\rho_i,\rho_k)|\le\theta<1\) at \(t=0\).
The same strict bound, with a slightly larger \(\theta<1\), holds
on a fixed complex parameter disk for all real angles.  Every
complete canceled pair in \eqref{neg:canceled} is bounded there.

For completeness, the Schur identity used here is
\(\operatorname{Pf}(\mathcal A(y_i,y_k))_{i,k=1}^{N}
=\prod_{i<k}\mathcal A(y_i,y_k)\) for even \(N\).
It follows by induction in \(N\): as rational functions of \(y_N\),
the two sides have the same simple residues at \(y_N=1/y_k\).
In checking the residues, the other paired factors cancel by
\(\mathcal A(a,b)\mathcal A(a,1/b)=1\), and the remaining product
is the induction hypothesis.  Their difference is constant because
both are bounded at infinity, and that constant is zero at
\(y_N=y_{N-1}\).  The identity then extends from generic tuples
by rational continuation.

Apply this identity only to the \(y\)-variables, adjoining the
ghost variable \(-1\) because \(n\) is odd.  The entry
\(\mathcal A(y_i,-1)\) is the removable constant one.
Each of the \(n!!\) matchings has \((n-1)/2\) real edges.
Combine each matched \(y\)-factor with the corresponding
\(\rho\)-factor by \eqref{neg:canceled}.  The remaining
\((n-1)^2/2\) \(\rho\)-factors are at most \(\theta\) in modulus.
Thus, with the one-body factors and normalization included,
\begin{equation}
 \sup_{|t|<r_0,\ q\in\mathcal T_n}|H_n(t,q)|
 \le E_n,\qquad
 E_n=\exp\{-\kappa n^2+C n\log(n+1)+C\}
 \label{neg:gaussian}
\end{equation}
for some \(\kappa>0\).  Initially the Pfaffian calculation can be
made off its apparent reciprocal-pair divisors; the canceled
identity extends the estimate to those divisors.  There is no
division by a possibly zero \(\rho\)-factor.  The same estimate
holds when every angle is in a sufficiently small fixed complex
neighborhood of \(\pi\).  The latter version will control the
deformation of the local integration cycle.

Choose shortest lifts \(q_i=\pi+\delta_i\bmod2\pi\), with
\(\delta_i\in[-\pi,\pi]\), and put
\[
 S(q)=\sum_{i=1}^n\delta_i^2 .
 \]
The odd momentum constraint says
\(\sum_i\delta_i=w\pi\), where \(w\) is odd.  Therefore
\begin{equation}
 S(q)\ge\frac{\pi^2}{n}.
 \label{neg:momentum-gap}
\end{equation}
On the inward ray, compactness of the branch and
\(\Im W(-\epsilon,\pi+\delta)\ge
c(\epsilon+\operatorname{dist}(\delta,2\pi\Z)^2)\) imply
\begin{equation}
 \Re g(-\epsilon,\delta)
 \ge c\bigl(\epsilon+\operatorname{dist}(\delta,2\pi\Z)^2\bigr).
 \label{neg:attenuation}
\end{equation}
Indeed \(\Im W=\sinh(\Re g)\sin(\Im g)\), and the relevant roots
remain in a fixed compact set.  All parameter derivatives of \(g\)
are bounded on a fixed complex neighborhood.

It follows from \eqref{neg:momentum-gap} that \(C_n\) is holomorphic
on \(|t|<c/n^2\).  On this disk,
\(\Re\Gamma_n\ge c/n\) after decreasing \(c\), and hence
\(|\coth(\Gamma_n/2)|\le Cn\).  The torus volume is absorbed in
the linear exponential in \eqref{neg:gaussian}.  Cauchy's estimate
therefore gives
\begin{equation}
 |[t^j]C_n(t)|\le (Cn^2)^j\,CnE_n .
 \label{neg:small-disk}
\end{equation}
More generally, the disk of radius \(c(\epsilon+n^{-2})\) about
\(-\epsilon\) has
\(\Re\Gamma_n\ge c(n\epsilon+n^{-1})\).  The same argument bounds
the \(j\)th derivative at \(-\epsilon\) by
\(j!(Cn^2)^j CnE_n\), uniformly for small \(\epsilon\).
On compact subsets of \(\mathcal U^-\), the kernel is bounded and
\eqref{neg:gaussian} proves normal convergence.  Dominated
convergence for each fixed \(j\) now proves
\begin{equation}
 a_j=\sum_{\substack{n\ge1\\n\ {\rm odd}}}[t^j]C_n(t),
 \qquad
 \sum_n |[t^j]C_n(t)|<\infty .
 \label{neg:jet-sum}
\end{equation}
This establishes the existence of all the jets, but does not give
a positive radius for their formal Taylor series.

\subsection{An exact reduction to simple-pole periods}
\label{neg:pole-reduction}

The shrinking gaps are encoded by nearby simple poles. We isolate
their locations exactly and place the remaining analytic terms in a
single remainder, whose coefficient growth is small enough for the
later comparison.

Fix a sufficiently small \(\sigma>0\).  There is a fixed \(c_1>0\)
such that on \(|t|<c_1/n\) the part with \(S\ge\sigma\) is analytic
with a bounded kernel.  To see this, use
\(\Re\Gamma_n(0,q)\ge cS\) and the perturbation bound
\(|\Gamma_n(t,q)-\Gamma_n(0,q)|\le Cn|t|\).
Choose \(c_1\) to be a sufficiently small multiple of \(\sigma\).

In the region \(S<\sigma\), set
\[
 D_n(t,q)=\Gamma_n(t,q)-n\ii\phi .
 \]
The Taylor expansion of \(g\) at \((0,0)\) gives
\(|D_n(t,q)|\le C(S+n|t|)\).
If \(p\nmid n\), the distance of \(n\phi\) from \(2\pi\Z\) is at
least \(2\pi/p\).  Reduce \(\sigma,c_1\) so that the perturbation
is less than \(\pi/p\).  This part is also analytic, with bounded
kernel, on \(|t|<c_1/n\).

It remains to consider odd multiples of \(p\) and \(S<\sigma\).
Here \(n\phi\in2\pi\Z\), and
\[
 \coth(\Gamma_n/2)=\frac2{D_n}+\mathcal R(D_n),
 \]
where \(\mathcal R\) is holomorphic on a fixed disk containing all
these values for \(|t|<c_1/n\).
There is a unique small zero \(t=T_n(q)\) of \(D_n(t,q)\).
This statement has a fixed-disk quantitative form.  With
\(a_0=B(0)/(\ii s_\phi)\ne0\), one has
\[
 \frac1n\partial_tD_n(t,q)=a_0+O(|t|+\sigma)
 \]
on a fixed disk.  By reducing that disk and \(\sigma\), its
distance from \(a_0\) is less than \(|a_0|/4\).  Rouch\'e's
theorem on a circle of radius \(CS/n\), followed by uniqueness
from this derivative estimate, gives
\begin{equation}
 |T_n(q)|\le CS/n,\qquad
 D_n(t,q)=n(t-T_n(q))A_n(t,q),
 \label{neg:pole-factor}
\end{equation}
where \(A_n\) is holomorphic on a fixed smaller disk and bounded
away from zero.  One can take
\[
 A_n(t,q)=\frac1n\int_0^1
   \partial_tD_n(T_n(q)+u(t-T_n(q)),q)\,\dd u .
 \]
All such disks and bounds are independent of \(n\) and \(q\).

Put \(J_n=H_n/A_n\).  The exact identity
\begin{equation}
 \frac{2H_n(t,q)}{D_n(t,q)}
 =\frac{2J_n(T_n(q),q)}{n(t-T_n(q))}
  +\frac2n
     \frac{J_n(t,q)-J_n(T_n(q),q)}{t-T_n(q)}
 \label{neg:division}
\end{equation}
separates the pole from an analytic divided difference.
The second term is bounded by \(CE_n\), by fixed-disk Cauchy
estimates.  The \(\mathcal R\) term has the same bound on
\(|t|<c_1/n\).  Combining these observations with the exterior
and nonresonant regions gives
\begin{align}
 a_j&=
 \sum_{\substack{n\ {\rm odd}\\p\mid n}}A^{\rm pole}_{n,j}+R_j,
 \label{neg:all-poles}\\
 A^{\rm pole}_{n,j}
 &=-\int_{S<\sigma}
   \frac{2H_n(T_n(q),q)}
        {\partial_t\Gamma_n(T_n(q),q)}
   T_n(q)^{-j-1}\,\dd q ,
 \label{neg:pole-coefficient}\\
 |R_j|
 &\le \sum_{n\ge1}(Cn)^j
             \exp\{-\kappa n^2+C n\log(n+1)+C\}\\
 &\le
 \exp\{\tfrac12 j\log(j+2)+Cj+C\}.
 \label{neg:regular-remainder}
\end{align}
For each \(n\), the coefficient identity is first justified on
the smaller disk \(|t|<c/n^2\); no pole is crossed there.
The last inequality follows by absorbing \(Cn\log(n+1)\) into
half the Gaussian and maximizing \(j\log n-\kappa n^2/2\).
This reduction isolates all small pole locations before any local
change of angular variables is made.

\subsection{Bounds for every winding and every particle index}
\label{neg:all-windings}

Before evaluating the leading windings, we need estimates valid for
every winding and particle index. These bounds will control the
parts of the infinite sum that lie outside the asymptotic window.

The zero in \eqref{neg:pole-factor} has a useful more precise
estimate.  Set
\[
 v_i=1-\cos\delta_i,\qquad \bar v=\frac1n\sum_i v_i .
 \]
At \(t=\bar v\), the linear terms in the expansion of
\(\operatorname{arcosh}(b+B(\bar v)(\bar v-v_i))\) cancel on
averaging.  The remaining average is bounded by
\(C\,n^{-1}\sum_i(v_i-\bar v)^2\).  The uniform nonzero
derivative used in \eqref{neg:pole-factor} therefore yields
\[
 T_n=\bar v+
       O\!\left(\frac1n\sum_i(v_i-\bar v)^2\right).
 \]
Since \(v_i\ge0\), \(\max_i v_i\le C\sigma\), and
\(\bar v=(S/(2n))(1+O(\sigma))\), this implies
\begin{equation}
 |T_n(q)|\ge(1-C\sigma)\frac{S}{2n},
 \qquad
 |\partial_t\Gamma_n(T_n(q),q)|\ge cn .
 \label{neg:pole-lower}
\end{equation}

Retain the exact central numerator size
\begin{equation}
 H_{0,n}=
 \frac{|B(0)|^{n^2/2}}
 {n!(2\pi)^{n-1}2^{n(n-1)}s_\phi^{n^2}} .
 \label{neg:H0}
\end{equation}
Dividing the sine factors by their linear terms and all the
hyperbolic sine factors by \(\ii s_\phi\) gives
\begin{equation}
 |H_n(T_n(q),q)|
 \le H_{0,n}\Delta(\delta)^2e^{CnS},
 \qquad
 \Delta(\delta)=\prod_{i<k}(\delta_i-\delta_k).
 \label{neg:pole-density-upper}
\end{equation}
Here is the relevant accounting.  The logarithm of the anisotropy
ratio costs \(Cn^2|T_n|\le CnS\).  The sum of the absolute
one-body phase deviations is at most \(CS\), so all pair
denominators together cost \(CnS\).  Finally
\(\sum_{i<k}|\delta_i-\delta_k|^2\le CnS\) controls the sine
divided differences.  Thus no undetermined constant to the power
\(n^2\) replaces the displayed normalization.

For an odd winding \(w\), write
\begin{equation}
 \delta_i=\frac{w\pi}{n}+\eta_i,\qquad
 \sum_i\eta_i=0,\qquad
 Q(\eta)=\sum_i\eta_i^2,\qquad
 m_{n,w}=\frac{\pi^2w^2}{n}.
 \label{neg:winding}
\end{equation}
Then \(S=Q+m_{n,w}\) and \(\Delta(\delta)=\Delta(\eta)\).
The winding region is exactly \(Q<\sigma-m_{n,w}\) when it is
nonempty.  In particular there are at most \(C\sqrt n\) windings.
Define
\begin{equation}
 \alpha_n=\frac{n^2-1}{2},\qquad
 P_n=\prod_{r=1}^{n-1}\frac{r!}{2^r},\qquad
 G_n=\frac{n!\pi^{(n-1)/2}P_n}{\sqrt n}.
 \label{neg:gaussian-constant}
\end{equation}
On the real zero-sum plane, with measure
\(\dd\eta_1\cdots\dd\eta_{n-1}\),
\begin{align}
 \int e^{-Q(\eta)}\Delta(\eta)^2\,\dd\eta&=G_n,
 \label{neg:gaussian-period}\\
 \int\frac{\Delta(\eta)^2}{(Q(\eta)+m)^{j+1}}\,\dd\eta
 &=G_n m^{\alpha_n-j-1}
       \frac{\Gamma(j+1-\alpha_n)}{\Gamma(j+1)}
       \quad(j+1>\alpha_n).
 \label{neg:beta-period}
\end{align}
To verify the constants, the unrestricted Gaussian Vandermonde
integral is
\(n!\prod_{r=0}^{n-1}(\sqrt\pi\,r!/2^r)
 =n!\pi^{n/2}P_n\).
This follows by writing the Vandermonde as a determinant of monic
Hermite polynomials and integrating the determinant product.
The monic squared norm is \(\sqrt\pi\,r!/2^r\), obtained by
integration by parts.  In the decomposition
\(x_i=\eta_i+u\), \(\sum\eta_i=0\), the ordinary-coordinate
Jacobian is \(n\), and the mean integral is
\(\int e^{-nu^2}\dd u=\sqrt{\pi/n}\).  This proves
\eqref{neg:gaussian-period}.  Homogeneity gives radial density
proportional to \(Q^{\alpha_n-1}\dd Q\), proving
\eqref{neg:beta-period} by the ordinary beta integral.

Let \(A^{\rm pole}_{n,j,w}\) denote the contribution of the \(w\)th
winding to \eqref{neg:pole-coefficient}.  Equations
\eqref{neg:pole-lower}--\eqref{neg:beta-period} imply
\begin{align}
 |A^{\rm pole}_{n,j,w}|
 &\le \frac{CH_{0,n}}n e^{Cn\sigma}
    \left(\frac{2n}{1-C\sigma}\right)^{j+1}
    G_n m_{n,w}^{\alpha_n-j-1}
    \frac{\Gamma(j+1-\alpha_n)}{\Gamma(j+1)}
 \quad(j+1>\alpha_n).
 \label{neg:winding-bound}
\end{align}
The dependence on \(w\) here is the explicit factor
\(|w|^{-2(j+1-\alpha_n)}\).

There is also a bound without any condition on \(j/\alpha_n\).
Use \(|T_n|\ge c/n^2\), and use homogeneity to write
\[
 \int_{Q<\sigma}\Delta(\eta)^2\,\dd\eta
       =G_n\frac{\sigma^{\alpha_n}}{\Gamma(\alpha_n+1)}.
 \]
The elementary factorial bounds in \eqref{neg:gaussian-constant}
then give
\begin{equation}
 \log\!\left(\sum_w|A^{\rm pole}_{n,j,w}|\right)
 \le 2j\log n-\frac12n^2\log n+Cn^2+Cj+C .
 \label{neg:all-index-bound}
\end{equation}
In detail,
\(\log P_n=\tfrac12n^2\log n+O(n^2)\) and
\(\log\Gamma(\alpha_n+1)=n^2\log n+O(n^2)\);
the remaining factors in \(H_{0,n}G_n\) have logarithm \(O(n^2)\).
The number of windings costs only \(O(\log(n+1))\).
This proves \eqref{neg:all-index-bound} also when \(n\) is much
larger than the Taylor order.

\subsection{An exact symmetric pole coordinate}
\label{neg:normal-form}

The next change of variables makes the pole location exactly
quadratic while preserving the collision hyperplanes of the
Vandermonde. This exactness is needed because the Taylor order grows
with the particle number. The construction also keeps its domains
and estimates independent of that dimension.

We now evaluate the first windings accurately enough for comparison
between adjacent admissible particle numbers.  All assertions in
this subsection hold for sufficiently large odd multiples of \(p\).
Put
\begin{equation}
 d_n=\frac\pi n,\qquad
 \tau_n=1-\cos d_n,\qquad B_n=B(\tau_n),\qquad
 V_n=\{v\in\C^n:\ \sum_i v_i=0\},
 \label{neg:center}
\end{equation}
and use the \(w=1\) chart \(\delta_i=d_n+v_i\).  Distinguish the
holomorphic quadratic form from the norm used in estimates:
\[
 Q(v)=\sum_i v_i^2,\qquad q(v)=\sum_i|v_i|^2 .
 \]
All transposes and gradients in holomorphic identities on \(V_n\)
use the complex bilinear extension of its real Euclidean structure.
The integration measure remains
\(\dd v_1\cdots\dd v_{n-1}\), with \(v_n=-\sum_{i<n}v_i\).

Define
\begin{equation}
 c_n=\frac{\ii B_n}{2s_\phi},\qquad
 h_n=\cos d_n-\frac{B_n b}{s_\phi^2}\sin^2d_n
       =1+O(n^{-2}).
 \label{neg:hessian}
\end{equation}
Direct differentiation, keeping the actual function \(B(t)\), gives
\begin{equation}
 g(\tau_n,d_n)=\ii\phi,\qquad
 g_t(\tau_n,d_n)=-2c_n,\qquad
 \frac12g_{\delta\delta}(\tau_n,d_n)=c_n h_n .
 \label{neg:center-derivatives}
\end{equation}
For example, \(W_t=B(t)+B'(t)(t-1+\cos\delta)\); the
\(B'\) term vanishes at this center.  Also
\(g_{\delta\delta}=W_{\delta\delta}/\sinh g
-(\cosh g)W_\delta^2/\sinh^3g\), which gives precisely
\eqref{neg:hessian}.

The pole equation in this chart is
\[
 \sum_i\bigl(g(t,d_n+v_i)-\ii\phi\bigr)=0 .
 \]
Its solution, denoted again by \(T_n(v)\), is holomorphic and
permutation invariant on a fixed ball \(\|v\|_2<r\), independent
of \(n\).  Moreover
\begin{equation}
 T_n(0)=\tau_n,\qquad
 |T_n(v)-\tau_n|\le\frac{Cq(v)}n,\qquad
 \|\nabla T_n(v)\|_2\le\frac{C\sqrt{q(v)}}n,\qquad
 \|D^2T_n(v)\|\le\frac Cn .
 \label{neg:implicit-estimates}
\end{equation}
Here and below operator norms are induced by the Euclidean norm.
To prove these estimates, subtract the constant and linear terms
of \(g(t,d_n+v_i)\) in \(v_i\).  The linear terms sum to zero,
and the remaining sum is bounded by \(Cq(v)\).  The \(t\)
derivative of the sum is bounded below by \(cn\), uniformly on
this fixed neighborhood.  Rouch\'e's theorem, or the quantitative
implicit-function theorem, gives the unique solution with
\(|T_n-\tau_n|\le Cq/n\).  Uniqueness patches the local solutions
over the whole ball.  If
\(P=I-\mathbf1\mathbf1^T/n\), its gradient is
\[
 \nabla T_n
 =-\frac{P(g_\delta(T_n,d_n+v_i))_{i=1}^n}
          {\sum_i g_t(T_n,d_n+v_i)} .
 \]
The numerator has norm at most \(C\|v\|_2\); another
differentiation gives the Hessian bound.  These arguments use
fixed scalar derivative bounds for \(g\), not a Cauchy radius
that shrinks with the dimension.

Introduce removable divided differences
\begin{align}
 \mathfrak a_n(t)
   &=\frac{g(t,d_n)-\ii\phi}{t-\tau_n},\\
 \mathfrak e_n(t,z)
   &=g(t,d_n+z)-g(t,d_n)-g_\delta(t,d_n)z,\\
 \mathfrak h_n(t,z)
   &=-\frac{2\mathfrak e_n(t,z)}
             {h_n\mathfrak a_n(t)z^2}.
 \label{neg:divided-functions}
\end{align}
The function \(\mathfrak a_n\) is uniformly bounded away from zero.
Evenness of \(g\) in its second variable implies
\begin{align}
 \mathfrak h_n(\tau_n,0)&=1,\\
 |\mathfrak h_n(t,z)-1|
   &\le C\bigl(|t-\tau_n|+|d_n||z|+|z|^2\bigr),\\
 |\partial_z\mathfrak h_n(t,z)|
   &\le C(|d_n|+|z|),\qquad
 |\partial_z^2\mathfrak h_n|+|\partial_t\mathfrak h_n|\le C .
 \label{neg:divided-bounds}
\end{align}
For clarity, Taylor's integral formula writes
\(\mathfrak e_n/z^2\) as an average of
\(g_{\delta\delta}(t,d_n+uz)\).  Its first \(z\) derivative
is an average of \(g_{\delta\delta\delta}\), an odd function
of its angular argument, and hence is \(O(|d_n|+|z|)\).
This proves the dimension-free bounds above.

Summing the divided differences at the pole gives the exact identity
\begin{equation}
 T_n(v)-\tau_n
   =\frac{h_n}{2n}
       \sum_i v_i^2\mathfrak h_n(T_n(v),v_i).
 \label{neg:separable-energy}
\end{equation}
We turn its right-hand side into a quadratic form by an explicit
permutation-equivariant change of variables.  In particular, the
construction never divides by \(Q(v)\), which has a complex null
cone.

First set
\begin{equation}
 U_i(v)=v_i\sqrt{\mathfrak h_n(T_n(v),v_i)},\qquad
 w=G(v)=P U(v),
 \label{neg:projected-map}
\end{equation}
using the square root near one.  On a sufficiently small fixed ball,
\[
 \|DG-I\|\le
 C\bigl(|d_n|\|v\|_\infty+\|v\|_\infty^2+q(v)/n\bigr)<\tfrac14 .
 \]
Thus \(G\) and its inverse are holomorphic on smaller fixed balls,
with bounded derivatives.  They commute with coordinate
permutations.  Because \(T_n(v)\) is common to all components,
one-variable divided differences give
\begin{align}
 \frac{U_i-U_k}{v_i-v_k}&=1+\mathfrak r_{ik}(v),\\
 |\mathfrak r_{ik}(v)|
 &\le C\bigl(
 |T_n-\tau_n|+|d_n|(|v_i|+|v_k|)
                   +|v_i|^2+|v_k|^2\bigr).
 \label{neg:pair-divided}
\end{align}
The quotients extend holomorphically across \(v_i=v_k\).
Projection subtracts a common mean and does not change
component differences.  Summing logarithms of the near-one
factors yields
\begin{equation}
 \left|\log\frac{\Delta(G(v))}{\Delta(v)}\right|
 \le C\bigl(n^2|T_n-\tau_n|
              +n|d_n|\sum_i|v_i|+nq(v)\bigr)
 \le C\bigl(\sqrt{nq(v)}+nq(v)\bigr).
 \label{neg:G-vandermonde}
\end{equation}
The quotient in this formula means its holomorphic extension.

The Jacobian estimate has no extra dimension loss.  The
derivative of \(G\) is a compressed diagonal matrix plus a
rank-one matrix.  The trace norm of the diagonal perturbation
is bounded by
\[
 C\bigl(n|T_n-\tau_n|+|d_n|\sum_i|v_i|+q(v)\bigr),
 \]
and that of the rank-one perturbation is at most \(Cq(v)/n\).
The convergent matrix logarithm therefore gives
\begin{equation}
 |\log\det_{V_n}DG(v)|
 \le C\bigl(|d_n|\sqrt{nq(v)}+q(v)\bigr).
 \label{neg:G-jacobian}
\end{equation}
The determinant on \(V_n\) equals the change-of-variable
determinant in the constrained coordinates, since the same
constant linear coordinate transformation is used on the
source and target.

\emph{Restoring the discarded mean.}
Projection in \eqref{neg:projected-map} enforces the zero-sum
constraint, but leaves a scalar mean in the quadratic energy.
In the \(w\) coordinates let
\[
 \mu(w)=\frac1n\sum_i U_i(G^{-1}(w)).
 \]
At fixed \(t\),
\[
 U_i=A_0(t)v_i+A_1(t)v_i^2+A_2(t,v_i)v_i^3,
 \qquad A_1(t)=O(d_n),
 \]
with uniformly bounded analytic coefficients.  The first terms
sum to zero.  Differentiating this expansion and using
\eqref{neg:implicit-estimates} gives, with \(q=q(w)\),
\begin{equation}
 |\mu|\le\frac Cn(|d_n|+\sqrt q)q,\qquad
 \|\nabla\mu\|\le\frac Cn(|d_n|+\sqrt q)\sqrt q,\qquad
 \|D^2\mu\|\le\frac Cn(|d_n|+\sqrt q).
 \label{neg:mean-estimates}
\end{equation}
For the last assertion, the second derivative of \(G^{-1}\)
is bounded by \(C(|d_n|+\|w\|_2)\), as follows directly by
differentiating \eqref{neg:projected-map}.  The chain-rule
terms therefore obey the same bound.

Define the holomorphic vector field
\begin{equation}
 \ell(w)=\int_0^1\nabla\mu(u w)\,\dd u\in V_n.
 \label{neg:ell}
\end{equation}
Then \(\mu(w)=w^T\ell(w)\), and
\[
 \|\ell(w)\|\le\frac Cn(|d_n|+\sqrt q)\sqrt q,\qquad
 \|D\ell(w)\|\le\frac Cn(|d_n|+\sqrt q).
 \]
Equivariance under the transposition of coordinates \(i,k\)
shows that
\begin{equation}
 \ell_i(w)-\ell_k(w)=(w_i-w_k)L_{ik}(w),\qquad
 |L_{ik}(w)|\le\frac Cn(|d_n|+\sqrt q).
 \label{neg:ell-pairs}
\end{equation}
Indeed compare \(w\) with the vector with these two entries
exchanged and integrate the derivative along the segment.
The segment remains in the ball.  The corresponding
one-variable divided difference also proves holomorphic
extension at the coincidence hyperplane.

Now set
\begin{equation}
 R(w)=\bigl(I+n\ell(w)\ell(w)^T\bigr)^{1/2}w
     =w+\frac{n\,\mu(w)\ell(w)}
                   {1+\sqrt{1+n\ell(w)^T\ell(w)}}.
 \label{neg:rank-one}
\end{equation}
The matrix square root is the branch near the identity.
Its perturbation has norm at most
\(C(|d_n|+\sqrt q)^2q/n\).  The map is uniformly
biholomorphic on a smaller fixed ball.  Exact multiplication
gives
\begin{equation}
 Q(R(w))=Q(w)+n\mu(w)^2=\sum_i U_i(v)^2 .
 \label{neg:rank-energy}
\end{equation}
This identity uses bilinear transposes and remains valid on
the complex null cone.

By \eqref{neg:ell-pairs}, the relative change of each pair
difference under \(R\) is at most
\(C(|d_n|+\sqrt q)^2q/n\).
The total logarithmic cost for the Vandermonde is therefore
at most \(Cnq\).  Differentiating \eqref{neg:rank-one}
shows that the trace norm of \(DR-I\) is at most
\(C(|d_n|+\sqrt q)^2q\).  For example, the term involving
\(D\ell\) has operator norm \(O((|d_n|+\sqrt q)^2q/n)\);
its trace norm costs at most the dimension \(n-1\).
All the other terms are rank one.  This proves the same
bound for the logarithmic Jacobian.

\emph{The exact quadratic coordinate.} Combining the projected
map and the rank-one correction, define
\begin{equation}
 \Psi_n(v)=\sqrt{h_n}\,R(G(v)).
 \label{neg:Psi}
\end{equation}
It is a permutation-equivariant biholomorphism on fixed
balls, \(D\Psi_n(0)=\sqrt{h_n}I\), and it preserves all
coincidence divisors.  Equations
\eqref{neg:separable-energy} and \eqref{neg:rank-energy}
prove the exact normal form
\begin{equation}
 T_n(v)=\tau_n+\frac{Q(\Psi_n(v))}{2n}.
 \label{neg:exact-quadratic}
\end{equation}
After extracting the constant \(h_n^{-\alpha_n}\), the
Jacobian--Vandermonde ratio for its inverse has a
holomorphic logarithm bounded by
\(C(\sqrt{nq}+nq)\).  The exponent follows because the
degree of \(\Delta^2\), plus the dimension of the plane,
is \(n(n-1)+(n-1)=2\alpha_n\).

\subsection{The complete residue density and its integration cycle}
\label{neg:density-cycle}

The normal form changes both the density and the integration cycle.
We first control the complete residue density relative to its center
value, then justify moving the transformed cycle back to a real
plane without crossing a pole.

In the first-winding chart put
\[
 \varrho_n(v)=
 -\frac{2H_n(T_n(v),\pi+d_n+v)}
        {\partial_t\Gamma_n(T_n(v),\pi+d_n+v)} .
 \]
Here the arguments of \(H_n\) and \(\Gamma_n\) mean
\(q_i=\pi+d_n+v_i\), so that \(\delta_i=d_n+v_i\).
The quotient
\(\varrho_n/\Delta(v)^2\) is holomorphic and nonzero on a
fixed small ball.  Its value at zero is
\begin{equation}
 \varrho_{n,0}=
 -\frac{2B_n^{\alpha_n-1/2}}
 {n\,n!(2\pi)^{n-1}2^{n(n-1)}s_\phi^{2\alpha_n}} .
 \label{neg:center-residue}
\end{equation}
To check its phase and sign, the numerator divided by
\(\Delta^2\) at the center has the factor
\((\ii s_\phi)^{-n^2}\).
Since \(n\) is odd, \(\ii^{-n^2}=-\ii\), while
\(\partial_t\Gamma_n=-\ii nB_n/s_\phi\).
The minus sign in \eqref{neg:center-residue} also includes
the minus sign from the Taylor expansion of \((t-T_n)^{-1}\).

In the coordinates \(w=\Psi_n(v)\), the complete residue form is
\begin{equation}
 \varrho_n(v)\,\dd v
  =\varrho_{n,0}h_n^{-\alpha_n}
       \Delta(w)^2(1+\mathcal E_n(w))\,\dd w,
 \qquad
 |\mathcal E_n(w)|
 \le e^{C(\sqrt{nq(w)}+nq(w))}-1 .
 \label{neg:full-density}
\end{equation}
Here the error is a relative error, after dividing by the
explicit center value.  To prove it, take logarithms only
after canceling \(\Delta^2\).  The anisotropy ratio costs
\(Cn^2|T_n-\tau_n|\).  Moreover
\[
 \sum_i|g(T_n,d_n+v_i)-\ii\phi|
 \le C\bigl(|d_n|\sum_i|v_i|+q(v)
                                  +n|T_n-\tau_n|\bigr).
 \]
Pair denominators cost at most \(Cn\) times this sum;
the one-body denominators cost only the sum itself.
The sine divided differences cost
\(C\sum_{i<k}|v_i-v_k|^2\le Cnq(v)\).
Finally the linear terms cancel in the derivative sum, giving
\[
 \frac{\partial_t\Gamma_n(T_n(v),\pi+d_n+v)}
      {n g_t(\tau_n,d_n)}
      =1+O(q(v)/n).
 \]
Together with \eqref{neg:implicit-estimates} and the
Jacobian--Vandermonde estimates, these are exactly the
logarithmic bound in \eqref{neg:full-density}.

We next justify replacing the original real cycle by a real
cycle in the exact \(w\) coordinates.  Choose a small fixed
radius \(a>0\).  For large \(n\), the real ball \(Q(v)<a^2\)
lies strictly in the first-winding region
\(Q(v)<\sigma-\pi^2/n\).  The omitted shell has
\(|T_n|\ge c/n\), by \eqref{neg:pole-lower}, and hence its
contribution to the \(j\)th coefficient is bounded by
\[
 (Cn)^j\exp\{-\kappa n^2+C n\log(n+1)+C\}.
 \]
Use \eqref{neg:gaussian} for the numerator and absorb the
volume and other factors into the indicated constants.

By choosing \(a\) small and then \(n\) large, one has
\(\|\Psi_n(v)-v\|_2\le\eta\|v\|_2\), with a fixed
\(\eta<1/4\), on a slightly larger ball.  For real \(v\),
the straight homotopy
\(w_u=(1-u)v+u\Psi_n(v)\) satisfies
\begin{equation}
 \Re Q(w_u)\ge(1-2\eta-\eta^2)Q(v)>0
                    \quad(v\ne0).
 \label{neg:cycle-positivity}
\end{equation}
This follows by writing \(w_u=v+e\),
\(\|e\|\le\eta\|v\|\), and estimating
\(2v^Te+e^Te\).  The pole
\(\tau_n+Q(w)/(2n)=0\) is therefore avoided throughout
the homotopy.

Apply Stokes' theorem to the holomorphic \((n-1)\)-form
\[
 \varrho_{n,0}h_n^{-\alpha_n}
 \frac{\Delta(w)^2(1+\mathcal E_n(w))}
      {(\tau_n+Q(w)/(2n))^{j+1}}
       \,\dd w_1\wedge\cdots\wedge\dd w_{n-1}.
 \]
The difference between its integrals over the image of the
real ball and over the real ball itself is the integral over
the lateral boundary of the homotopy.  On that boundary,
\eqref{neg:cycle-positivity} gives
\(|\tau_n+Q(w)/(2n)|\ge c/n\).
The inverse image under \(\Psi_n\) remains in a fixed
complex ball around the all-endpoint configuration.  The
complex-angle version of \eqref{neg:gaussian} bounds its
complete numerator.  The implicit derivative is at least
\(cn\); the Jacobians and the parameterized boundary volume
cost at most \(\exp(Cn\log(n+1)+C)\).
For the latter bound it is enough to use bounded operator
norms for the maps and their inverses and the elementary
Euclidean ball and sphere volume bounds in dimension \(n-1\).
Consequently the entire boundary error has the bound
\begin{equation}
 |b_{n,j}|\le
 (Cn)^j\exp\{-\kappa n^2+C n\log(n+1)+C\}.
 \label{neg:cycle-error}
\end{equation}
This argument estimates the exact residue form, not an
approximate quadratic denominator.

\subsection{A uniform beta evaluation of the first two windings}
\label{neg:beta-evaluation}

The real cycle permits an exact beta integral. The task is to show
that the relative density error remains small in the simultaneous
large-particle, high-Taylor-order limit. The two minimal windings
then have the same complex coefficient and add.

For \(j+1>\alpha_n\), \eqref{neg:beta-period} yields
\begin{equation}
 \int_{V_n\cap\R^n}
 \frac{\Delta(w)^2}
      {(\tau_n+Q(w)/(2n))^{j+1}}\,\dd w
 =
 G_n(2n)^{\alpha_n}\tau_n^{\alpha_n-j-1}
       \frac{\Gamma(j+1-\alpha_n)}{\Gamma(j+1)}.
 \label{neg:first-beta}
\end{equation}
The gamma quotient is retained exactly.  Replacing it by
\(j^{-\alpha_n}\) would introduce a logarithmic error of
order \(\alpha_n^2/j\), which is not uniformly small in
the regime required below.

We prove that \eqref{neg:full-density} has relative integrated
error
\begin{equation}
 O\!\left(\frac n{\sqrt j}+\frac{n^2}{j}\right)
 \quad\hbox{as }n\longrightarrow\infty,\quad n^2/j\longrightarrow0.
 \label{neg:relative-beta-error}
\end{equation}
Under the normalized positive measure in
\eqref{neg:first-beta}, the variable
\(x=Q/(2n\tau_n)\) has density proportional to
\[
 x^{\alpha_n-1}(1+x)^{-j-1}\,\dd x .
 \]
Since \(nQ=2n^2\tau_n x\asymp x\),
the error in \eqref{neg:full-density} is bounded by
\(e^{C(\sqrt x+x)}-1\).  On \(0\le x\le\delta\), for any
fixed small \(\delta>0\), this is at most \(C(\sqrt x+x)\).
The exact mean
\[
 \mathbb E x=\frac{\alpha_n}{j-\alpha_n}
 \]
and Cauchy--Schwarz show that its integral is bounded by
\(C(\sqrt{\alpha_n/j}+\alpha_n/j)\).

It remains to estimate the tail, where the preceding exponential
cannot be integrated over the whole half-line.  The fixed local
ball has \(x\le Ln\) for a constant \(L\).
Uniformly on \(\delta\le x\le Ln\), when \(n^2/j\to0\),
\[
 C(\sqrt x+x)\le\frac j4\log(1+x).
 \]
Write
\(Z(\alpha,j)=\Gamma(\alpha)\Gamma(j+1-\alpha)/\Gamma(j+1)\).
The tail integral divided by the original normalization is
at most
\[
 (1+\delta)^{-j/4}
       \frac{Z(\alpha_n,j/2)}{Z(\alpha_n,j)}.
 \]
Because \(\alpha_n\) is an integer,
\[
 \frac{Z(\alpha_n,j/2)}{Z(\alpha_n,j)}
 =\prod_{r=0}^{\alpha_n-1}
        \frac{j-r}{j/2-r}\le4^{\alpha_n}
                  \qquad(\alpha_n\le j/4).
 \]
Thus this tail is \(\exp(-cj+O(\alpha_n))\).
The same bound, without the exponential error factor,
controls extending the main integral from the fixed real
ball to the whole real plane.  This proves
\eqref{neg:relative-beta-error} with constants uniform in
the specified regime.

The winding \(w=-1\) is identical to \(w=1\): replace all
lifts by their negatives, use evenness of \(g\), and use the
squared sine factors.  The ordinary real measure has the same
orientation since \(n-1\) is even.  The two leading terms
therefore add.  Substituting \eqref{neg:center-residue} and
\eqref{neg:gaussian-constant} into \eqref{neg:first-beta}
gives the following exact-factor model:
\begin{align}
 A^{\rm pole}_{n,j,1}+A^{\rm pole}_{n,j,-1}
 &=M(n,j)(1+e_{n,j})+b_{n,j},
 \label{neg:local-model}\\
 |e_{n,j}|&\le C\left(\frac n{\sqrt j}+\frac{n^2}{j}\right),
 \label{neg:local-error}\\
 M(n,j)&=
 -\frac{2^{2-\alpha_n}B_n^{\alpha_n-1/2}P_n
                         n^{\alpha_n-3/2}}
        {\pi^{(n-1)/2}s_\phi^{2\alpha_n}}\,
 h_n^{-\alpha_n}\tau_n^{\alpha_n-j-1}
       \frac{\Gamma(j+1-\alpha_n)}{\Gamma(j+1)} .
 \label{neg:model}
\end{align}
The term \(b_{n,j}\) satisfies \eqref{neg:cycle-error}.
The power \(B_n^{\alpha_n-1/2}\), including its varying
complex phase, is part of the model.

\subsection{Exact finite differences and a single dominant index}
\label{neg:dominance}

We compare the moduli of the exact-factor models at neighboring
admissible particle numbers. Strict discrete concavity isolates a
single maximum, and the Taylor order is chosen so that the selected
index sits at that maximum. This yields noncancellation without any
assumption on the phases of the other indices.

Let
\[
 L(n,j)=\log|M(n,j)|,\qquad d=2p .
 \]
The admissible indices are the odd multiples of \(p\), so their
spacing is \(d\).  Elementary Stirling bounds in
\eqref{neg:model} give, as \(n\to\infty\) and \(n^2/j\to0\),
\begin{equation}
 L(n,j)=2j\log n-\frac{n^2}{2}\log j+c_*n^2
         -j\log(\pi^2/2)
         +O(n\log n+n^4/j+j/n^2).
 \label{neg:L-asymptotic}
\end{equation}
For example,
\[
 \log P_n=\frac12n^2\log n
             -\left(\frac34+\frac12\log2\right)n^2
             +O(n\log n),
 \qquad
 \log\tau_n=\log(\pi^2/2)-2\log n+O(n^{-2}),
 \]
while the logarithm of the gamma quotient is
\(-\alpha_n\log j+O(\alpha_n^2/j)\).
These formulas give exactly the value of \(c_*\) in
\eqref{neg:cstar}.

The adjacent-index comparison requires finite differences of
\(L\). We compute them from the exact factors, since the error
bound in \eqref{neg:L-asymptotic} gives no derivative control.
The resulting identity is
\begin{equation}
 L(n+d,j)-2L(n,j)+L(n-d,j)
 =d^2\left(-\frac{2j}{n^2}-\log j+2c_*\right)
    +O_d(n^2/j+j/n^4+n^{-2}).
 \label{neg:second-difference}
\end{equation}
Here is a direct verification with the required uniform
remainder.  First,
\begin{align}
 \log P_{n+d}-2\log P_n+\log P_{n-d}
 &=\sum_{k=0}^{d-1}\sum_{\ell=1}^{d}
        \log(n-d+k+\ell)-d^2\log2\\
 &=d^2(\log n-\log2)+O_d(n^{-2}).
 \label{neg:P-difference}
\end{align}
The offsets in this double sum have sum zero.
Next, because \(\alpha_n\) is an integer, for real
\(j>\alpha_{n+d}\) the exact identity is
\[
 \log\frac{\Gamma(j+1-\alpha_n)}{\Gamma(j+1)}
 =-\alpha_n\log j+
      \sum_{r=0}^{\alpha_n-1}-\log(1-r/j).
 \]
The upper and lower adjacent increments of \(\alpha_n\)
have lengths \(dn+d^2/2\) and \(dn-d^2/2\).
Pair the \(dn-d^2/2\) corresponding summands; their
arguments differ by \(O_d(n)\), and the derivative of
\(-\log(1-x/j)\) is at most
\(1/(j-\alpha_{n+d})\).
Their total difference is \(O_d(n^2/j)\).
The \(d^2\) unpaired summands cost the same amount.
Thus this factor contributes
\(-d^2\log j+O_d(n^2/j)\).

For a precise accounting of all remaining factors, put
\[
 c_\tau=\pi^2/2,\qquad
 r_n=\log\tau_n-\log c_\tau+2\log n,\qquad
 K=-\tfrac12\log2+\tfrac12\log|B(0)|
                      -\log s_\phi+\tfrac12\log c_\tau .
 \]
Up to a constant independent of \(n,j\), the logarithm
\(L(n,j)\) is exactly
\begin{align*}
 &2j\log n-j\log c_\tau-\alpha_n\log j
   +\log P_n-\tfrac12n^2\log n+Kn^2
                 -\tfrac n2\log\pi+\log n\\
 &\quad+
 \sum_{r=0}^{\alpha_n-1}-\log(1-r/j)
 +(\alpha_n-\tfrac12)\Re\log(B_n/B(0))
 -\alpha_n\Re\log h_n+(\alpha_n-j-1)r_n .
\end{align*}
The last three terms have convergent expansions in \(n^{-2}\),
and
\[
 r_n=-\frac{\pi^2}{12n^2}+O(n^{-4}).
 \]
Their second differences are
\(O_d(n^{-2}+j/n^4)\), as is seen either term by term in
these convergent expansions or by differentiating them on
a fixed neighborhood of \(n^{-1}=0\).
Also
\[
 \Delta_d^2(-\tfrac12n^2\log n)
     =-d^2(\log n+\tfrac32)+O_d(n^{-2}),\qquad
 \Delta_d^2(2j\log n)
     =-\frac{2d^2j}{n^2}+O_d(j/n^4).
 \]
Combine these with \eqref{neg:P-difference} and
\(2c_*=2K-\tfrac32-\log2\).  This proves
\eqref{neg:second-difference} directly from the exact
model.

In particular, throughout the intermediate range
\begin{equation}
 \log j<n<\frac{\sqrt j}{(\log j)^{1/4}},
 \label{neg:middle-range}
\end{equation}
the second difference is at most \(-c d^2\log j\) for all
sufficiently large \(j\), whenever the neighboring indices
are also in that range.  Indeed \(n^2/j\to0\) uniformly,
and \(j/n^4\) is smaller than \(j/n^2\) by a factor at most
\((\log j)^{-2}\).  This is concavity over the entire
intermediate range.

Fix a large admissible \(n\), and temporarily let \(j\) be
real of order \(n^2\log n\).  Set
\[
 f_n(j)=L(n+d,j)-L(n-d,j).
 \]
The two gamma arguments differ by the integer \(2dn\).
Differentiating their exact product gives
\begin{align}
 f_n'(j)
 &=\log\frac{\tau_{n-d}}{\tau_{n+d}}
   -\sum_{r=1}^{2dn}
       \frac1{j+1-\alpha_{n-d}-r}\\
 &=\frac{4d}{n}-\frac{2dn}{j}
       +O_d(n^{-3}+n^3/j^2)>0 .
 \label{neg:crossing-derivative}
\end{align}
Exact first differences also give
\begin{equation}
 \frac{f_n(j)}{2d}
 =\frac{2j}{n}-n\log j+2c_*n
       +O_d(\log n+n^3/j+j/n^3).
 \label{neg:first-difference}
\end{equation}
For detail, the quotient \(P_{n+d}/P_{n-d}\) is a
product of \(2d\) factorials; Stirling's bound on these
individual factorials gives its centered first difference
\(n\log n-(1+\log2)n+O_d(\log n)\).
The centered first difference of the gamma correction sum
is \(O_d(n^3/j)\).  The convergent \(n^{-2}\) expansions
above contribute \(O_d(n^{-3}+j/n^3)\).
The displayed leading terms then follow exactly as in the
second-difference calculation.

At
\[
 j=n^2\left(\log n+\tfrac12\log\log n-c_*\pm\epsilon\right)
 \]
the expression in \eqref{neg:first-difference} has opposite
signs for each fixed \(\epsilon>0\) and all large \(n\).
There is a unique crossing in this interval by
\eqref{neg:crossing-derivative}.  Round it to the nearest
integer \(j_n\).  We obtain
\begin{equation}
 j_n=n^2\left(\log n+\tfrac12\log\log n-c_*+o(1)\right),
 \qquad f_n(j_n)=O_d(1/n).
 \label{neg:selected-order}
\end{equation}
These \(n\) lie strictly inside \eqref{neg:middle-range}.

Let \(\lambda_k=L(n+kd,j_n)\) for the admissible indices
in that range.  Their successive first differences decrease
by at least \(c d^2\log j_n\).  Equation
\eqref{neg:selected-order} gives
\(\lambda_1-\lambda_{-1}=O_d(1/n)\).
Consequently
\[
 \lambda_0-\lambda_{\pm1}\ge c'd^2\log j_n,
 \qquad
 \lambda_k-\lambda_0
       \le-c''d^2 k^2\log j_n+O_d(|k|/n)\quad(k\ne0).
 \]
Summing this discrete Gaussian proves
\begin{equation}
 \sum_{\substack{m\ne n,\ m\ {\rm odd},\ p\mid m\\
                         m\ {\rm satisfies}\ \eqref{neg:middle-range}}}
       |M(m,j_n)|
 \le |M(n,j_n)|\,O(j_n^{-c'''d^2}).
 \label{neg:single-index}
\end{equation}
The complex phases of the different models play no role in
this inequality.

\subsection{Exhausting the remaining indices and windings}
\label{neg:exhaustion}

The comparison must cover the complete coefficient, including
indices far from the selected one. We divide the remainder into
small indices, analytic and cycle errors, the infinite large-index
tail, higher windings, and the other indices in the intermediate
range. Each class is compared with the same selected model.

We compare every part of \eqref{neg:jet-sum} with the selected
model.  First, \eqref{neg:L-asymptotic} and
\eqref{neg:selected-order} imply
\begin{align}
 \log|M(n,j_n)|
 &=j_n\log j_n-j_n\log\log j_n
          +j_n\bigl(\log(4/\pi^2)-1\bigr)+o(j_n),
 \label{neg:model-growth}\\
 |M(n,j_n)|^{1/j_n}
 &\sim\frac{4j_n}{e\pi^2\log j_n}.
 \label{neg:model-root}
\end{align}
Indeed \eqref{neg:selected-order} gives
\[
 n^2=\frac{j_n}{\tfrac12\log j_n-c_*+o(1)}.
 \]
Substituting this in the four leading terms of
\eqref{neg:L-asymptotic} gives \eqref{neg:model-growth};
each error there is \(o(j_n)\).

The following estimates cover all the complementary terms.

\paragraph{Small particle indices.}
For \(m\le\log j\), the full coefficient bound
\eqref{neg:small-disk} gives
\begin{equation}
 \sum_{\substack{m\le\log j\\m\ {\rm odd}}}
       |[t^j]C_m(t)|
 \le \exp\{2j\log\log(j+2)+Cj+C\}.
 \label{neg:small-indices}
\end{equation}
This is negligible compared with \eqref{neg:model-growth}
when \(j=j_n\).  Equivalently, after the pole reduction one
can apply the same bound to the poles plus their analytic
remainders, since their difference has the smaller bound
\eqref{neg:regular-remainder}.

\paragraph{All analytic and cycle remainders.}
The sum of the regular remainders in
\eqref{neg:regular-remainder} and all the terms
\eqref{neg:cycle-error} is bounded by
\begin{equation}
 \exp\{\tfrac12j\log(j+2)+Cj+C\}.
 \label{neg:all-remainders}
\end{equation}
This too is negligible at the selected orders.  This estimate
includes the entire nonresonant particle class \(p\nmid m\)
outside the small-index range.

\paragraph{All sufficiently large particle indices.}
For
\[
 m\ge\frac{\sqrt j}{(\log j)^{1/4}},
 \]
use \eqref{neg:all-index-bound}, which has no restriction
on \(m^2/j\).  With \(x=m^2\), its exponent is at most
\[
 j\log x-\frac14x\log x+Cx+Cj+C .
 \]
Its derivative in \(x\) is negative on
\(x\ge j/\sqrt{\log j}\), for all large \(j\), and at the
lower endpoint its value is
\begin{equation}
 j\log j-\frac12j\log\log j
            -\frac14j\sqrt{\log j}
            +o(j\sqrt{\log j}).
 \label{neg:large-indices}
\end{equation}
For example its derivative is at most \(-c\log j\)
there for large \(j\), since
\(j/x\le\sqrt{\log j}\) and
\(\tfrac14\log x\sim\tfrac14\log j\).
Summing the resulting decreasing exponential over the
discrete values \(x=m^2\) costs at most a harmless polynomial
factor.  Thus \eqref{neg:large-indices} bounds the logarithm
of the whole infinite large-index tail up to
\(o(j\sqrt{\log j})\).
Its gap below \eqref{neg:model-growth} tends to infinity
faster than \(j\).

\paragraph{Higher odd windings in the intermediate range.}
Let \(M_0(m,j)>0\) be the modulus of the expression
\eqref{neg:model} with \(B_m\) replaced by \(B(0)\),
\(h_m\) by \(1\), and \(\tau_m\) by
\(\pi^2/(2m^2)\).  The exact factors give
\begin{equation}
 \left|\log\frac{M_0(m,j)}{|M(m,j)|}\right|
 \le C(1+j/m^2).
 \label{neg:M0}
\end{equation}
For the \(\tau\) factor this follows from
\(\log(\tau_m/(\pi^2/(2m^2)))=O(m^{-2})\);
the \(B\) and \(h\) factors cost \(O(1)\).
Comparison of \eqref{neg:winding-bound} with this
explicit \(M_0\) yields, for each odd \(|w|\ge3\),
\[
 |A^{\rm pole}_{m,j,w}|
 \le C M_0(m,j)e^{Cm\sigma}
       (1-C\sigma)^{-j-1}
       |w|^{-2(j+1-\alpha_m)} .
 \]
The constants \(C\) here do not depend on \(m,j,w\).
Choose the fixed \(\sigma\) still smaller, if necessary,
so that
\(-\log(1-C\sigma)<\tfrac12\log3\).
Throughout \eqref{neg:middle-range},
\[
 \alpha_m=o(j),\qquad m=o(j),\qquad j/m^2=o(j)
 \]
uniformly.  The number of possible windings is at most
\(C\sqrt m\).  Hence for some \(c>0\),
\begin{equation}
 \sum_{\substack{w\ {\rm odd}\\|w|\ge3}}
       |A^{\rm pole}_{m,j,w}|
 \le |M(m,j)|e^{-cj}
 \label{neg:higher-windings}
\end{equation}
uniformly in that range.  The choice of a smaller
\(\sigma\) merely changes the fixed constants in the
earlier arguments.

\paragraph{The first windings of every intermediate index.}
Equation \eqref{neg:local-error} tends uniformly to zero
in \eqref{neg:middle-range}; it is at most
\(C(\log j)^{-1/4}\).
At \(j=j_n\), \eqref{neg:single-index} shows that the sum
of the moduli of all intermediate models other than
\(M(n,j_n)\) is \(o(|M(n,j_n)|)\).
The sum of all their relative errors is bounded by
\[
 o(1)\sum_m|M(m,j_n)|=o(|M(n,j_n)|).
 \]
The same is true of all higher windings by
\eqref{neg:higher-windings}.

These classes exhaust every odd particle number and every
winding, together with all analytic and cycle remainders.
Combining them in \eqref{neg:jet-sum} proves the actual
full-sum coefficient asymptotic
\begin{equation}
 a_{j_n}=M(n,j_n)(1+o(1)).
 \label{neg:full-asymptotic}
\end{equation}
In particular, this is a statement about the sum of the
complex coefficients, not merely a comparison of individual
singularity sets.

\subsection{Nonextension and uniformity}
\label{neg:conclusion}

Equations \eqref{neg:model-root} and
\eqref{neg:full-asymptotic} prove
\eqref{neg:growth-statement}, so
\(\limsup_j|a_j|^{1/j}=\infty\).
If \(F\) admitted a holomorphic extension to a disk about
zero from the physical side \(\mathcal U^-\), its derivatives
on the inward negative ray would agree with those of that
extension.  By \eqref{neg:jets}, its Taylor coefficients at
zero would therefore be the \(a_j\).  Cauchy's estimate on
any smaller extension disk would bound
\(|a_j|\le C R^{-j}\), a contradiction.
If \(P(0)\ne0\), a holomorphic extension of \(PF\) would,
after division by the nonvanishing holomorphic factor
\(P\), give an extension of \(F\).  Thus a local physical
normalization by such a factor cannot remove the obstruction.

Finally, all the constants used above arise from a common
holomorphic neighborhood for \(B\), a lower bound for
\(-\Im B(0)\), and fixed scalar bounds for the dispersion
and its derivatives.  A holomorphic family with
\(\Im B(0,\eta)\le-\beta<0\) on a relatively compact
parameter polydisk supplies these uniformly by Cauchy's
estimate.  The square roots and logarithms near one used
in the normal form also have common neighborhoods.
The exact coefficient model continues to use the actual
values \(B(\tau_n,\eta)\); nothing in the argument freezes
the anisotropy.  The crossing argument and all range
 estimates are uniform on each such compact family, with
 \(c_*=c_*(\eta)\) and the selected orders \(j_n(\eta)\).
This completes the proof of Theorem~\ref{neg:theorem}.


\section{The isotropic restriction: a unified even--odd proof}
\label{iso:section}

The isotropic line is collinear in the joint edge geometry. We
therefore prove its local obstructions directly from the complete
form-factor integrals. The analytic estimates treat the two phases
together, with particle parity as the phase-dependent ingredient.
First we identify the full-site coefficients, then evaluate the first
singular order, and finally control every higher order uniformly
near the boundary.

Write $\zeta$ for the auxiliary exterior variable, and put
\[
 E(\zeta)=\zeta+\zeta^{-1},\qquad
 D(x,y;\zeta)=E(\zeta)-\tfrac12(x+x^{-1}+y+y^{-1}).
\]
For the low-temperature restriction $\zeta=s$; for the high-temperature
restriction $\zeta=s^{-1}$. This is a change of variable, not a duality
identification between the two susceptibilities. Define
\begin{align}
 T_N^\sigma(\zeta)
 &=\frac1{N!}\oint\cdots\oint
 \frac{Z^{-1}+Y^{-1}}{(1-\sigma Z)(1-\sigma Y)}
 \prod_{i<j}P_{ij}\prod_i R(y_i;\zeta)\frac{\dd y_i}{2\pi\ii},
 \label{iso:reduced}\\
 f_N(\zeta)
 &=\frac1{N!}\oint\cdots\oint
 (ZY)^{-1}\prod_{i<j}P_{ij}\prod_iR(y_i;\zeta)\frac{\dd y_i}{2\pi\ii},
 \label{iso:onsite}
\end{align}
where $\sigma\in\{1,-1\}$, $Y=\prod_i y_i$, $Z=\prod_i z_i$, and
\[
 W=E(\zeta)-\tfrac12(y+y^{-1}),\quad
 z=W-\sqrt{W^2-1},\quad R=\frac{2z^2}{1-z^2},\quad
 P_{ij}=-\frac{(y_i-y_j)^2z_i z_j}{y_i y_j(1-z_i z_j)^2}.
\]
On real $\zeta>1$ the root satisfies $0<z<1$ on the angular circle;
elsewhere it is the continued interior root. Put $T_N=T_N^1$ and
$C_N=f_N+2T_N$.

These expressions retain the full-site normalization. Their double-contour
lift, with $X=\prod_i x_i$, is
\begin{equation}
 T_N^\sigma=\frac1{N!}\oint\cdots\oint
 \frac{X^{-1}+Y^{-1}}{(1-\sigma X)(1-\sigma Y)}
 \prod_{i<j}\frac{x_i-x_j}{1-x_ix_j}
             \frac{y_i-y_j}{1-y_iy_j}
 \prod_i\frac{\dd x_i\,\dd y_i}{(2\pi\ii)^2D(x_i,y_i;\zeta)}.
 \label{iso:double}
\end{equation}
For $f_N$ replace the first rational factor by $(XY)^{-1}$. Choose a
common radius $r<1$ with $|z(y;\zeta)|<r$ on the $y$ circle. The only
enclosed dispersion pole is then $x=z$, with residue
$2z^2/(1-z^2)$. At $x=0$, $D^{-1}=-2x+O(x^2)$ cancels the
apparent numerator pole. All pair and product poles remain outside after
each successive residue. Subtracting two dispersion equations gives
\[
 (z_i-z_j)(1-z_i z_j)
 =-\frac{z_i z_j}{y_i y_j}(y_i-y_j)(1-y_i y_j),
\]
and hence the canceled expression for $P_{ij}$. The full lattice
multiplier is exactly
\begin{equation}
 \frac{(1+X)(1+Y)}{XY(1-X)(1-Y)}
 =\frac1{XY}+\frac{2(X^{-1}+Y^{-1})}{(1-X)(1-Y)}.
 \label{iso:sites}
\end{equation}
Thus $f_N$ is the onsite coefficient, $T_N$ is the indicated half of the
offsite coefficient, and $C_N$ is the full lattice coefficient. Equivalently,
\begin{equation}
 C_N(\zeta)=\frac1{N!}\int_{\sum q_i=0\ (2\pi)}
 \prod_i\frac1{\sinh\gamma_i}\,
 \frac{1+Z}{1-Z}
 \prod_{i<j}\left[
 \frac{\sin((q_i-q_j)/2)}{\sinh((\gamma_i+\gamma_j)/2)}
 \right]^2\prod_{i<N}\frac{\dd q_i}{2\pi},
 \label{iso:angular}
\end{equation}
where $\cosh\gamma_i=E(\zeta)-\cos q_i$ and $z_i=e^{-\gamma_i}$.
The independent measure has $N-1$ coordinates, with
$q_N=-\sum_{i<N}q_i$ modulo $2\pi$.
For completeness, this equality can be checked before any boundary limit.
On real $\zeta>1$ let $q_N(m,n)$ be $1/N!$ times the $N$-angle integral with
multiplier $Z^mY^n$, measure
$\prod_i\dd q_i/(2\pi\sinh\gamma_i)$, and the displayed squared pairs.
Angle reversal gives $q_N(m,n)=q_N(m,-n)$; exchanging the two coordinate
sets in \eqref{iso:double} gives $q_N(m,n)=q_N(n,m)$ for $m,n\ge0$.
Any permutation sign of the top form is canceled by the same permutation
of the product orientation, for either parity. A common small angle
strip has $\Re\gamma_i\ge c>0$, so simultaneous angle shifts give, at
fixed $N,\zeta$, a bound $C_{N,\zeta}e^{-N(cm+\eta|n|)}$.
Consequently the spatial sum is absolutely convergent. Summing $n\in\Z$
enforces $\sum q_i=0$ modulo $2\pi$, and summing the reflected $m$
coordinate gives $(1+Z)/(1-Z)$. Alternatively,
\[
 \frac{Y+Z}{(1-Y)(1-Z)}
 =\sum_{m\ge1,n\ge0}Y^mZ^n+
  \sum_{m\ge0,n\ge1}Y^mZ^n
\]
counts each positive axis once and each open positive quadrant twice.
Thus $f_N+2T_N$ counts exactly every lattice site, establishing
\eqref{iso:angular} without omitting the onsite term.

\subsection{Compact convergence and a dense arithmetic family}

Compact convergence defines the holomorphic sum inside its domain.
The arithmetic construction then selects boundary points where one
particle order is the first possible resonance and its two collision
charts are uniquely determined.

For even $N$ the Schur identity states that
\[
 \Pf\left[\frac{v_i-v_j}{1-v_iv_j}\right]
 =\prod_{i<j}\frac{v_i-v_j}{1-v_iv_j}.
\]
For odd $N$, append $v_{N+1}=-1$; its last upper column is the removable
constant one, and the same identity gives the product over the $N$ real
indices. One proof puts $a_i=(1+v_i)/(1-v_i)$ and compares residues at
$a_N=-a_i$ inductively. Both rational functions have the same residues,
are bounded at infinity, and vanish at $a_N=a_{N-1}$.

On a compact exterior neighborhood choose fixed radii $r<1$ with all
interior roots bounded by $q<r$ and bounded away from zero. The global
inverse products cost only $C^N$. Hadamard's inequality gives
$|\Pf A|\le C^{m/2}(m-1)^{m/4}$ for an even matrix of size $m$.
Applying this to both matrices, including the ghost index if necessary,
gives
\begin{equation}
 |T_N^\sigma|+|f_N|+|C_N|
 \le C_0\frac{C_1^N N^{(N+1)/2}}{N!}.
 \label{iso:normal}
\end{equation}
The bound is summable. Cauchy's estimate on a larger compact neighborhood
proves local normal convergence of every fixed derivative. The constants
in this compact statement are not asserted uniform at the unit circle.

\begin{lemma}[The selected first orders]\label{iso:arithmetic}
Let $p\ge11$ be prime, choose distinct
$a,b\in\{1,\ldots,(p-1)/2\}$, and set
\[
 \alpha=2\pi a/p,\quad\beta=2\pi b/p,\quad
 E_*=\cos\alpha+\cos\beta=2\cos\theta,
 \quad\zeta_*=e^{\ii\theta}.
\]
For the even family choose $\alpha,\beta\in(0,\pi/2)$ and put
$N_0=2p$, $k=N_0^2/2-1$. For the odd family choose both angles in
$(0,\pi/2)$ or both in $(\pi/2,\pi)$ and put
$N_0=p$, $k=(p^2-1)/2$.
The first order of the respective parity satisfying a Nickel equation is
$N_0$, and its unordered cosine pair is uniquely
$\{\cos\alpha,\cos\beta\}$. No endpoint configuration is resonant at any
order. The even family is dense in the upper-right arc; the odd families
are dense in both upper quadrants.
\end{lemma}

\begin{proof}
Let $\omega=e^{2\pi\ii/p}$, $d=(p-1)/2$,
$u_j=\omega^j+\omega^{-j}$, and
$F=\mathbb Q(\omega+\omega^{-1})$. Every nonzero-index $u_j$ has
trace $-1$. Thus $v=u_a+u_b$ has trace $-2$ and is irrational:
if rational it would be an integer, whereas its trace would force
$v=-2/d$. If $p\nmid n$, then
$\mathbb Q(\omega)\cap\mathbb Q(e^{2\pi\ii/n})=\mathbb Q$.
Indeed the compositum has degree
$\varphi(pn)=\varphi(p)\varphi(n)$, so the intersection has degree one.
A lower order of the required parity would contradict this fact.

The only rational relation between $1,u_1,\ldots,u_d$ is a multiple of
$1+\sum_j u_j=0$: lift the relation to a polynomial of degree at most
$p-1$ and use the minimal polynomial $1+X+\cdots+X^{p-1}$.
At order $p$ all cosines are $u_j/2$ or $1$; at order $2p$ they are
$\pm u_j/2$ or $\pm1$. An equality of two cosine pairs has at most four
nonconstant positions. Since $d\ge5$, a zero coefficient forces the
multiple of the relation to vanish, proving uniqueness.

For $E_*>0$, an endpoint resonance requires a root of unity $\xi$ with
$\xi+\xi^{-1}=v-2$. Its trace is $-2-2d<-2d$, impossible because
every conjugate of a root-of-unity trace belongs to $[-2,2]$.
For $E_*<0$ the possible endpoint equation is
$\xi+\xi^{-1}=v+2$. Direct expansion gives
$\operatorname{Tr}_F(v^2)=4d-6$. If all conjugates of $v+2$ were in
$[-2,2]$, every conjugate of $v$ would be in $[-4,0]$, where
$v^2\le-4v$. Taking traces would imply $4d-6\le8$, contrary to
$d\ge5$. Approximating any angle in either relevant half interval by
two distinct adjacent prime-grid angles proves density.
\end{proof}

\begin{lemma}[Exponential endpoint separation]\label{iso:separation}
If a fixed algebraic unit-modulus number $\xi$ is not a root of unity,
there is $A>0$ such that
$|1-\xi^N|,|1+\xi^N|\ge e^{-AN}$ for every positive integer $N$.
\end{lemma}
\begin{proof}
Write its primitive irreducible polynomial as
$a_d\prod_{\nu=1}^d(X-\xi_\nu)$ with $\xi_1=\xi$.
Its nonzero integer resultant with $X^N-1$ has modulus at least one.
Bounding every other conjugate by
$|\xi_\nu^N-1|\le2\max(1,|\xi_\nu|)^N$ proves the first assertion.
The inequality $|1+\xi^N|\ge |1-\xi^{2N}|/2$ proves the second after
increasing $A$.
\end{proof}

\subsection{The first complete singular coefficient}

At a selected point, the local calculation has two parts: bounded
contributions from the complementary configurations, and explicit
residues from the two collision charts. Their complex phases must be
matched before the leading coefficient can be used for the full sum.

Put $\zeta_\eps=(1+\eps)\zeta_*$ and keep the integration radius
$r=e^{-c_0\eps}$ fixed during every complex derivative. On a sufficiently
small disk of radius $c\eps$, $\Im W\ge c'\eps$ for every angle.
The interior roots belong to the lower unit half-disk and satisfy
$|z|<r$ if $c_0$ is small enough. These statements hold for either
upper quadrant.

\begin{lemma}[Fixed-order complements]\label{iso:fixed}
At each point in Lemma~\ref{iso:arithmetic}, every fixed derivative of
each lower coefficient of the selected parity is bounded on the exterior
radial approach. At order $N_0$, the complement of the two lower
double-contour configurations
\[
 (x_i,y_i)=(e^{-\ii\beta},e^{-\ii\alpha})\quad\hbox{for all }i,
 \qquad\hbox{and their exchange}
\]
has bounded derivatives of every fixed order. The onsite coefficient has
bounded fixed derivatives for every fixed $N$.
\end{lemma}

\begin{proof}
The potentially singular factors on the unit double torus have gradient
vectors consisting of the two positive all-ones product vectors, the
positive reciprocal-pair vectors, and the dispersion vectors with entries
$(\Im x_i,\Im y_i)$. The sign of dispersion damping is fixed by
$\Im E(\zeta_\eps)>0$. A nonnegative zero combination cannot contain
a pair vector: its two positive entries would require both corresponding
imaginary parts to be negative, whereas reciprocal circle coordinates
have opposite imaginary parts.

If both product vectors occur, every coordinate is lower and the ratios
$\sin A/\sin B$ agree along $\cos A+\cos B=E_*$. Along that curve
\[
 \frac{\dd}{\dd A}\frac{\sin A}{\sin B}
 =\frac{E_*(1-\cos A\cos B)}{\sin^3B}
\]
has a strict constant sign. All pairs therefore coincide, and both common
coordinates are $N$th roots of unity. Lemma~\ref{iso:arithmetic} leaves
exactly the specified configurations at $N_0$ and none below it.
With only one product vector, the other coordinates must all be $+1$
when $E_*>0$, or all $-1$ when $E_*<0$; this is the excluded endpoint
resonance. With neither product vector, a nonzero dispersion vector cannot
cancel because $E_*\notin\{-2,0,2\}$.

On the complement, a finite smooth partition supplies a real direction
strictly separating all active vectors from zero on each chart. Write
its $\ell$ singular factors as damped positive-half-line exponential
integrals, and put $R=\sum_{i=1}^\ell\xi_i$. The exponent $g$ has
$\Re g\le0$, $|e\cdot\nabla g|\ge cR$, and fixed angular jets
$O(R)$. A fixed $j$th parameter derivative costs $O(R^j)$.
Integrating by parts $\ell+j+1$ times with
$(e\cdot\nabla)/(e\cdot\nabla g)$ leaves radial majorant $CR^{-2}$
for $R\ge1$, including all coefficient and cutoff derivatives. The
small-$R$ part is bounded directly. Omitting both product factors proves
the onsite assertion. This argument is at fixed $N,j$ only.
\end{proof}

\begin{proposition}[Nonzero first amplitude]\label{iso:first}
There is $L_*\ne0$ such that
\[
 T_{N_0}^{(k)}(\zeta_\eps)=
 \begin{cases}
 2L_*\eps^{-1/2}+o(\eps^{-1/2}),&N_0\text{ even},\\
 2L_*\eps^{-1}+o(\eps^{-1}),&N_0\text{ odd}.
 \end{cases}
\]
For even $N_0$ all lower derivatives are bounded. For odd $N_0$ the
$(k-1)$st derivative is $O(1+|\log\eps|)$ and all smaller derivatives
are bounded. The corresponding leading coefficient for $C_{N_0}$ is
$4L_*$.
\end{proposition}

\begin{proof}
Use fixed smooth weights depending only on the $y$ angles around the two
specified configurations, with disjoint supports. All $x$ integrations
remain unweighted, so the residues giving \eqref{iso:reduced} apply to
each chart exactly. Lemma~\ref{iso:fixed} controls the full complementary
double integral.

In the first chart write
$y_i=r e^{\ii(-\alpha+u_i)}$, $v=\sum_i u_i$,
$u_i=v/N+\eta_i$, $\sum_i\eta_i=0$.
The absolute real Jacobian in coordinates
$(v,\eta_1,\ldots,\eta_{N-1})$ is one. We use positive real product
measures here; equivalently, the induced domain orientation compensates
the determinant $(-1)^{N-1}$ of this coordinate change.
Write $z_i=e^{-\ii\phi_i}$. At the center,
\[
 b=\phi_u=\frac{\sin\alpha}{\sin\beta}>0,\quad
 d=-\tfrac12\phi_{uu}
 =\frac{H_{\alpha\beta}}{2\sin^3\beta}\ne0,\qquad
 H_{\alpha\beta}=E_*(1-\cos\alpha\cos\beta).
\]
Also
\[
 \phi=\beta-\ii A_0\eps+bu-du^2+
 O(\eps^2+\eps|u|+|u|^3),\quad
 A_0=\frac{2\sin\theta-c_0\sin\alpha}{\sin\beta}>c_0.
\]
Replace the smooth mean cutoff by a fixed hard segment; its difference
has separated product phases if the shape radius is sufficiently small.
Move that segment downward in a fixed rectangle. The $Y=1$ pole is at
$v=-\ii Nc_0\eps$; downward motion increases the damping of $Z$.
The vertical sides have a fixed phase gap and the bottom fixed damping,
so all their fixed derivatives are bounded. The clockwise closure and
$\partial_v(1-Y)=-\ii$ give the actual residue multiplier $+2\pi$.
No nonholomorphic weight is moved through a complex contour.

At the residue the $y_i$ lie on the unit circle and
\begin{equation}
 1-Z=Q\eps-\ii d\sum_i\eta_i^2+
 O(\eps^2+\eps|\eta|+|\eta|^3),\quad
 Q=\frac{2N\sin\theta}{\sin\beta}>0,
 \label{iso:first-denominator}
\end{equation}
with nonlinear bound $|1-Z|\ge c(\eps+\sum_i\eta_i^2)$.
To see this bound, modulus damping is at least $c\eps$; permutation
symmetry cancels the linear shape term, and the real phase Hessian has a
strict sign determined by $-H_{\alpha\beta}$. Its integral remainder
controls the phase outside $\sum\eta_i^2=O(\eps^2)$, while damping
controls the remaining region.

The exact numerator is $\Delta(\eta)^2 A(\zeta,\eta)$ with $A$ analytic
and nonzero at the center. Fixed parameter derivatives preserve the
Vandermonde. Each pair contributes
$-(\eta_i-\eta_j)^2/(4\sin^2\beta)$, and each one-body angle measure
contributes $y_0z_0/(2\pi\ii\sin\beta)$. Since $y_0^Nz_0^N=1$,
their combined phase is
$(-1)^{N(N-1)/2}\ii^{-N}=\ii^{-N^2}$.
After the actual mean residue the leading constrained value is
\[
 \frac{2\ii^{-N^2}}{N!(2\pi)^{N-1}2^{N(N-1)}\sin^{N^2}\beta}
 \frac{\Delta(\eta)^2\,\dd\eta_1\cdots\dd\eta_{N-1}}
 {Q\eps-\ii d\sum_i\eta_i^2}.
\]
After $k$ derivatives and $\eta=\sqrt\eps\,w$, only the top pole
power contributes at scale $\eps^{-1/2}$ for even $N$ and
$\eps^{-1}$ for odd $N$. The radial majorants have tails $R^{-2}$
and $R^{-3}$ respectively. Lower pole powers vanish after multiplication
by the indicated singularity scale; splitting at $R=\sqrt\eps$ gives
the logarithm at the odd lower derivative. This is dominated convergence
after the mean residue, not of the unintegrated absolute mean--shape
integral.

For $a=(N^2-1)/2$, let
\[
 c_N=\frac1{\sqrt N}\int_{\sum\omega_i=0,\ |\omega|=1}
                         \Delta(\omega)^2\,\dd\omega>0.
\]
The remaining radial period is
\[
 \int\frac{\Delta(w)^2\,\dd w_1\cdots\dd w_{N-1}}
 {(Q-\ii d\sum_iw_i^2)^{k+1}}
 =\frac{c_N}{2}Q^{a-k-1}(-\ii d)^{-a}
 B(a,k+1-a),
\]
where the quadratic coefficient is continued from positive real part.
Here $k+1-a$ equals $1/2$ or $1$, and every factor is nonzero.
For odd $N=p$, $a=k\equiv0\pmod4$, so $d^k=|d|^k$ also on the
negative-$E_*$ family. The pole derivative is $Qe^{-\ii\theta}$.
In particular the complete one-chart coefficient is
\begin{equation}
 L_*=
 \frac{\ii^{-N^2}(-1)^k k!\,c_N B(a,k+1-a)\,
       2^a(2N\sin\theta)^{a-1}e^{-\ii k\theta}}
 {N!(2\pi)^{N-1}2^{N(N-1)}(-\ii H_{\alpha\beta})^a}.
 \label{iso:amplitude}
\end{equation}
Indeed the sine exponent is $-N^2-(a-1)+3a=0$.
For even $N$, $H_{\alpha\beta}>0$ and the last power uses argument
$-\pi/2$. For odd $N$, $a=k$ is divisible by four, so that power is
$|H_{\alpha\beta}|^k>0$ for either sign of $E_*$.
The factor $\ii^{-N^2}$ equals one for even $N$ and $-\ii$ for odd
$N$. Every remaining factor is symmetric in $\alpha,\beta$.
This checks the equality of the two complex coefficients, including
orientation and phase, rather than just equality of their moduli.
The two charts therefore add. Finally
$C_N=f_N+2T_N$ and Lemma~\ref{iso:fixed} give the factor four.
\end{proof}

\subsection{A uniform boundary-tail theorem}

The first singular coefficient is useful only if all higher particle
orders are smaller at the same derivative order. The theorem below
provides that comparison for the complete differentiated tail; its
proof keeps the onsite terms and both signed kernels.

Negating $(\zeta,y,z)$ gives, directly from the forms and their measures,
\begin{equation}
 T_N^\sigma(-\zeta)=T_N^{\sigma(-1)^N}(\zeta),\qquad
 f_N(-\zeta)=(-1)^Nf_N(\zeta).
 \label{iso:sign}
\end{equation}
Conjugation follows from the real germ. Consequently it suffices to
estimate either signed kernel at a selected point with $0<E_*<2$.
Set $b_*=\arccos(E_*-1)$ and $\xi=e^{-\ii b_*}$.
Lemma~\ref{iso:arithmetic}, also applied after the reflection of the
negative family, and Lemma~\ref{iso:separation} give the needed
$|1\pm\xi^N|\ge e^{-AN}$.

\begin{theorem}[The entire differentiated tail]\label{iso:tail}
For either parity, its selected first order $N_0$, and every fixed
$0\le j\le k$,
\[
 \sum_{\substack{N>N_0\\N\text{ of the selected parity}}}
 |T_N^{(j)}(\zeta_\eps)|=o(\eps^{-1/2}),\qquad
 \sum_{\substack{N>N_0\\N\text{ of the selected parity}}}
 |C_N^{(j)}(\zeta_\eps)|=o(\eps^{-1/2}).
\]
The estimates include all particle orders and both upper quadrants
needed by the odd family.
\end{theorem}

We divide the proof into contour deformation, pair suppression,
large-order estimates, and a separate estimate for the intermediate
orders. The latter uses differentiation on the real integration
domain and coarea bounds. The parameter disks are specified before
the corresponding Cauchy estimates are applied.
Fix the point and $k$ first. Put
\[
 H=\log(1/\eps),\qquad 0<\alpha_0<\frac1{2(k+2)},\qquad
 \lambda_*=\eps^{\alpha_0}.
\]
All contour functions below are fixed smooth periodic functions, with
bounded jets of every order required through $k$.

Choose $m$ supported near the lower branch angle $-b_*$ and equal to one
on a neighborhood of every lower branch-cutoff support. Choose $a$
supported in $[\delta,\pi-\delta]$, equal to one on
$[2\delta,\pi-2\delta]$, and choose $p$ supported in $(0,\pi)$ and
equal to one on a neighborhood of $\operatorname{supp}a$.
The upper inverse branch is inside $a=p=1$, and the supports of $m,p$
are disjoint. For a small fixed $\tau>0$ define
\[
 P=\sum_i p(\theta_i),\quad
 h_i=-2\tau p(\theta_i)+\frac{\tau P}{2N}m(\theta_i),\quad
 y_i(\lambda,\theta)=r e^{\ii\theta_i+\lambda h_i},\quad
 \psi=\prod_i(1-a(\theta_i)).
\]
Write $U_{\lambda,\zeta}$ for the pullback of the complete normalized
form \eqref{iso:reduced}, including $1/N!$.

\begin{lemma}[Exact weighted deformation]\label{iso:deformation}
For every $N$ and either $\sigma$,
\begin{align*}
 T_N^\sigma&=F_N+K_N+S_N,\\
 F_N&=\int(1-\psi)U_{1,\zeta},\qquad
 K_N=\int\psi U_{0,\zeta},\\
 S_N&=-2\ii\tau\int_0^1\sum_q
             \int(\partial_{\theta_q}\psi)U_{\lambda,\zeta}\,\dd\lambda.
\end{align*}
Every current term has one named upper index $q$ with
$p_q=1$, $p'_q=m_q=m'_q=0$; all other indices are in
$\operatorname{supp}(1-a)$.
\end{lemma}

\begin{proof}
The total radial shift satisfies $\sum h_i\le-3\tau P/2$, so $|Y|<1$.
Moreover,
\[
 \partial_{\log|y|}\Im W=-\cosh(\log|y|)\sin\theta.
\]
Both inward upper motion and outward lower motion increase $\Im W$.
All roots remain interior; the product and canceled-pair denominators
remain nonzero throughout the closed homotopy.

After dividing row $i$ of the angular Jacobian by $y_i$, its matrix is
\[
 M_{ij}=d_i\delta_{ij}+\frac{\lambda\tau}{2N}m_i p'_j,
 \quad d_i=\ii+\lambda\left(-2\tau p'_i+
                                      \frac{\tau P}{2N}m'_i\right).
\]
Since $m_i p'_i=0$, its determinant is $\prod_i d_i$.
At a named current index its row and column are diagonal with $d_q=\ii$;
the homotopy replacement column multiplies the determinant by $2\ii\tau$.
Cartan's formula on the closed torus gives
$\partial_\lambda\int\psi U_\lambda
=-2\ii\tau\sum_q\int(\partial_q\psi)U_\lambda$.
Combine this with invariance of the unweighted closed integral.
\end{proof}

\paragraph{Support geometry and pair suppression.}
Here and below ``branch'' means the lower branch, never its excluded
upper inverse. Choose nested branch cutoffs: an outer one supported in
$|u|\le\delta_B$, equal to one on $|u|\le\delta_B/2$, and an inner
one supported in $|u|\le2r_0$, equal to one on $|u|\le r_0$, where
$u=\theta+b_*$ and $r_0=\delta_B/(2L^2)$ with $L$ sufficiently large.
All these supports lie in $m=1$. Telescope the product of the outer
cutoffs. A complementary term has an outer compact index with
$|u_q|\ge\delta_B/2$, including all cutoff jets. Split the other
coordinates by the inner cutoff and compact left/right labels. The latter
can be made periodic by returning their transition inside the upper
$a=1$ plateau, where the full $K$ and current weights vanish with all
jets. Thus no cutoff differentiation introduces a new support type.

The partition has at most $C^N$ labels. Its types are: all branch; a
compact left root with a modulus gap; a separated compact/true-branch
pair; and all compact right. A named current has no all-branch sector.
If it has a branch index, choose the named upper index and a fixed branch
index as the separated pair. The nested widths ensure that the branch
slope dominates the compact slope by at least four; preliminary upper
endpoint neighborhoods fix the latter bound before their final shrinking.

For clarity the quadratic pair estimate used throughout is proved next.
At $\eps=\lambda=0$, lower angles lie in $[-\pi,0]$; their roots are
real in $[z_0,1]$ on $[-\pi,-b_*]$ and are $e^{-\ii v}$,
$0\le v\le b_*$, on $[-b_*,0]$, where
$z_0=E_*+1-\sqrt{(E_*+1)^2-1}>0$.
The identity
\begin{equation}
 |1-vw|^2-|v-w|^2
 =(1-|v|^2)(1-|w|^2)+4\Im v\Im w
 \label{iso:schur-gap}
\end{equation}
and its angular version imply $|P_{ij}|\le1$.
Equality is possible only for $y_i=1$ paired with a left angle, or
$y_i=-1$ paired with a right angle, with either ordering.
Choose fixed overlapping groups
\[
 B=[-b_*-w,-b_*+w],\quad
 L=[-\pi,-b_*-\ell_L],\quad R=[-b_*+\ell_R,0],
 \qquad0<\ell_L,\ell_R<w.
\]
There is no equality configuration in a same-group square.
For $B^2$ the angular Schur factor is at most
$\sin w/\min(\sin(b_*-w),\sin(b_*+w))<1$ after shrinking $w$;
the other groups have a strict root Schur bound. At the double branch
the apparent singularity is removable: the inverse root coordinate is
$y_\zeta(\phi)=G_\zeta(\phi^2)$ and
\begin{equation}
 P_{ij}=(\phi_i-\phi_j)^2 A_\zeta(\phi_i,\phi_j),
 \qquad A_{\zeta_*}(0,0)=-\frac1{4\sin^2b_*},
 \label{iso:regular-pair}
\end{equation}
with analytic bounded fixed jets. This follows by factoring both
$y_i-y_j$ and $1-e^{-\ii(\phi_i+\phi_j)}$; the common
$\phi_i+\phi_j$ factor cancels.

On the closed lower support, $1-z_i z_j=0$ forces the already resolved
double root $z_i=z_j=1$: the alternative $z_i=z_j=-1$ is excluded by
$W_0+1\ge E_*>0$. Away from a smaller double-branch chart, this
denominator therefore has a fixed positive minimum. The regular
expression \eqref{iso:regular-pair} and the original rational expression
agree on their overlaps, and give a finite cover with bounded fixed
jets. This also proves continuity at the common endpoints, where the
canceled pair is zero. No reciprocal-root pole remains on that cover.

Choose a strict same-group bound $q<q_1<q'<1$ and cross-group slack
$\eta>0$ satisfying $\log(1+\eta)<-\log(q')/4$.
Project the upper slivers $[0,2\delta]$ and $[\pi-2\delta,\pi]$ to their
respective endpoints, assigning them to $R,L$. If $\bar\theta_i$ denotes
this projection (and equals $\theta_i$ for lower angles), then uniformly
in the entire independent box $0\le\lambda,P/N\le1$,
\[
 |y_i-e^{\ii\bar\theta_i}|+
 |W_i-W_0(\bar\theta_i)|
 \le C(\eps+\tau+\delta)+C|\zeta-\zeta_\eps|.
\]
Compact-root changes are Lipschitz, and branch changes are one-half
H\"older by the square-root factorization and
\eqref{iso:regular-pair}. The finite cover therefore gives a complete-pair
change at most
\[
 C\sqrt{\eps+\tau+\delta+|\zeta-\zeta_\eps|}.
\]
Choose $\delta,\tau,\eps_0$ to make the radial-center change smaller
than $\min(q_1-q,\eta/2)$; reserve $q'-q_1$ and $\eta/2$ for the
parameter disks. This proves both same-group and cross-group bounds
on the slivers, uniformly in the actual coupled occupancy.

The named current index needs only bounded pairs. At its right limiting
endpoint the root is $\xi=e^{-\ii b_*}$, and for every limiting lower root
$z$, $|1-\xi z|=|e^{\ii b_*}-z|\ge\sin b_*>0$.
At the left endpoint the root is $z_0<1$ and
$|1-z_0z|\ge1-z_0$. These margins persist on the chosen neighborhoods.
The order of choices is essential: the branch/upper-sliver limiting
moduli are
\[
 \frac{\sin((b_*+a)/2)}{\sin((b_*-a)/2)}
 =1+a\cot(b_*/2)+O(a^2),\qquad
 \frac{\cos((b_*-a)/2)}{\cos((b_*+a)/2)}
 =1+a\tan(b_*/2)+O(a^2).
\]
They can exceed one. Keeping the compact-group margins fixed before
shrinking the slivers is what preserves the strict same-group bound;
no claim that every sliver pair contracts is used.

For $n=N$ lower variables in $K$, or $n=N-1$ in a current, the number
of same-group pairs is at least $n^2/6-n/2$ and the number of cross pairs
at most $n^2/3$. After deleting at most $j$ factors, apply the bounds
directly to the remaining product; no division by a possibly zero pair
is allowed. The result is
\begin{equation}
 \prod_{(i,l)\notin\mathcal E}|P_{il}|
 \le C_j^N e^{-\kappa N^2},\qquad |\mathcal E|\le j,
 \label{iso:quadratic}
\end{equation}
for some fixed $\kappa>0$. Named pairs cost only $C^N$.

\paragraph{The selected contour, including arbitrarily large $N$.}
Telescope $1-\psi$ so that a selected index has $p_q=1$ and hence
$P\ge1$. At $\lambda=1$ it has a fixed interior root gap, and
$\sum h_i\le-3\tau P/2$ supplies a fixed $Y$ gap.
Every true branch has $\Im W\ge cP/N$, so its residue factor is
bounded by $C\sqrt N$. These facts persist on a common disk
$|\zeta-\zeta_\eps|\le c/N$. Branch roots remain in the lower unit
half-disk; compact roots move by $O(1/N)$, and the fixed anchor absorbs
their total inflation $e^{CN(c/N)}$ when $c$ is small.

Let $M$ be the number of compact indices. Divide their radial-center roots
into a strict interior group and a compact lower-arc group. Same-group
Schur factors are at most a fixed number below one. For a branch root
$w$ on the actual disk and a compact radial-center root $z_0$, the
denominator $|1-z_0w|$ has a fixed positive lower bound, and
\[
 \frac{z-w}{1-zw}-\frac{z_0-w}{1-z_0w}
 =\frac{(z-z_0)(1-w^2)}{(1-zw)(1-z_0w)}.
\]
Thus mixed factors are at most $1+C/N$; only the compact root has been
displaced in this comparison. It follows, even with the weaker
$1+C/\sqrt N$ budget, that the entire root Pfaffian product is at most
$C^N e^{-aM^2}$.

Every $y$-pair involving a true branch is bounded: its reciprocal upper
partner has inward shift $-2\tau$, dominating the outward shift at most
$\tau/2$. Compact pairs have the occupancy-independent majorant
\[
 \left|\frac{y_i-y_j}{1-y_iy_j}\right|
 \le\frac C{\eps+\dist(\theta_i+\theta_j,2\pi\Z)}.
\]
Indeed the sum of the two shifts is nonpositive at reciprocity and is
$O(\dist(\theta_i+\theta_j,2\pi\Z))$ nearby. A matched two-angle
integral therefore costs $O(H+1)$.
Expand the angular Pfaffian into bounded minors and a matching on an
even compact subset. Hadamard bounds the minor by
$C^N N^{(N+1-2r)/4}$, the subset and matching count by $N^{3r}$,
and the disjoint singular pair integrals by $(C(H+1))^r$.
Including the one-body cost $C^NN^{N/2}$ and $N!$, the complete
selected integral is bounded, for each assignment, by
\[
 C^N N^{-N/4+1/4}(CN^3(H+1))^M e^{-aM^2}.
\]
The $C^N$ assignments and named-index choices are absorbed in $C^N$.
Completing the square in $M$, and applying Cauchy only to the complete
integral with its fixed angular weights, gives
\begin{equation}
 \sum_N|F_N^{(j)}(\zeta_\eps)|
 \le C_j\exp(C_j\log^2(H+2)).
 \label{iso:selected-bound}
\end{equation}
Here $e^{C_j\log^2(N+1)}N^{-N/4+j+1/4}$ is summable.
Pointwise assignments need not be holomorphic, because none is separately
differentiated. The argument applies to every $N$, not just a moving
finite window, and also with the product poles deleted.

\paragraph{The original and protected current disks.}
On $|\zeta-\zeta_\eps|\le c\eps$, the undeformed imaginary margin
and its monotone radial improvement keep all roots lower and interior.
The two signed product gaps are at least $c'\eps$.
Near a true branch set $\mu=\lambda P/N$ and $u=\theta+b_*$.
For $y=e^{v+\ii\theta}$, $v=-c_0\eps+\tau\mu/2$, the exact formulas
\[
 \Re W=\Re E-\cosh v\cos\theta,\qquad
 \Im W=\Im E-\sinh v\sin\theta
\]
give an imaginary margin $c(\eps+\mu)$ and real branch-center
displacement $O(\eps^2+\mu^2+|\zeta-\zeta_\eps|)$.
Splitting at $|u|=K(\eps+\mu)$, first using the imaginary margin and
then the linear real term, proves the occupancy-independent majorant
\begin{equation}
 |R(y;\zeta)|\le C(|u|+\eps)^{-1/2}.
 \label{iso:body}
\end{equation}
Its integral is bounded. Compact factors and Jacobians cost $C^N$.
Branch-coordinate motion on the disk is $O(\sqrt\eps)$; choosing
$\eps_0$ small retains the reserved pair slack, so
\eqref{iso:quadratic} holds throughout the disk. Cauchy therefore gives
\begin{equation}
 |K_N^{(j)}|+|J_{N,j}(\lambda)|
 \le C_j^N\eps^{-j-2}e^{-\kappa N^2},\qquad0\le\lambda\le1,
 \label{iso:original-bound}
\end{equation}
where $J_{N,j}$ is the already differentiated, angularly integrated
current. This bound will be used only for large $N$.

For a named current and $\lambda\ge\lambda_*$, $P\ge1$ protects a
larger disk $|\zeta-\zeta_\eps|\le c\lambda_*/N^2$.
Its branch margin is at least $c(\eps+\lambda P/N)$.
Choose its disk constant after $\tau$ so that the parameter displacement
uses at most one quarter of this margin. Compact roots move by
$O(\lambda_*/N^2)$; the named upper root has modulus at most
$e^{-c\lambda}$ and absorbs the total compact inflation.
Both product gaps are therefore at least $c'\lambda$.
The same split proving \eqref{iso:body}, now including the disk
displacement, gives the same integrable majorant. Pair-coordinate motion
uses the remaining fixed slack in \eqref{iso:quadratic}.
Cauchy and then integration in $\lambda$ yield
\begin{equation}
 \sum_N\int_{\lambda_*}^1|J_{N,j}(\lambda)|\,\dd\lambda
 \le C_j\lambda_*^{-j-2}
 =C_j\eps^{-\alpha_0(j+2)}.
 \label{iso:large-current}
\end{equation}
The convergent sum $\sum_NC^NN^{2j}e^{-\kappa N^2}$ includes all
orders. No large-current estimate is used to justify the original disk.

\subsection{Real-domain differentiation in the intermediate window}
\label{iso:lie-section}

Direct differentiation raises the powers of the product poles.
Instead, a tangential correction to the parameter derivative keeps
those kernels simple. The resulting derivatives fall on the
numerator and cutoffs; the Vandermonde zeros compensate for the
resulting collision losses. Vanishing puncture fluxes justify every
integration by parts before the uniform estimates are taken.

The remaining orders satisfy $N_0+2\le N<D\sqrt H$, with fixed $D$.
For the small current $0\le\lambda\le\lambda_*$, all finite compact
jet errors $N^{C_j}(\eps+\lambda_*)$ tend uniformly to zero in this
window. We prove the required bound
\begin{equation}
 |K_N^{(j)}|+\int_0^{\lambda_*}|J_{N,j}(\lambda)|\,\dd\lambda
 \le e^{C_jN\log(N+1)}(H+N)^{C_j}.
 \label{iso:middle-bound}
\end{equation}

\begin{lemma}[Differentiation with simple kernels]\label{iso:lie}
Let $w$ be a fixed smooth real angular weight, and let $V$ be a
parameter-dependent complex vector field. Put
$\mathcal D=\partial_\zeta-\mathcal L_V$ on top forms. If all
puncture fluxes vanish at every preceding stage, then
\[
 \partial_\zeta^j\int w\Omega_\zeta
 =\int\mathcal D^j(w\Omega_\zeta).
\]
If $(\partial_\zeta-V\cdot\nabla)Y=
(\partial_\zeta-V\cdot\nabla)Z=0$, both signed product kernels
remain simple at every order.
\end{lemma}
\begin{proof}
The Lie derivative of a top form is a divergence. Integrate it on the
fixed real domain, with punctures removed, and use the vanishing flux.
Repeat on the entire preceding top form, including all derivatives of
$V$, its divergence, and the previous coefficients. Complex coefficients
are permitted by linearity of real integration by parts; no flow of a
smooth cutoff into a complex domain is assumed.
\end{proof}

On an original branch chart set $\phi=\arccos W$ and
$\chi_\zeta:u\mapsto\phi$. The branch one-body form is regular:
\[
 \frac{\dd y}{2\pi\ii}\frac{z}{\sinh\gamma}
 =-\frac{z}{\pi(1-y^{-2})}\,\dd\phi.
\]
Together with \eqref{iso:regular-pair}, the regular numerator is
$\Delta(\phi)^2A_\zeta(\phi)$ with analytic bounded fixed jets.
Put $a_i=E'/W_{u_i}$, $b_i=\partial_{u_i}\phi_i$, $A=\sum_i a_i$.
For a selected pair $p,q$ use
\[
 V_i=a_i\ (i\ne p,q),\qquad
 V_p=a_p+\frac{Ab_q}{b_p-b_q},\qquad
 V_q=a_q-\frac{Ab_p}{b_p-b_q}.
\]
Then $\sum V_i=0$ and $\sum b_iV_i=\partial_\zeta\sum\phi_i$.
The residual regular-coordinate field
$B=\partial_\zeta\chi_\zeta-(\chi_\zeta)_*V$ has only
\[
 B_p=-\frac{Ab_pb_q}{b_p-b_q},\qquad B_q=-B_p.
\]
Because $b(\phi)=B_\zeta(\phi^2)/\phi$ with $B_\zeta$ analytic
and $B_\zeta(0)\ne0$, one has the exact analytic factorization
\[
 \phi_qB_\zeta(\phi_p^2)-\phi_pB_\zeta(\phi_q^2)
 =(\phi_q-\phi_p)
       \{B_\zeta(0)+O(|\phi_p\phi_q|)\}.
\]
The expression in braces is bounded away from zero on a smaller fixed
chart. Thus $b_pb_q/(b_p-b_q)$ is an analytic bounded coefficient
divided by $\phi_q-\phi_p$, with bounded required jets.
All fixed coefficient jets have polynomial
$N$-cost apart from this explicit selected difference pole.
For smooth real $w$ the exact pullback identity is
\begin{equation}
 (\partial_\zeta-\mathcal L_V)(w\chi_\zeta^*\Omega)
 =w\chi_\zeta^*(\partial_\zeta+\mathcal L_B)\Omega
 -(V\cdot\nabla_u w)\chi_\zeta^*\Omega.
 \label{iso:pullback}
\end{equation}
It follows by differentiating the Jacobian determinant and subtracting
the two divergence formulas. The real cutoff is never holomorphically
extended.

After $j$ operations, each term has a derivative
$\partial_u^\mu w$, a numerator derivative
$\partial_\zeta^\ell\partial_\phi^\nu(\Delta^2A)$, the unchanged
simple kernels, and a coefficient of selected pole order $m$ satisfying
\begin{equation}
 m+|\mu|+|\nu|\le2j,\qquad
 \ell+|\mu|+|\nu|\le j.
 \label{iso:filtration}
\end{equation}
The filtration follows by checking each product-rule term. A
fixed-$\phi$ parameter derivative of a coefficient adds no difference pole; $B\cdot\nabla$ on a coefficient,
or multiplication by $\operatorname{div}B$, adds at most two; a
numerator-coordinate derivative adds one derivative and at most one
pole; a numerator-parameter derivative adds only one derivative; and a
cutoff derivative adds one cutoff derivative and at most one pole.
These are all product-rule terms in \eqref{iso:pullback}.
There are $N^{O(j)}$ descendants. At most $j$ complete pair factors are
changed in each descendant, so \eqref{iso:quadratic} survives on the
others. One must not divide by the original Vandermonde, which would
create poles at unselected collisions.

To select pairs smoothly, fix an integer $M>2k+2$ and use
\[
 w_{pq}=\frac{(u_p-u_q)^{2M}}{\sum_{i<l}(u_i-u_l)^{2M}},\quad
 d_u=\max_{i<l}|u_i-u_l|.
\]
Off full equality these sum to one and obey
\[
 |\partial_u^\mu w_{pq}|
 \le C_kN^{C_k}\frac{|u_p-u_q|^{2M-|\mu|}}{d_u^{2M}}.
\]
The analytic inverse branch is \(u=U_\zeta(\phi^2)\), with its
fixed contour constants included. Its bounded derivative gives
\[
 |u_p-u_q|\le C|\phi_p^2-\phi_q^2|
                  \le C'|\phi_p-\phi_q|.
\]
The rational weights therefore absorb the
selected poles, losing at most $2j$ powers of shape radius. A further
power is reserved for a transition or flux estimate.
At fixed $\eps,N$, a selected-diagonal tube of radius $h$ has flux
$O(h^{2M+1-2j})$ on a separated support. At full equality, the shape
surface measure gives flux exponent
\[
 N(N-1)-2j-1+N-2=N^2-2j-3>0
 \quad(N\ge N_0+2,\ j\le k).
\]
Thus all puncture fluxes vanish, in this order, at every preceding Lie
stage. These checks justify the identities before taking boundary-uniform
bounds.

\paragraph{Branch geometry used in coarea.}
For the original contour, near $u=0$,
\[
 1-W=a_*u-\ii b_*'\eps+O(u^2+\eps|u|+\eps^2),
 \quad a_*=\sin b_*>0,\quad b_*'=2\sin\theta-c_0a_*>0.
\]
The exact factorization $\phi=\sqrt{2(1-W)}\,h(1-W)$, $h(0)=1$,
gives
\[
 |\phi'|\asymp(|u|+\eps)^{-1/2},\quad
 \int|\phi'|\,\dd u\le C,\quad
 -\Im\phi\ge c\sqrt{|u|+\eps}\quad(u\le0).
\]
For $u\ge0$, $x'=\partial_u\Re\phi\asymp(u+\eps)^{-1/2}$,
and $-x''\ge cu^{-3/2}$ for $u\ge C_0\eps$.
If $M_1/\eps\to\infty$ and $0\le m_1\le M_1/2$, then
\[
 x'(m_1)-x'(M_1)\ge c(m_1+\eps)^{-1/2},\quad
 \frac{|\phi'(m_1)\phi'(M_1)|}{x'(m_1)-x'(M_1)}\le CM_1^{-1/2}.
\]
For $M_1/2<m_1<M_1$, with $m_1\gg\eps$, the last ratio is at most
$C\sqrt{M_1}/(M_1-m_1)$. These follow from the leading square-root
derivative, using its $u^{-3/2}$ curvature in the comparable case.
On a true branch plateau of the small current the same factorization
holds with scale $\eps+\lambda P/N$. In particular
$\Re\phi_u\ge(2/3)|\phi_u|$ for $u\ge0$, and
\begin{equation}
 \frac{|\phi'_u|}{|1-\sigma Z|}
 \le\frac C{|u|+\eps+\lambda P/N}\quad(u<0).
 \label{iso:negative-branch}
\end{equation}
The possible real branch-center displacement is quadratic in this scale;
it is retained in deriving these bounds. A plateau variation fixes $P$.

\paragraph{All-branch microcore.}
Choose $b_N=c_*e^{-(A+4)N}/N$ and localize by
$\prod_i\chi(u_i/b_N)$, with $\chi=1$ on $[-1/2,1/2]$ and
supported in $[-1,1]$. Uniformly for $N<D\sqrt H$,
$N\eps e^{AN}\to0$ and $\eps/b_N^m\to0$ for each fixed $m$.
Therefore $|1-\sigma Y|\ge\tfrac12e^{-AN}$ on the core and all
cutoff jets. The one-phase field $V_i=a_i$ freezes every regular
$\phi_i$, hence $Z$, and preserves the Vandermonde polynomial.
Its $j$th derivative has only $e^{O_j(N)}$ cutoff and $Y$-kernel losses
and $e^{O_j(N^2)}$ analytic product cost.
The branch image has radius $C\sqrt{b_N}$ in a fixed cone; choosing
$c_*$ small makes $N\sqrt{b_N}$ small and
$|1-\sigma Z|\ge c\sum_i|\phi_i|$ (a weaker bound also valid for
$\sigma=-1$). The Laplace identity for the reciprocal of this sum and
\[
 \int_{\rm core}e^{-t|\phi|}|\dd\phi|
 \le C\min(\sqrt{b_N},t^{-1})
\]
bound the core derivative by
\begin{equation}
 e^{C_jN^2}b_N^{(N^2-1)/2}.
 \label{iso:microcore}
\end{equation}
Its logarithm is $-(A+4)N^3/2+O_j(N^2\log(N+1))$, and the series
is bounded independently of $\eps$.

\paragraph{All-branch complement.}
Telescope the complement of the microcore. A named anchor satisfies
$|u_a|\ge b_N/2$ on every cutoff jet. Put
$\bar u=N^{-1}\sum u_i$, $\rho^2=\sum(u_i-\bar u)^2$ and split
smoothly at $\rho_N=e^{-BN}$, where $B$ is fixed large enough that
$2\rho_N\le b_N/8$ for every $N\ge N_0+2$.

On $\rho\ge\rho_N$, maximal separation is at least
$\rho_N/\sqrt N$. The Lie losses and cutoff jets cost
$e^{O_j(N\log(N+1))}$ and unchanged pairs retain quadratic
suppression. If all coordinates are nonnegative, select the actual
minimum $m_1$ and maximum $M_1$ only in this final absolute integral.
For fixed other variables the phase map
\[
 (m_1,M_1)\longmapsto
 (m_1+M_1+\text{constant},\ x(m_1)+x(M_1)+\text{constant})
\]
has positive Jacobian $x'(m_1)-x'(M_1)$ and multiplicity one on each
extreme-index cell. The preceding branch estimates absorb the two
branch measures at cost $e^{O(N)}N^{O(1)}$; here $M_1\ge b_N/2$,
and in the comparable case $M_1-m_1\ge\rho_N/\sqrt N$.
The two unwrapped phase intervals have bounded length for fixed other
variables.

If some negative coordinate has magnitude at least
$\delta_N=e^{-C_1N}\le b_N/16$, then $Z$ has gap $e^{-O(N)}$;
integrate the remaining $Y$ kernel in the named anchor, whose measure
costs at most $e^{O(N)}$. If every negative coordinate has magnitude
below $\delta_N$, the anchor is positive; use one negative coordinate
with \eqref{iso:negative-branch} and use the anchor for the $Y$ kernel.
Both integrations cost at most $e^{O(N)}(H+N)$.

On $\rho\le2\rho_N$, the anchor gives
$|\bar u|\ge3b_N/8$, so all coordinates have the same sign and
$|u_i|\ge b_N/4$. The branch measures and the conversion of the
Vandermonde to the angular shape radius cost at most
\[
 C^{N^2}b_N^{-N^2/2}
 =\exp\left(\tfrac{A+4}{2}N^3+\tfrac12N^2\log N+O(N^2)\right).
\]
The remaining shape power after Lie losses and coarea is
\[
 L_N=N^2-2j-4\ge
 \begin{cases}4N_0+2,&N_0\text{ even},\\4N_0+1,&N_0\text{ odd}.
 \end{cases}
\]
For negative mean the $Z$ gap is $e^{-O(N)}$ and a single mean phase
suffices. For positive mean the phase map
$(\bar u,\rho)\mapsto(N\bar u+\text{constant},\sum_i x(u_i))$
has Jacobian at least $cN\rho$, by strict concavity and the zero first
shape derivative. It has multiplicity one on each shape ray.
Bound $\rho^{L_N}$ by $(2\rho_N)^{L_N}$ before integrating the
phase kernels. Both unwrapped phase lengths are $O(N)$, giving an
explicit factor $CN^2(H+N)^2$. The logarithm of the total bound is
\[
 \left(\tfrac{A+4}{2}-B\right)N^3+\tfrac12N^2\log N
 +O_j(N^2)+B(2j+4)N,
\]
apart from that phase factor. Choose $B>A+4$ after fixing the
anchor-neighborhood constants. This is summable; finitely many small
orders have fixed, possibly large, constants. The same positive power
already proved the required full-equality puncture fluxes vanish.

\emph{How the phase integrations are estimated.}
For the simple-kernel integrations just used, the elementary inequality
\[
 |1-e^{-a+\ii t}|^{-1}
 \le C\bigl(\min(1,a)+\dist(t,2\pi\Z)\bigr)^{-1}
\]
gives $C(1+\ell)(1+|\log a|)$ on an interval of length $\ell$.
For two phases, divide out the specified real Jacobian, enlarge to the
containing rectangle, and apply this estimate twice. Periodic wrapping
is accounted for by interval lengths, separately from the stated
unwrapped multiplicity. A coarea radius is never integrated a second
time after being exchanged for a phase.

\paragraph{Mixed and compact sectors.}
Use regular $\phi_i$ on true branch indices $J$, and original angles
on compact indices. Let $q$ be the named compact anchor, choose
$j_0\in J$, and put
\[
 A=\sum_{i\in J}a_i,\quad D_0=\sum_{i\notin J}\partial_\zeta\phi_i,
 \quad X_0=\frac{D_0+b_qA}{b_{j_0}-b_q}.
\]
Set $V_i=a_i$ for $i\in J\setminus\{j_0\}$,
$V_{j_0}=a_{j_0}+X_0$, $V_q=-A-X_0$, and all others zero.
This freezes $Y,Z$ and every unselected branch phase. The sole residual
branch velocity is
$-(D_0+b_qA)/(1-b_q/b_{j_0})$.
The nested support has $|b_q/b_{j_0}|\le1/4$, so this coefficient and
all fixed jets have polynomial $N$-cost in regular coordinates; there
is no selected-difference pole. On current supports the chosen indices
are on true plateaus, giving $V\cdot\nabla P=V\cdot\nabla\sum m_i=0$
identically, including parameter derivatives. The hybrid Jacobian is
triangular; its branch measure is precisely the regular one displayed
above, and its compact factor retains the full $\prod_i d_i$ determinant.

For positive branch coordinate, the phase Jacobian in the compact
anchor and selected branch is
$\Re(b_{j_0}-b_q)\ge(2/3-1/4)|b_{j_0}|$; its fixed-first-phase map is
injective and absorbs the selected branch measure. If a branch
coordinate is negative, use \eqref{iso:negative-branch} instead.
With a compact left root, $Z$ has a fixed gap: take $V_i=a_i$ on
branches and $V_q=-A$, freezing $Y$ and each branch root, and
differentiate the separated $Z$ kernel at polynomial cost.
These alternatives prove \eqref{iso:middle-bound} for the mixed and
left sectors using \eqref{iso:quadratic} and the simple phase estimates.

On an all compact right interval put
$g(\theta)=\arccos(E_*-\cos\theta)$. Direct differentiation gives
\[
 g''(\theta)=-\frac{E_*(1-\cos\theta\cos g)}{\sin^3g}\le-c_R<0.
\]
The coupled-contour finite jets differ from their limiting values by
$N^{O(j)}(\eps+\lambda)$, uniformly small in this window. Thus
$G=\Re\sum_i\phi_i$ retains a strictly negative Hessian.
Let $f_i=\partial_{\theta_i}\log Y$,
$g_i=\partial_{\theta_i}\sum\phi_i$, and
$Q_0=\partial_\zeta\sum\phi_i$. By permutation symmetry and smooth
division,
\[
 f_pg_q-f_qg_p=(\theta_p-\theta_q)c_{pq},\qquad|c_{pq}|\ge c_R/2.
\]
The field $V_p=-Q_0f_q/(f_pg_q-f_qg_p)$,
$V_q=Q_0f_p/(f_pg_q-f_qg_p)$ freezes both kernels. The same
five-case induction and rational weights lose at most two shape powers
per Lie order. Coupled occupancy is included: keep each complete pair
as a scalar function of all angles. A derivation changes at most one
new pair per product-rule step, even when its occupancy derivative
moves many roots. Thus each descendant changes at most $j$ pairs,
has $N^{O(j)}$ terms, and retains \eqref{iso:quadratic} on the others.

For maximal separation at least $N^{-K_0}$ all selected losses are
polynomial. The corresponding real coarea assertion also holds for
the genuinely coupled contour. Fix the other angles and the mean of
an ordered extreme pair. By permutation symmetry
\(G(\ldots,m+d,\ldots,m-d,\ldots)\) is even in \(d\), and
its second derivative is at most \(-2c_R'\) for a fixed
\(c_R'>0\). Integrating from \(d=0\) gives
\[
 |\partial_pG-\partial_qG|
       \ge c_R'|\theta_p-\theta_q|.
\]
Consequently the phase map \((\theta_p,\theta_q)\mapsto
(\theta_p+\theta_q,G)\) has multiplicity one on each ordered
extreme-index cell and Jacobian at least \(c_R'\) times the maximal
separation. Its coarea cost is polynomial in \(N\); unwrapped
phase lengths \(O(N)\) give a further
\(N^{O(1)}(H+\log N)^2\) factor by the simple-kernel estimates.
Below that scale use $\theta_i=\mu+\rho\omega_i$,
$\sum\omega_i=0$, $\sum\omega_i^2=1$.
Then $\rho\le\sqrt N\,N^{-K_0}$, not merely $N^{-K_0}$.
The phase map $(\mu,\rho)\mapsto(N\mu,G)$ has Jacobian at least
$c_RN\rho/2$ and multiplicity one on each ray. The numerator costs
$C^{N^2}\rho^{N(N-1)}$. After the same conservative losses its
remaining coarea power is $L_N>0$ above. A fixed sufficiently large
$K_0>1$ absorbs $C^{N^2}$ and the polynomial field costs. The phase
rectangle has side lengths $O(N)$ and attenuation floors $cN\eps$,
so its integral costs $CN^2(H+\log N)^2$. Rational-weight and
full-equality fluxes vanish as already checked. This proves
\eqref{iso:middle-bound} in the last support type.

\paragraph{Onsite terms and completion of the tail.}
For $f_N$ repeat the exact weighted deformation with both product poles
deleted. The selected estimate \eqref{iso:selected-bound} remains
valid. For the other pieces use only the one-phase branch field
$V_i=(\partial_\zeta\phi_i)/(\partial_{\theta_i}\phi_i)$ and zero
velocity on compact indices. It freezes each branch phase, has bounded
analytic coefficients, and preserves occupancy on the branch plateaus.
There are no selected poles or punctures. Every fixed derivative changes
at most $j$ pairs; regular branch lengths are uniformly integrable by
\eqref{iso:body}. Thus the remaining pieces are bounded by
$C_j^NN^{C_j}e^{-\kappa N^2}$ for every $N$ and
$0\le\lambda\le1$. Consequently
\begin{equation}
 \sum_N|f_N^{(j)}(\zeta_\eps)|
 \le C_j\exp(C_j\log^2(H+2))=o(\eps^{-\delta})
 \quad\hbox{for every }\delta>0.
 \label{iso:onsite-bound}
\end{equation}

Differentiate the locally normally convergent series first. Only
afterwards split its already differentiated terms and current at
$N=D\sqrt H$ and $\lambda_*$.
Choose a common $\kappa$ in the pair estimates and then
$D$ with $\kappa D^2>k+2$.
The selected sum is $e^{O_j(\log^2(H+2))}$.
The large current is $O_j(e^{\alpha_0(j+2)H})=o(e^{H/2})$.
The entire remaining high-order tail from \eqref{iso:original-bound}
is at most
\[
 \exp\bigl((j+2-\kappa D^2)H+O_j(\sqrt H)\bigr)=o(e^{H/2}).
\]
Summing \eqref{iso:middle-bound} over all intermediate orders costs
$\exp(O_j(\sqrt H\log(H+2)))=o(e^{H/2})$.
This proves the assertion for signed $T_N^\sigma$ at positive $E_*$;
\eqref{iso:sign} and conjugation cover the other odd family.
Equation \eqref{iso:onsite-bound} and $C_N=f_N+2T_N$ prove the
full-site assertion of Theorem~\ref{iso:tail}.

\paragraph{Order of constants and quantifiers.}
The point, $N_0,k$, radial margin and all required derivative orders are
fixed first. Next fix branch and compact margins and the nested slope
separation; then the same-group pair bound, its cross slack, and only
then the endpoint slivers and $\tau$. Choose disk constants after these
data, preserving the branch and named-anchor gaps. Next fix
$\alpha_0,D,M,K_0$ and the algebraic separation constant $A$.
Choose the microcore prefactor and determine its anchor-neighborhood
costs before choosing $B,C_1$ large enough for the inclusions at every
$N\ge N_0+2$ and the negative cubic coefficient. Finally choose one
$\eps_0$ small enough for all the preceding estimates. The elementary
limits $\eps^r e^{CN}N^m\to0$ for $N\le D\sqrt H$ and fixed
$r>0,C,m$ make this single choice uniform in that window.
Finitely many small-order constants may be large but remain independent
of $\eps$. No particle or current window endpoint is differentiated.

\subsection{Dense local obstructions and the exact physical identities}

We now restore the physical normalization. The first singular
coefficient and the complete tail bound combine before any density
argument is used. Exact coefficient identities then relate the
isotropic functions to the equal-coupling restriction of the joint
functions.

The full-site form-factor identification gives, on the respective real
physical intervals,
\begin{align}
 \mathcal X_L(s)&=(1-s^{-4})^{1/4}
                  \sum_{N\ge2,\ N\ {\rm even}}C_N(s),
                  &&s>1,\label{iso:physical-L}\\
 \mathcal X_H(s)&=\frac{(1-s^4)^{1/4}}s
                  \sum_{N\ge1,\ N\ {\rm odd}}C_N(1/s),
                  &&0<s<1.\label{iso:physical-H}
\end{align}
The low-temperature formula has the plus-phase subtraction already
performed: its omitted zero-particle term is $M_+^2$, and its onsite
contribution is $1-M_+^2$. Equation \eqref{iso:sites} therefore also
gives the exact alternative
$\mathcal X_L=1-M_+^2+2M_+^2\sum_{N\ge2,\ N\ {\rm even}}T_N$.
The high-temperature onsite sum is one, and $C_1=1/(E-2)$ checks the
normalization. The branches in \eqref{iso:physical-L}--\eqref{iso:physical-H}
are positive on their stated real intervals.

\begin{proposition}[Dense obstructions for the physical isotropic sums]
\label{iso:local-obstruction}
The full normalized even exterior sum in \eqref{iso:physical-L}
and the full normalized odd interior sum in \eqref{iso:physical-H}
have no local holomorphic extension at a dense set of points of
\(|s|=1\). At their selected points the nonzero leading derivatives
are those of Proposition~\ref{iso:first}, multiplied by the physical
nonvanishing factors. These assertions use the full-tail
bounds of Theorem~\ref{iso:tail}.
\end{proposition}
\begin{proof}
At every selected point, Proposition~\ref{iso:first}, the bounded lower
orders, and Theorem~\ref{iso:tail} give a nonzero leading derivative of
the complete physical sum. All lower derivatives are smaller than that
leading scale, so multiplication by its local analytic nonzero phase
prefactor preserves it by the product rule. A holomorphic extension
through such a point would bound that fixed finite derivative on the
short radial segment, a contradiction. The selected points are dense on
the required arcs. For the even sum, the contour transformation
$(s,x,y)\mapsto(-s,-x,-y)$ gives evenness, and conjugation completes
all four quadrants. For the odd sum both signs were treated explicitly;
conjugation completes a dense subset of the circle. The inversion
\(s=1/\zeta\) is locally biholomorphic at each unit-circle point,
so the odd conclusion concerns the physical interior germ. This proves
the stated local obstructions. Their passage to all boundary points and
their role in the exact slice and isotropic corollaries are established
in Section~\ref{main:closure}.
\end{proof}

The high-temperature value at $s=0$ is the removable value one.
This can also be checked
without the joint representation: uniformly in the real angular
variables, $z=s/2+O(s^2)$, each one-body factor is $O(s)$ and each
squared pair is $O(s^2)$, giving $C_N(1/s)=O(s^{N^2})$ with a
bound $C^{N^2}|s|^{N^2}/N!$ for sufficiently small $|s|$.
The $N=1$ coefficient is $s/(1-s)^2$, and every higher odd order
vanishes after division by $s$. The exterior normalized
low-temperature germ is single valued; the energy-field susceptibility
is interpreted on the physical continuation of
$\beta_{\rm phys}=\operatorname{arsinh}(s)/(2J)$.
At every selected nonaxis point this factor is analytic and nonzero,
so the same dense obstructions hold for $\chi=\beta_{\rm phys}\mathcal X$.
No global single-valued inverse temperature on the exterior is asserted.

Finally put $s=\sinh z$ on the equal-coupling restriction.
The central cell
\[
 |\Im z|<\pi/2,\qquad
 |\Re z|<\operatorname{arsinh}(\cos\Im z)
\]
maps biholomorphically onto $|s|<1$. The right cell
\[
 \Re z>\operatorname{arsinh}(|\cos\Im z|)
\]
is the universal cover of $|s|>1$ under the same map, with deck
translations $z\mapsto z+2\pi\ii$.
One direct verification writes $w=e^z$ and solves $w-w^{-1}=2s$:
for $|s|>1$ exactly one root has $|w|>1$, since a unit root would give
$s\in\ii[-1,1]$. Exponentiation supplies the stated covering.
At the selected nonramified boundary points the local isotropic
obstructions pull back directly. These identities fix the entire
even/odd sums and the physical factor $\beta=z/(2J)$ used in the
proof of the slice theorem and its isotropic descendants below.

\section{Proofs of the joint, slice and isotropic theorems}
\label{main:closure}

All local estimates are now available. We first prove the joint
boundary theorem, then prove nonextension on each fixed-ratio line
using its own transverse geometry, and finally pass through the
isotropic covering map. This order keeps restriction and descent
logically separate.

\subsection{The complete joint obstruction}

\begin{proof}[Proof of Theorem~\ref{main:joint}]
Theorem~\ref{geo:exact-sigma} determines the spectral set at every
coupling, including collinear and zero-real-part degeneracies.
Theorem~\ref{geo:components} identifies the physical components
and their contractibility and \(p\)-bounds. Proposition~\ref{geo:pseudoconvex}
gives an explicit plurisubharmonic exhaustion, independent of the
susceptibility. Proposition~\ref{ana:normal}, with the physical
interface and its coefficientwise angular identification, proves
joint holomorphy and normal convergence of the actual physical sums.

Take a regular face point whose free coefficient is
\(b=\cos(2\pi a/\ell)\), with \(\ell\) an odd prime and
\(0<a<\ell/2\). By Proposition~\ref{geo:local-coordinates} the
variables \((C,B)\), after exchanging transfer axes if needed,
are independent holomorphic local coordinates, \(\Im B\ne0\),
and the physical domain is the full complementary side of that edge.
At such a point \(S,T\ne0\), and the edge identity
\eqref{geo:p-edge} implies \(p^2\ne1\). Thus the locally chosen
physical prefactor is analytic and nonzero.

At a positive endpoint, restrict to any holomorphic transversal with
nonramified edge coordinate \(u=C-B-b\). Its actual anisotropy
is a holomorphic function \(B(u)\). Conjugating the data when
needed gives \(\Im B(0)>0\), and the physical side contains
\(u=\ii\eps\). Proposition~\ref{ana:positive} proves a nonzero
singular derivative of the entire odd or even sum, including the
complete differentiated higher tail and physical prefactor. The
same proposition handles a negative endpoint for the even sum by
the exact angular shift and, if necessary, conjugation.

For a negative endpoint of the odd sum use
\[
             t=\frac{C+B-b}{B}.
\]
This is also nonramified on the stipulated transversal, because
\(B(0)\ne0\). In its lower-half-root orientation the actual
\(B(t)\) has \(\Im B(0)<0\), and the local physical side is
\(\Im(B(t)t)>0\). Theorem~\ref{neg:theorem} gives all finite
boundary jets of the complete odd sum and a selected sequence
whose coefficient roots diverge. Its exact full-sum asymptotic
\eqref{neg:full-asymptotic} excludes cancellation by every other
particle index and winding. Division by the nonvanishing physical
prefactor proves the same nonextension for the physical response.
The opposite orientation is its complex conjugate. The local uniform
complex-parameter hypotheses and constants are stated and proved in
Sections~\ref{ana:edges} and \ref{neg:section}; no estimate for
real fixed ratios is promoted to a joint estimate without proof.

In particular choose a constant-\(B\) transversal. If the joint
physical function extended to a neighborhood of the selected point,
its restriction to that disk would give the forbidden one-variable
extension. Thus the transversal obstruction gives ambient nonextension at
each selected point. The selected points are dense
on regular faces by Proposition~\ref{geo:local-coordinates}, and
regular faces are dense in the entire physical boundaries by
Proposition~\ref{geo:joint-density}. Corollary~\ref{geo:propagate-obstruction}
then excludes componentwise extension at every boundary point,
including corners, collinear points and vanishing transfer-coordinate
charts. This proves the joint assertion and its transverse-curve
strengthening. Proposition~\ref{iso:local-obstruction}, proved
from the full sums in Section~\ref{iso:section}, supplies the
additional local assertion on the collinear equal-coupling restriction.

Finally consider analytic continuation in an unramified Riemann domain
over \(\C^2\) from one of these physical germs. As long as the
projected continuation path remains in its domain, uniqueness of
analytic continuation makes every continued germ the germ of the
single-valued function already constructed there. If a path first
leaves that domain, its first exit is a boundary point. A local
biholomorphic chart of the Riemann domain at the continued point
would yield a holomorphic function on a base neighborhood agreeing
with the physical function on an accumulating domain-side component.
The proved componentwise obstruction forbids this. Compactness of
the finite path interval guarantees a finite first exit whenever the
path reaches outside. Thus neither a monodromy construction nor a
different unramified physical sheet bypasses the boundary. This
uses the standard notion of holomorphic continuation, not a
definition tied to the spectral set.
\end{proof}

\subsection{All physical ratios, including the collinear one}

\begin{lemma}[Actual complex-ratio neighborhoods]\label{main:complexratio}
At any regular selected point of a positive-ratio slice where its edge
coordinate has nonzero derivative, the transverse representation and
the boundary estimates hold in a neighborhood with \(\lambda\)
complex. The uniformity is on common compact parameter polydisks,
not merely pointwise on the real-ratio family.
\end{lemma}
\begin{proof}
The boundary temperature coordinate \(z_0\) is nonzero. Thus
\((z,\lambda)\mapsto(z,w=\lambda z)\) has determinant \(z_0\)
and is a local biholomorphism. Its composition with the regular
\((C,B)\) coordinate map is holomorphic. If the positive-edge
parameter is \(u=C-B-b\), its derivative in \(z\) is nonzero
by hypothesis; the holomorphic implicit-function theorem gives
\(z=z(u,\lambda)\) and hence \(B=B(u,\lambda)\) on a product
polydisk. At a negative edge use \(t=(C+B-b)/B\) in the same way.
Shrinking the polydisk preserves the strict sign of \(\Im B\),
the nonvanishing physical prefactor and the other strict face
conditions. Cauchy estimates on a slightly larger polydisk give
the common holomorphic bounds required in
Proposition~\ref{ana:positive} and Theorem~\ref{neg:theorem}.
Their proofs therefore apply to the actual complex-ratio family.
In the negative case the selected integers may depend on the
family parameter, as explicitly stated there. No boundary estimate
is claimed uniform through the collinear degeneration at
\(\lambda=1\); that restriction is treated by its own proved
full-sum estimates. Joint holomorphy on interior compact sets already
includes complex-ratio neighborhoods at \(\lambda=1\).
\end{proof}

\begin{proof}[Proof of Theorem~\ref{main:slices}]
The exact open-set equalities are Theorem~\ref{geo:exact-pullbacks}.
The restrictions are jointly constructed physical functions by
Proposition~\ref{ana:normal}, and are holomorphic on precisely the
stated pullbacks. A boundary point of either pullback maps to the
corresponding joint boundary: it is a limit of domain points and
cannot map to the joint interior by the same open-set equality.

First fix \(\lambda\ne1\). Proposition~\ref{geo:slice-density}
proves, separately on this actual complex line, density of regular
boundary points at which the local edge map has nonzero derivative
and its free phase is one of the chosen prime-grid values. This
proof includes every finite positive ratio, whether rational or
irrational. In particular it does not rely on intersecting an
arbitrary ambient dense set with the line. At each such point the
temperature line is a holomorphic curve \(B=B(u)\) or \(B=B(t)\)
of exactly the kind treated in the proof of Theorem~\ref{main:joint}.
Apply Proposition~\ref{ana:positive} or Theorem~\ref{neg:theorem}
to this actual curve, retaining its varying anisotropy. These local estimates prove
nonextension of the restricted complete physical sum. The
one-variable collar in Lemma~\ref{geo:collar} and
Corollary~\ref{geo:propagate-obstruction} give every boundary point.

Now let \(\lambda=1\). The exact coefficient and normalization
identities in \eqref{iso:physical-L}--\eqref{iso:physical-H}
identify the restriction with the full even and odd integrals of
Section~\ref{iso:section}; its local obstructions were proved in
Proposition~\ref{iso:local-obstruction}. Lemma~\ref{geo:isotropic-cover} maps the two
slice domains to the exterior and disk, respectively. At every
boundary point with \(\cosh\zeta\ne0\) this is a local
biholomorphism. The selected obstructing points pull back densely
on the nonramified boundary arcs. The remaining points with
\(\cosh\zeta=0\) are limits of those arcs. The same local collar
therefore proves nonextension everywhere on both slice boundaries.
This argument uses no limit as \(\lambda\to1\).

All nonzero even orders were retained after subtraction of the
zero-particle \(M^2\), and the exact full-site multiplier retains
the origin and both axes. Section~\ref{ana:physical} proves the
physical energy-field equality \(\chi=\beta\mathcal X\) with
the plus phase fixed in LT. Here \(\beta=\zeta/(2J_x)\) is an
entire function of the slice coordinate and is nonzero on both
boundaries: zero lies in the HT interior and is strictly outside
the LT closure. Multiplication or division by it cannot change
local boundary extendability. The result consequently establishes the claims in both response conventions.
\end{proof}

\subsection{Derivation of the complete isotropic theorems from the slice theorem}

\begin{proof}[Proof of Corollary~\ref{main:isotropic}]
Apply the now proved Theorem~\ref{main:slices} with \(\lambda=1\).
The identities \(S=T=s\), \(p=s^2\), \(B=1\), and
\[
 \mathcal D(\zeta_1,\eta_1;z,z)
 =s\left[s+s^{-1}
 -\tfrac12(\zeta_1+\zeta_1^{-1}+\eta_1+\eta_1^{-1})\right]
\]
identify the parent coefficients term by term with the full isotropic
coefficients. On the real LT interval \(\sqrt p=s\), so each
factor \(\sqrt p/\mathcal D\) is exactly its isotropic counterpart.
The even prefactor becomes \((1-s^{-4})^{1/4}\). On the real HT
interval the dispersion is \(s+s^{-1}-\cos q\) and the odd
prefactor is \((1-s^4)^{1/4}/s\), giving the exterior auxiliary
variable \(1/s\). Normal convergence and the identity theorem
extend these identities to their physical complex domains.

For HT, the biholomorphism in Lemma~\ref{geo:isotropic-cover}
transports the restricted holomorphic function to the disk. The
removable value at zero is one, independently checked at the end
of Section~\ref{iso:section}. For LT, the normalized termwise
identities show invariance under every deck translation
\(z\mapsto z+2\pi\ii k\); the full normalized function therefore
descends from \(\Omega_{L,1}\) to a single-valued holomorphic
function on \(|s|>1\). It is the positive real LT germ and has
the stated connected onsite contribution.

Suppose one of these descended functions extended through a point
\(s_0\) of the unit circle. Choose a boundary preimage \(z_0\)
on its physical slice. Composing the proposed extension with the
entire map \(\sinh\) would extend the slice function near
\(z_0\), in agreement on the physical side. This contradicts
Theorem~\ref{main:slices}. This composition argument also works
at the ramified boundary values \(s_0=\pm\ii\); no local inverse
at those values is required. Hence the entire unit circle is a
natural boundary for both normalized germs.

Finally \(\beta=z/(2J)\) on the prescribed physical temperature
continuation yields the energy-field claims. In the disk it is
the branch \(\operatorname{arsinh}(s)/(2J)\). In the exterior it
is interpreted on the proved covering. At all selected nonaxis
points it is analytic and nonzero, so the local singular
derivative remains nonzero after the product rule; density gives
all remaining boundary points. This proves the HT and connected LT descendants with their
distinct germs and complete odd/even sums. Their natural-boundary
conclusions follow from the estimates proved in this paper.
\end{proof}

\section{Two examples on the scope of continuation}
\label{main:verification}

Two exact examples clarify the scope of these results. First the
decoupled restriction \(K_y=0\) is entire, as proved in
Remark~\ref{ana:axis}; it refutes an unrestricted principle that
ambient nonextension passes to every complex line. The positive-ratio
slice theorem instead uses its explicit transverse-curve argument and
separate collinear proof. Second, \(\Sigma\) cannot be replaced by
\(\{|p|=1\}\): at
\[
 z=\ii\pi/2,\qquad w=\operatorname{arsinh}(1/2)
\]
one has \(A=0,S=\ii,T=1/2\), so \(a=b=0\) represents a
spectral point, while \(|p|=1/2\). In fact \(F_1=-1/4,F_2=1\),
so this point lies in the interior of \(\Sigma\). Thus the modulus level set is strictly smaller than the spectral
set; the example makes no assertion of a physical full-sum
singularity at this interior spectral point.



\begin{thebibliography}{99}
\bibitem{BL} A. I. Bugrij and O. Lisovyy,
\emph{Correlation function of the two-dimensional Ising model on a
finite lattice. II}, Theoretical and Mathematical Physics \textbf{140}
(2004), 987--1000. \href{https://arxiv.org/abs/0708.3643v1}{arXiv:0708.3643v1}.
\bibitem{IL} N. Iorgov and O. Lisovyy,
\emph{Ising correlations and elliptic determinants},
Journal of Statistical Physics \textbf{143} (2011), 33--59.
\href{https://doi.org/10.1007/s10955-011-0154-6}{doi:10.1007/s10955-011-0154-6}.
\href{https://arxiv.org/abs/1012.2856v2}{arXiv:1012.2856v2}.
\bibitem{ONGP} W. P. Orrick, B. G. Nickel, A. J. Guttmann and J. H. H. Perk,
\emph{The susceptibility of the square lattice Ising model: New developments},
Journal of Statistical Physics \textbf{102} (2001), 795--841.
\href{https://doi.org/10.1023/A:1004850919647}{doi:10.1023/A:1004850919647}.
\href{https://arxiv.org/abs/cond-mat/0103074v1}{arXiv:cond-mat/0103074v1}.
\bibitem{BHM} S. Boukraa, S. Hassani and J.-M. Maillard,
\emph{Holonomic functions of several complex variables and singularities
of anisotropic Ising n-fold integrals}, Journal of Physics A:
Mathematical and Theoretical \textbf{45} (2012), 494010.
\href{https://doi.org/10.1088/1751-8113/45/49/494010}{doi:10.1088/1751-8113/45/49/494010}.
\href{https://arxiv.org/abs/1207.1784v2}{arXiv:1207.1784v2}.
\bibitem{Raoufi} A. Raoufi,
\emph{Translation-invariant Gibbs states of Ising model: General setting},
preprint (2017).
\href{https://arxiv.org/abs/1710.07608v1}{arXiv:1710.07608v1}.
\bibitem{CGN} F. Camia, C. Garban and C. M. Newman,
\emph{The Ising magnetization exponent on \(\mathbb Z^2\) is \(1/15\)},
preprint (2013).
\href{https://arxiv.org/abs/1205.6612v4}{arXiv:1205.6612v4}.

\end{thebibliography}
\end{document}

